\documentclass[12pt]{TA-ITS}
\usepackage{booktabs}
\usepackage{longtable}
\usepackage{array}
\usepackage{ragged2e}
\usepackage[T1]{fontenc}
\usepackage{adjustbox}
\usepackage{array}
\usepackage{xcolor}
\usepackage{colortbl}
\usepackage{pdflscape}

\definecolor{headerblue}{RGB}{31,73,125}
\definecolor{rowgray}{RGB}{242,242,242}

% -- Definisi warna --
\definecolor{codebg}{RGB}{248, 249, 250}
\definecolor{codeframe}{RGB}{210, 215, 220}
\definecolor{kw}{RGB}{0, 112, 193}
\definecolor{str}{RGB}{163, 21, 21}
\definecolor{cmt}{RGB}{106, 153, 85}
\definecolor{num}{RGB}{181, 101, 29}
\definecolor{builtin}{RGB}{64, 120, 192}
\definecolor{titlecolor}{RGB}{30, 80, 150}
\definecolor{sectioncolor}{RGB}{40, 95, 165}
\definecolor{headerbg}{RGB}{230, 238, 250}
% ============================================================
% lampiran_setup.tex
% File ini SUDAH INCLUDE semua package yang dibutuhkan.
% Cukup letakkan di preamble dokumen utama:
%
%     \input{lampiran_setup}
%
% sebelum \begin{document} -- tidak perlu tambah package lain.
% ============================================================

% ── Package (hanya load jika belum ada) ────────────────────
\makeatletter
\@ifpackageloaded{xcolor}{}{\usepackage{xcolor}}
\@ifpackageloaded{listings}{}{\usepackage{listings}}
\@ifpackageloaded{caption}{}{\usepackage{caption}}
\@ifpackageloaded{mdframed}{}{\usepackage{mdframed}}
\@ifpackageloaded{amsmath}{}{\usepackage{amsmath}}
\makeatother

% ── Warna ───────────────────────────────────────────────────
\definecolor{codebg}{RGB}{248, 249, 250}
\definecolor{codeframe}{RGB}{210, 215, 220}
\definecolor{kw}{RGB}{0, 112, 193}
\definecolor{str}{RGB}{163, 21, 21}
\definecolor{cmt}{RGB}{106, 153, 85}
\definecolor{builtin}{RGB}{64, 120, 192}
\definecolor{headerbg}{RGB}{230, 238, 250}
\definecolor{breakcol}{RGB}{160, 160, 160}

% ── Style listing Python ────────────────────────────────────
\lstdefinestyle{pycode}{
    language=Python,
    backgroundcolor=\color{codebg},
    basicstyle=\ttfamily\footnotesize,
    keywordstyle=\color{kw}\bfseries,
    stringstyle=\color{str},
    commentstyle=\color{cmt}\itshape,
    numberstyle=\tiny\color{gray},
    emph={True,False,None,self},
    emphstyle={\color{builtin}\bfseries},
    numbers=left,
    numbersep=8pt,
    stepnumber=1,
    tabsize=4,
    showstringspaces=false,
    breaklines=true,
    breakatwhitespace=false,
    breakindent=16pt,
    postbreak=\mbox{\textcolor{breakcol}{$\hookrightarrow$}\space},
    columns=fullflexible,
    keepspaces=true,
    frame=single,
    framerule=0.5pt,
    rulecolor=\color{codeframe},
    xleftmargin=22pt,
    xrightmargin=0pt,
    framexleftmargin=22pt,
    framexrightmargin=0pt,
    captionpos=b,
    aboveskip=8pt,
    belowskip=4pt,
}
\makeatother
% SETELAH ACC
\pagestyle{TA}

% Singkatan Tugas Akhir
\newcommand{\ta}{penelitian Tugas Akhir}
\newcommand{\TA}{Penelitian Tugas Akhir}

% Metode & model
\newcommand{\method}{}
\newcommand{\smethod}{}

\begin{document}
    %%%%%%%%%%%%%%%%%%%%%%%%%%%%%%%%%%%%%%%%%%%%%%%%%%%%%%%%%%%%%%%%%%%%%%%%%%%%%%%%
% DATA-DATA UNTUK PERINTAH KHUSUS
% Masukkan data-data sesuai panduan berikut untuk otomatisasi format Tugas Akhir
%%%%%%%%%%%%%%%%%%%%%%%%%%%%%%%%%%%%%%%%%%%%%%%%%%%%%%%%%%%%%%%%%%%%%%%%%%%%%%%%

% Judul Tugas Akhir dalam Bahasa Indonesia & Inggris

\Judul{Optimasi \textit{Hyperparameter} pada \textit{Tree-based Models} Menggunakan Algoritma Genetika untuk Prediksi Pendapatan Layanan SMS Internasional}
\Title{Hyperparameter Optimization on Tree-based Models Using Genetic Algorithm for Predicting Revenue of SMS International Service}

% Nama mahasiswa yang mengerjakan Tugas Akhir
\Nama{Jiryan Farokhi}

% Nama alias atau nama kecil; jika ada
\NamaKecil{} % Bisa dikosongin, gapapa

% NRP mahasiswa yang mengerjakan Tugas Akhir
\NRP{5002221102}

% Kode mata kuliah Tugas Akhir
\KodeTA{SM234801}

% Asal departemen dalam Bahasa Indonesia & Inggris
\Departemen{Matematika} 
\Department{Mathematics}

% Asal laboratorium, bidang studi atau peminatan; tidak perlu (tahun 2024)
% \Laboratorium{Pembelajaran Mesin dan Big Data}

% Asal fakultas dalam Bahasa Indonesia & Inggris
\Fakultas{Sains dan Analitika Data} % Tanpa kata `fakultas`-nya
\Faculty{Science and Data Analytics} % Tanpa kata `faculty of`-nya`

% Singkatan fakultas dalam Bahasa Indonesia & Inggris
\SingkatanFakultas{FSAD}
\SingkatanFakultasEng{SCIENTICS}

% Tanggal disetujuinya Tugas Akhir
\TanggalDisetujui{6}{ Juni }{2026} % {<tanggal>}{<bulan>}{<tahun>}

% Profil kepala departemen
\Kadep{Dr. Didik Khusnul Arif, S.Si., M.Si.} % Nama kepala departemen
\NIPKadep{19730930 199702 1 001} % NIP kepala departemen

% Dosen Pembimbing
\DosenPembimbing{Dr. Drs. Bandung Arry S., MI.Komp} % Pembimbing pertama
                {} % Pembimbing kedua; abaikan, jika tidak ada
                {} % Pembimbing ketiga; abaikan, jika tidak ada
\NIPPembimbing{19630605 198903 1 003} % NIP Pembimbing pertama
              {} % NIP Pembimbing kedua; abaikan, jika tidak ada
              {} % NIP Pembimbing ketiga; abaikan, jika tidak ada

% Dosen Penguji
\DosenPenguji{Dr. Budi Setiyono, S.Si, MT} % Penguji pertama
             {Muhammad Yusuf Hakim Widianto, S.Si., M. Si., M.Sc., Ph.D} % Penguji kedua
             {Prof. Subchan, S.Si, M.Sc, Ph.D} % Penguji ketiga
\NIPPenguji{19720207 199702 1 001} % NIP Penguji pertama
           {199106202024061001} % NIP Penguji kedua
           {19710513 199702 1 001} % NIP Penguji ketiga
    \def\topseparator{-.2\paperheight}
\def\titleimagepadding{0.75in}
\def\companyimagepadding{-\titleimagepadding/1.75}
\def\textpadding{0.6in}
\def\ITSsize{1.5in}
\def\companysize{0.4in}

\pagenumbering{gobble}
% Kita gunakan tikzpicture agar lebih mudah mengaturnya (sekaligus bisa menerapkan hitungan matematis)
\makeatletter % Anggap `@` sebagai suatu simbol dalam perintah
\thispagestyle{empty}
{
\begin{tikzpicture}[remember picture, overlay]

    %%%%%%%%%%%%%%%
    % Daerah Konten
    %%%%%%%%%%%%%%%
    \fill[itsblue] (current page.south west) rectangle (current page.north east);
    \node[anchor=north west, text width=\paperwidth-1.8in, color=white, font=, align=justify] at ([xshift = \titleimagepadding/0.7, yshift=\topseparator-\textpadding] current page.north west) 
        {
            
            \setstretch{1.5} % Atur spasi antar baris
            \noindent \hspace{-2.5mm} % Dipastikan tidak ada indentasi; edit spasi
            \textsf{\fontsize{\coverfontsize}{\fontbaselineskip}\selectfont{\textbf{\MakeUppercase{Tugas Akhir --- \kode@TA}}}}

            \vspace{0.05\paperheight}
            \textsf{\fontsize{\titlefontsize}{\fontbaselineskip}\selectfont{\textbf{\MakeUppercase{\@judul}}}}

            \vspace{0.075\paperheight}
            \textsf{\fontsize{\coverfontsize}{\fontbaselineskip}\selectfont{\textbf{\MakeUppercase{\@nama}}}}

            \textsf{\fontsize{\coverfontsize}{\fontbaselineskip}\selectfont{NRP. \@NRP}}

            \vspace{0.07\paperheight}
            \textsf{\fontsize{\coverfontsize}{\fontbaselineskip}\selectfont{Dosen Pembimbing}}
            
            \ifthenelse{\equal{\pbb@tiga}{}\OR\equal{\pbb@tiga}{ }} % Cek apakah ada pembimbing ketiga atau tidak
                {% Kalau tidak ada, cek ada pembimbing kedua atau tidak
                    \ifthenelse{\equal{\pbb@dua}{}\OR\equal{\pbb@dua}{ }}
                        {% Kalau tidak ada juga, cukup tulis pembimbing pertama
                            \textsf{\fontsize{\coverfontsize}{\fontbaselineskip}\selectfont{\textbf{\pbb@satu}\\ NIP. \NIP@pbb@satu}}\\
                        }
                        {% Kalau ada, tulis kedua pembimbing saja
                            \textsf{\fontsize{\coverfontsize}{\fontbaselineskip}\selectfont{\textbf{\pbb@satu}\\ NIP. \NIP@pbb@satu}}
                            
                            \textsf{\fontsize{\coverfontsize}{\fontbaselineskip}\selectfont{\textbf{\pbb@dua}\\ NIP. \NIP@pbb@dua}}
                        }
                }
                {% Kalau ada, tulis ketiga pembimbingnya
                    \textsf{\fontsize{\coverfontsize}{\fontbaselineskip}\selectfont{\textbf{\pbb@satu}\\ NIP. \NIP@pbb@satu}}
                    
                    \textsf{\fontsize{\coverfontsize}{\fontbaselineskip}\selectfont{\textbf{\pbb@dua}\\ NIP. \NIP@pbb@dua}}
                    
                    \textsf{\fontsize{\coverfontsize}{\fontbaselineskip}\selectfont{\textbf{\pbb@tiga}\\ NIP. \NIP@pbb@tiga}}\\
                }

            \vspace{0.05\paperheight}
            \textsf{\fontsize{\regularfontsize}{\fontbaselineskip}\selectfont{\textbf{Program Sarjana}}}
            
            \textsf{\fontsize{\regularfontsize}{\fontbaselineskip}\selectfont{Departemen~\@departemen}}
            
            \textsf{\fontsize{\regularfontsize}{\fontbaselineskip}\selectfont{Fakultas~\@fakultas}}
            
            \textsf{\fontsize{\regularfontsize}{\fontbaselineskip}\selectfont{Institut Teknologi Sepuluh Nopember}}
            
            \textsf{\fontsize{\regularfontsize}{\fontbaselineskip}\selectfont{Surabaya}}
            
            \textsf{\fontsize{\regularfontsize}{\fontbaselineskip}\selectfont{2026}}

            % Kita sisakan spasi kosong, karena alasan yang jelas untuk spasi di cover
        };
    
    %%%%%%%%%%%%%%%
    % Daerah Gambar
    %%%%%%%%%%%%%%%
    \fill[white] (current page.north west) rectangle ++(\paperwidth, \topseparator);

    % Logo ITS
    \node[anchor=north west] at ([xshift=\titleimagepadding, yshift=\topseparator/2+\ITSsize/2] current page.north west) {
\includegraphics[height = \ITSsize]{tugas-akhir/assets/images/its blue.png}};

    % Logo Perusahaan
    % \node[anchor=north east] at ([xshift=\companyimagepadding, yshift=\topseparator/2+\companysize/2] current page.north east) {\includegraphics[height = \companysize]{LuDoYa-Internship/assets/images/report-need-stuffs/Logo SKPB-biru.png}};

\end{tikzpicture}}
\newpage
\makeatother % Kembalikan `@` sebagai simbol khusus
\makeatletter
\thispagestyle{empty}

{
\begin{tikzpicture}[remember picture, overlay]

    %%%%%%%%%%%%%%%
    % Daerah Konten
    %%%%%%%%%%%%%%%
    \fill[itsblue] ([yshift=\topseparator] current page.north west) rectangle ++(\paperwidth, -1cm);
    \node[anchor=north west, text width=\paperwidth-2.3in, color=black, font=, align=justify] at ([xshift = \titleimagepadding/0.7, yshift=\topseparator-\textpadding/0.7] current page.north west) {%
        \setstretch{1.5}\noindent
        {\fontsize{\coverfontsize}{\fontbaselineskip}\selectfont{\textbf{\MakeUppercase{Tugas Akhir --- \kode@TA}}}}

        \vspace{0.05\paperheight}
        {\fontsize{\titlefontsize}{\fontbaselineskip}\selectfont{\textbf{\MakeUppercase{\@judul}}}}

        \vspace{0.065\paperheight}
        {\fontsize{\coverfontsize}{\fontbaselineskip}\selectfont{\textbf{\MakeUppercase{\@nama}}}}
        
        {\fontsize{\coverfontsize}{\fontbaselineskip}\selectfont{NRP. \@NRP}}

        \vspace{0.07\paperheight}

        {\fontsize{\coverfontsize}{\fontbaselineskip}\selectfont{Dosen Pembimbing}}
        
        \ifthenelse{\equal{\pbb@tiga}{}\OR\equal{\pbb@tiga}{ }}
            {
                \ifthenelse{\equal{\pbb@dua}{}\OR\equal{\pbb@dua}{ }}
                    {
                        {\fontsize{\coverfontsize}{\fontbaselineskip}\selectfont{\textbf{\pbb@satu}\\ NIP. \NIP@pbb@satu}}
                        
                    }
                    {
                        {\fontsize{\coverfontsize}{\fontbaselineskip}\selectfont{\textbf{\pbb@satu}\\ NIP. \NIP@pbb@satu}}
                        
                        {\fontsize{\coverfontsize}{\fontbaselineskip}\selectfont{\textbf{\pbb@dua}\\ NIP. \NIP@pbb@dua}}
                    }
            }
            {
                {\fontsize{\coverfontsize}{\fontbaselineskip}\selectfont{\textbf{\pbb@satu}\\ NIP. \NIP@pbb@satu}}
                
                {\fontsize{\coverfontsize}{\fontbaselineskip}\selectfont{\textbf{\pbb@dua}\\ NIP. \NIP@pbb@dua}}
                
                {\fontsize{\coverfontsize}{\fontbaselineskip}\selectfont{\textbf{\pbb@tiga}\\ NIP. \NIP@pbb@tiga}}
                
            }
        
        \vspace{0.05\paperheight}

        {\fontsize{\regularfontsize}{\fontbaselineskip}\selectfont{\textbf{Program Sarjana}}}
                
        {\fontsize{\regularfontsize}{\fontbaselineskip}\selectfont{Departemen~\@departemen}}
        
        {\fontsize{\regularfontsize}{\fontbaselineskip}\selectfont{Fakultas~\@fakultas}}
        
        {\fontsize{\regularfontsize}{\fontbaselineskip}\selectfont{Institut Teknologi Sepuluh Nopember}}
        
        {\fontsize{\regularfontsize}{\fontbaselineskip}\selectfont{Surabaya}}

        \vspace{0.01\paperheight}
        {\fontsize{\regularfontsize}{\fontbaselineskip}\selectfont{2026}}
    };
    
    %%%%%%%%%%%%%%%
    % Daerah Gambar
    %%%%%%%%%%%%%%%
    \fill[white] (current page.north west) rectangle ++(\paperwidth, \topseparator);

    % Logo ITS
    \node[anchor=north west] at ([xshift=\titleimagepadding, yshift=\topseparator/2+\ITSsize/2] current page.north west) {
\includegraphics[height = \ITSsize]{tugas-akhir/assets/images/its.png}};
\end{tikzpicture}}
\newpage
\makeatother
\makeatletter
\thispagestyle{empty}

{
\begin{tikzpicture}[remember picture, overlay]

    %%%%%%%%%%%%%%%
    % Daerah Konten
    %%%%%%%%%%%%%%%
    \fill[itsblue] ([yshift=\topseparator] current page.north west) rectangle ++(\paperwidth, -1cm);
    \node[anchor=north west, text width=\paperwidth-2.3in, color=black, font=, align=justify] at ([xshift = \titleimagepadding/0.7, yshift=\topseparator-\textpadding/0.7] current page.north west) {%
        \setstretch{1.5}\noindent
        % {\fontsize{\coverfontsize}{\fontbaselineskip}\selectfont{\textbf{\MakeUppercase{Final Year Project --- \kode@TA}}}}
        {\large{\textbf{\MakeUppercase{Final Year Project --- \kode@TA}}}}

        \vspace{0.05\paperheight}
        % {\fontsize{\titlefontsize}{\fontbaselineskip}\selectfont{\textbf{\MakeUppercase{\@title@}}}}
        {\Large{\textbf{\MakeUppercase{\@title@}}}}

        \vspace{0.065\paperheight}
        % {\fontsize{\coverfontsize}{\fontbaselineskip}\selectfont{\textbf{\MakeUppercase{\@nama}}}}
        {\large{\textbf{\MakeUppercase{\@nama}}}}
        
        % {\fontsize{\coverfontsize}{\fontbaselineskip}\selectfont{Student ID \@NRP}}
        {\large{Student ID. \@NRP}}

        \vspace{0.07\paperheight}
        {\fontsize{\coverfontsize}{\fontbaselineskip}\selectfont{Advisor}}
        
        \ifthenelse{\equal{\pbb@tiga}{}\OR\equal{\pbb@tiga}{ }}
            {
                \ifthenelse{\equal{\pbb@dua}{}\OR\equal{\pbb@dua}{ }}
                    {
                        {\fontsize{\coverfontsize}{\fontbaselineskip}\selectfont{\textbf{\pbb@satu}\\ EIN. \NIP@pbb@satu}}
                        
                    }
                    {
                        {\fontsize{\coverfontsize}{\fontbaselineskip}\selectfont{\textbf{\pbb@satu}\\ EIN. \NIP@pbb@satu}}
                        
                        {\fontsize{\coverfontsize}{\fontbaselineskip}\selectfont{\textbf{\pbb@dua}\\ EIN. \NIP@pbb@dua}}
                    }
            }
            {
                {\fontsize{\coverfontsize}{\fontbaselineskip}\selectfont{\textbf{\pbb@satu}\\ EIN. \NIP@pbb@satu}}
                
                {\fontsize{\coverfontsize}{\fontbaselineskip}\selectfont{\textbf{\pbb@dua}\\ EIN. \NIP@pbb@dua}}
                
                {\fontsize{\coverfontsize}{\fontbaselineskip}\selectfont{\textbf{\pbb@tiga}\\ EIN. \NIP@pbb@tiga}}
                
            }
        
        \vspace{0.05\paperheight}

        {\fontsize{\regularfontsize}{\fontbaselineskip}\selectfont{\textbf{Undergraduate Program}}}
                
        {\fontsize{\regularfontsize}{\fontbaselineskip}\selectfont{\@department~Department}}
        
        {\fontsize{\regularfontsize}{\fontbaselineskip}\selectfont{Faculty of \@faculty}}
        
        {\fontsize{\regularfontsize}{\fontbaselineskip}\selectfont{Institut Teknologi Sepuluh Nopember}}
        
        {\fontsize{\regularfontsize}{\fontbaselineskip}\selectfont{Surabaya}}

        \vspace{0.01\paperheight}
        {\fontsize{\regularfontsize}{\fontbaselineskip}\selectfont{2026}}
    };
    
    %%%%%%%%%%%%%%%
    % Daerah Gambar
    %%%%%%%%%%%%%%%
    \fill[white] (current page.north west) rectangle ++(\paperwidth, \topseparator);

    % Logo ITS
    \node[anchor=north west] at ([xshift=\titleimagepadding, yshift=\topseparator/2+\ITSsize/2] current page.north west) {
\includegraphics[height = \ITSsize]{tugas-akhir/assets/images/its.png}};
\end{tikzpicture}}
\newpage
\makeatother

\pagenumbering{roman}
\setcounter{page}{1}
\addcontentsline{toc}{chapter}{Lembar Pengesahan}
\makeatletter
\begin{center}
    \tikz[remember picture,overlay] 
    \node[opacity=1,inner sep=0pt] at (current page.center)
        {
\includegraphics[width=\paperwidth,height=\paperheight]{tugas-akhir/assets/images/WM Lembar Pengesahan.jpg}};

    {\noindent\centering\textbf{\large\MakeUppercase{Lembar Pengesahan}}\par}

    \vspace{1cm}
    \textbf{\MakeUppercase{\@judul}}

    \vspace{0.75cm}
    \textbf{\MakeUppercase{Tugas Akhir}}

    \vspace{0.25cm}
    Diajukan untuk memenuhi salah satu syarat

    memperoleh gelar Sarjana Matematika pada

    Program Studi S-1 \@departemen

    Departemen~\@departemen

    Fakultas~\@fakultas

    Institut Teknologi Sepuluh Nopember

    \vspace{0.5cm}
    Oleh: \textbf{\@nama}

    NRP. \@NRP

    \vspace{0.5cm}
    Surabaya, ~\@bulan ~\@tahun

    Disetujui oleh Tim Penguji Tugas Akhir:

    \vspace{0.5cm}
    Pembimbing:

    \begin{enumerate}
        \ifthenelse{\equal{\pbb@tiga}{}\OR\equal{\pbb@tiga}{ }}
            {
                \ifthenelse{\equal{\pbb@dua}{}\OR\equal{\pbb@dua}{ }}
                    {
                        \item \pbb@satu ~ \hfill (\fillin[3cm])

                        NIP. ~\NIP@pbb@satu
                    }
                    {
                        \item \pbb@satu ~ \hfill (\fillin[3cm])

                        NIP. ~\NIP@pbb@satu
                        \item \pbb@dua ~ \hfill (\fillin[3cm])

                        NIP. ~\NIP@pbb@dua
                    }
            }
            {
                \item \pbb@satu ~ \hfill (\fillin[3cm])

                NIP. ~\NIP@pbb@satu
                \item \pbb@dua ~ \hfill (\fillin[3cm])

                NIP. ~\NIP@pbb@dua
                \item \pbb@tiga ~ \hfill (\fillin[3cm])

                NIP. ~\NIP@pbb@tiga
            }
    \end{enumerate}

    \vspace{0.5cm}
    Penguji:
    \begin{enumerate}
        \ifthenelse{\equal{\uji@tiga}{}\OR\equal{\uji@tiga}{ }}
            {
                \ifthenelse{\equal{\uji@dua}{}\OR\equal{\uji@dua}{ }}
                    {
                        \item \uji@satu ~ \hfill (\fillin[3cm])

                        NIP. ~\NIP@uji@satu
                    }
                    {
                        \item \uji@satu ~ \hfill (\fillin[3cm])

                        NIP. ~\NIP@uji@satu
                        \item \uji@dua ~ \hfill (\fillin[3cm])

                        NIP. ~\NIP@uji@dua
                    }
            }
            {
                \item \uji@satu ~ \hfill (\fillin[3cm])

                NIP. ~\NIP@uji@satu
                \item \uji@dua ~ \hfill (\fillin[3cm])

                NIP. ~\NIP@uji@dua
                \item \uji@tiga ~ \hfill (\fillin[3cm])

                NIP. ~\NIP@uji@tiga
            }
    \end{enumerate}

    \begin{flushright}
        \begin{minipage}{.4\textwidth}
            \begin{center}
                Mengetahui,

                Kepala Departemen \@departemen,

                Fakultas \@fakultas,

                \vspace{2cm}
                \underline{\@kadep}

                NIP. \NIP@kadep
            \end{center}
        \end{minipage}
    \end{flushright}
\end{center}
\makeatother
\makeatletter
\chapter*{Pernyataan Orisinalitas}

    \vspace{2cm}
    \noindent
    Yang bertanda tangan di sini:
    
    \vspace{1cm}
    \noindent
    \begin{tabular}{lcp{0.49\paperwidth}}
        Nama Mahasiswa / NRP & : & \@nama~/ \@NRP \\
        Departemen           & : & \@departemen   \\
        Dosen Pembimbing     & : & \ifthenelse{\equal{\pbb@tiga}{}\OR\equal{\pbb@tiga}{ }}
                                        {
                                            \ifthenelse{\equal{\pbb@dua}{}\OR\equal{\pbb@dua}{ }}
                                                {
                                                    \pbb@satu~/\protect\\&&\NIP@pbb@satu\\
                                                }
                                                {
                                                    \noindent
                                                       1. \pbb@satu \protect\\&&\hfill~/~\NIP@pbb@satu \\
                                                    && 2. \pbb@dua  \protect\\&&\hfill~/~\NIP@pbb@dua  \\
                                                }
                                        }
                                        {
                                            \noindent
                                               1. \pbb@satu \protect\\&&\hfill~/~\NIP@pbb@satu \\
                                            && 2. \pbb@dua  \protect\\&&\hfill~/~\NIP@pbb@dua  \\
                                            && 3. \pbb@tiga \protect\\&&\hfill~/~\NIP@pbb@tiga \\
                                        }
    \end{tabular}

    \vspace{1cm}
    \noindent
    dengan ini menyatakan bahwa Tugas Akhir dengan judul ``\textbf{\MakeUppercase{\@judul}}'' adalah hasil karya sendiri, bersifat orisinal, dan ditulis dengan mengikuti kaidah penulisan ilmiah.

    \vspace{5mm}
    \noindent
    Bilamana di kemudian hari ditemukan ketidaksesuaian dengan pernyataan ini, maka saya bersedia menerima sanksi sesuai dengan ketentuan yang berlaku di Institut Teknologi Sepuluh Nopember.

    \vspace{1cm}
    \begin{flushright}
        Surabaya, ~\@tanggal ~\@bulan ~\@tahun
    \end{flushright}

    \noindent Mengetahui,

    \vspace{1mm}
\begin{center}
    \begin{minipage}{.4\textwidth}
        \centering
        Dosen Pembimbing,

        \vspace{3cm}

        \underline{\pbb@satu}

        NIP. \NIP@pbb@satu
    \end{minipage}
    \hfill
    \begin{minipage}{.4\textwidth}
        \centering
        Mahasiswa,

        \vspace{3cm}

        \underline{\@nama}

        NRP. \@NRP
    \end{minipage}
\end{center}

\makeatother
    \chapter*{Abstrak}

\makeatletter
    \vspace{1cm}
    \begin{center}
        \noindent\textbf{\@judul}
    \end{center}

    \vspace{1cm}
    \begin{table}[htbp]
        \begin{tabular}{lcp{0.49\paperwidth}}
            Nama Mahasiswa / NRP & : & \@nama~/ \@NRP\\
            Departemen & : & \@departemen~{\singkatan@fakultas}-ITS\\
            Dosen Pembimbing & : & \pbb@satu
        \end{tabular}
    \end{table}

    \vspace{1cm}
    \noindent
    \textbf{Abstrak}

     Industri telekomunikasi memiliki tantangan dalam memprediksi pendapatan (\textit{revenue}) layanan SMS internasional yang memiliki karakteristik unik dan dipengaruhi oleh faktor eksternal. Meskipun pendekatan \textit{machine learning} telah banyak digunakan untuk prediksi \textit{revenue}, studi yang mengeksplorasi penggunaan Algoritma Genetika (\textit{Genetic Algorithm}) sebagai teknik optimasi hiperparameter pada model \textit{machine learning}, khususnya dalam konteks prediksi \textit{revenue} layanan SMS internasional, masih relatif terbatas. Penelitian Tugas Akhir ini mengimplementasikan Algoritma Genetika untuk optimasi hiperparameter pada model \textit{machine learning} khususnya \textit{tree based models} untuk mengatasi permasalahan prediksi \textit{revenue} layanan SMS internasional pada perusahaan X. Dataset primer yang digunakan berasal dari trafik SMS internasional perusahaan X pada tahun 2023 sampai 2024. Eksperimen dilakukan dengan membandingkan metode optimasi hiperparameter menggunakan Algoritma Genetika dengan tanpa menggunakan optimasi pada model machine learning berbasis pohon (\textit{tree-based models}). Hasil penelitian ini menunjukkan bahwa penerapan Algoritma Genetika untuk optimasi hiperparameter mampu untuk meningkatkan performansi \textit{tree based models} dalam memprediksi \textit{revenue} dengan WMAPE untuk \textit{XGBoost + GA} sebesar 16.29 \% dan \textit{Random Forest + GA} sebesar 20.97 \%  dibandingkan dengan \textit{baseline} yaitu WMAPE \textit{Random Forest} sebesar 25.37 \% dan WMAPE \textit{XGBoost} sebesar 27.12 \%.
    
    % Untuk kata kunci
    \katakunci{Pendapatan Layanan SMS Internasional, Tree-based Models, Optimasi Hyperparameter, Algoritma Genetika}
    
\makeatother
    \chapter*{Abstract}

\makeatletter
    \vspace{1cm}
    \begin{center}
        \noindent\textbf{\@title@}
    \end{center}

    \vspace{1cm}
    \begin{table}[htbp]
        \begin{tabular}{lcp{0.49\paperwidth}}
            Student Name / ID & : & \@nama~/ \@NRP\\
            Department & : & \@department~{\faculty@abbreviation}-ITS\\
            Advisor & : & \pbb@satu
        \end{tabular}
    \end{table}

    \vspace{1cm}
    \noindent
    \textbf{Abstract}

    The telecommunications industry faces challenges in predicting the revenue of international SMS services, which have unique characteristics and are influenced by external factors. Although machine learning approaches have been widely used for revenue prediction, studies exploring the use of Genetic Algorithm as a hyperparameter optimization technique in machine learning models, particularly in the context of international SMS service revenue prediction, remain relatively limited. This research implements Genetic Algorithm for hyperparameter optimization on machine learning models, specifically tree-based models, to address the problem of predicting international SMS service revenue at Company X. The primary dataset used is derived from international SMS traffic data of Company X covering the period of 2023 to 2024. Experiments were conducted by comparing hyperparameter optimization using Genetic Algorithm against baseline models without optimization on tree-based machine learning models. The results demonstrate that the application of Genetic Algorithm for hyperparameter optimization is capable of improving the performance of tree-based models in predicting revenue, achieving a WMAPE of 16.29\% for XGBoost + GA and 20.97\% for Random Forest + GA, compared to the baseline WMAPE of 25.37\% for Random Forest and 27.12\% for XGBoost.
    
    % For keywords
    \keywords{Revenue of SMS International Services, Tree-based Models, Hyperparameter Optimization, Genetic Algorithm}
    
\makeatother
    \chapter*{Kata Pengantar}
Puji syukur saya haturkan ke hadirat Allah Swt. karena berkat rahmat dan anugerah-Nya penulis dapat menyelesaikan \ta \ dengan judul \textbf{Optimasi Hyperparameter pada \textit{Tree-based Models} Menggunakan  Algoritma Genetika untuk Prediksi Pendapatan Layanan SMS Internasional}. 
Tugas Akhir ini disusun guna memenuhi salah satu persyaratan akademik untuk menyelesaikan Program Studi Sarjana di Departemen Matematika, Fakultas Sains dan Analitika Data, Institut Teknologi Sepuluh Nopember.

\TA~ini tidak terlepas dari dukungan, kerja sama, dan bantuan dari berbagai pihak. 
Oleh karena itu, penulis menyampaikan rasa terima kasih yang mendalam kepada: 

\begin{enumerate}
    \item Kedua orang tua, kakak, dan keluarga yang senantiasa memberikan motivasi dan doa selama proses perkuliahan hingga penyusunan Tugas Akhir ini. 
    \item Bapak Dr. Drs. Bandung Arry S., MI.Komp. selaku dosen pembimbing Tugas Akhir yang telah memberikan bimbingan dan arahan selama penyusunan Tugas Akhir ini. 
    \item Bapak Dr. Drs. Bandung Arry S., MI.Komp, Muhammad Yusuf Hakim Widianto, S.Si., M. Si., M.Sc., Ph.D, dan rof. Subchan, S.Si, M.Sc, Ph.D sebagai dosen penguji Tugas Akhir yang telah memberikan masukan dan saran guna meningkatkan kualitas Tugas Akhir ini. 
    \item Drs. Lukman Hanafi, M.Sc selaku dosen wali yang telah membimbing dan memberikan dukungan akademik selama perkuliahan. 
    \item Teman-teman seperjuangan mahasiswa yang telah memberikan dukungan dan motivasi selama perkuliahan. 
\end{enumerate}

Penulis menyadari bahwa Tugas Akhir ini masih memiliki banyak kekurangan. 
Kritik dan saran akan sangat diharapkan guna penyempurnaan penulisan di masa depan.
Semoga Tugas Akhir ini dapat menebar kebermanfaatan bagi semua pihak. 


\makeatletter
\begin{flushright}
    Surabaya, \@tanggal \@bulan \@tahun

    \vspace{2cm}
    \@nama
\end{flushright}
\makeatother
    \newpage
    \listoffigures
    \chapter*{Daftar Simbol}

\begin{table}[htbp]
    \begin{tabular}{lcl}
         $R_t$      & : & \textit{Return} emas pada waktu ke-$t$ \\[4pt]
        $P_t$       & : & Harga emas pada waktu ke-$t$ \\[4pt]
        $P_{t-1}$   & : & Harga emas pada waktu sebelum $t$ \\[4pt]
        $r$         & : & Jangkauan data \\[4pt]
        $x_{min}$   & : & Nilai minimum \\[4pt]
        $x_{maks}$  & : & Nilai maksimum \\[4pt]
        $l$         & : & Lebar masing-masing kelas \\[4pt]
        $k$         & : & Jumlah kelas \\[4pt]
        $\Pr$       & : & Operator probabilitas \\[4pt]
        $X_n$       & : & Kondisi saat $t = n$ \\[4pt]
        $P_{ij}$    & : & Probabilitas transisi dari \textit{state} $i$ menuju \textit{state} $j$ \\[4pt]
        $n_{ij}$    & : & Jumlah transisi dari \textit{state} $i$ menuju \textit{state} $j$ \\[4pt]
        $N_i$       & : & Total transisi yang berasal dari \textit{state} $i$ \\[4pt]
        $S$         & : & Himpunan \textit{state} pada DTMC \\[4pt]
        $\mathbb{N}$ & : & Himpunan bilangan asli \\[4pt]
        $\square$   & : & Operator temporal \textit{always} \\[4pt]
        $\Diamond$  & : & Operator temporal \textit{eventually} \\[4pt]
        $s_0$       & : & \textit{State} awal pada model DTMC \\[4pt]
        $AP$        & : & Himpunan proposisi atomik \\[4pt]
        $L$         & : & Fungsi pelabelan \textit{state} dalam DTMC \\[4pt]
        $\Phi$      & : & Formula \textit{state} dalam logika temporal \\[4pt]
        $\phi$      & : & Formula \textit{path} dalam logika temporal \\[4pt]
        $U$         & : & Operator \textit{until} pada LTL/CTL/PCTL \\[4pt]
        $U_{\le n}$ & : & Batas waktu maksimal $n$ untuk operator hingga (\textit{until}) \\[4pt]
        $P_J(\phi)$  & : & Operator probabilitas PCTL dengan interval $J$ \\[4pt]
    \end{tabular}
\end{table}
    \newpage
    \tableofcontents
    \newpage
    \setstretch{1.2}
    \pagenumbering{arabic}
    \setcounter{page}{1}
    \chapter{Pendahuluan}\label{ch:introduction}
    
    \section{Latar Belakang Masalah}\label{sec:background}

    Industri telekomunikasi global, khususnya pada segmen \textit{wholesale}, mengandalkan layanan pesan internasional seperti SMS \textit{Application-to-Person} (A2P) sebagai salah satu sumber pendapatan (\textit{revenue}) yang signifikan. Layanan SMS internasional memiliki karakteristik unik karena melibatkan banyak operator lintas negara, rute interkoneksi yang beragam, serta dipengaruhi oleh berbagai faktor eksternal seperti regulasi antarnegara, fluktuasi nilai tukar mata uang, dan pola musiman yang berbeda di setiap wilayah geografis \cite{10.1145/3690771.3690784}. Sebagai salah satu penyedia layanan \textit{wholesale} internasional, perusahaan X menghadapi kompleksitas serupa dalam mengelola pendapatan dari layanan ini.

    Tantangan utama yang dihadapi perusahaan X sebagai operator telekomunikasi saat ini adalah meningkatnya volume lalu lintas (\textit{traffic}) yang tidak selalu diiringi dengan peningkatan pendapatan. Fenomena ini disebabkan oleh kebocoran pendapatan (\textit{revenue leakage}) yang dipicu oleh berkembangnya teknologi alternatif, seperti \textit{Flash Calls} yang melewati (\textit{bypass}) jalur SMS tradisional, \textit{gray routes}, serta mode otentikasi alternatif yang lebih murah seperti VOTP \cite{telin2026leakage}. Berdasarkan laporan Mobilesquared dan VOX Carrier, akumulasi kebocoran pendapatan akibat \textit{gray routes} pada periode 2020--2024 mencapai USD 37,1 miliar, atau rata-rata USD 7,69 miliar per tahun \cite{mobilesquared2021}. Kondisi ini menunjukkan bahwa kontrol perusahaan penyedia layanan \textit{wholesale} internasional seperti perusahaan X semakin menurun sehingga pendapatan kedepan menjadi semakin sulit diprediksi.
    
    Prediksi \textit{revenue} yang akurat memiliki peranan penting bagi operator telekomunikasi untuk mendukung perencanaan operasional, alokasi kapasitas jaringan, penetapan strategi harga, dan pengambilan keputusan strategis jangka pendek maupun jangka panjang \cite{DiFrancescomarino2018GAHyperparameter}. Prediksi yang akurat dapat berfungsi sebagai instrumen deteksi dini, di mana deviasi yang signifikan antara nilai prediksi dan realisasi pendapatan dapat menjadi indikasi awal terjadinya penyimpangan atau kebocoran pendapatan. Pendapatan layanan telekomunikasi, termasuk layanan SMS internasional, dipengaruhi oleh berbagai faktor seperti volume trafik, negara tujuan, jenis layanan, serta karakteristik pelanggan. Hubungan antara  faktor-faktor tersebut bersifat non-linear dan dinamis, sehingga sulit jika menggunakan pendekatan statistik konvensional seperti \textit{Autoregressive Integrated Moving Average} (ARIMA) yang mengasumsikan hubungan linear antarvariabel \cite{ZHANG2003159}.
    
    Machine learning merupakan pendekatan yang efektif untuk menangani masalah prediksi pendapatan dengan kemampuannya menangkap pola non-linear dan hubungan kompleks dalam data historis \cite{DiFrancescomarino2018GAHyperparameter, Wu2022ParallelGA}. Model-model berbasis deep learning seperti Long Short-Term Memory (LSTM) dan Gated Recurrent Unit (GRU), serta model \textit{ensemble} seperti XGBoost dan Random Forest, telah banyak diaplikasikan dalam berbagai domain \textit{forecasting} termasuk telekomunikasi \cite{Wu2022ParallelGA, mehdary2024}. Meskipun model berbasis sekuensial seperti LSTM menunjukkan performa yang baik pada data deret waktu, model berbasis pohon (\textit{tree-based}) tetap menjadi pilihan populer untuk data tabular karena kemampuannya menangkap hubungan non-linear, robust terhadap \textit{outlier}, serta tidak memerlukan asumsi distribusi data tertentu \cite{grinsztajn2022treebasedmodelsoutperformdeep, Chen2016}. Selain itu, model \textit{tree-based} memiliki ruang hiperparameter yang relatif luas dan sensitif terhadap pengaturan parameter, sehingga kinerjanya sangat dipengaruhi oleh proses optimasi hiperparameter yang digunakan \cite{Probst2019Hyperparameters}.

    Optimasi hiperparameter menjadi tahapan penting dalam pengembangan model \textit{machine learning} yang menentukan kualitas akhir dari model prediksi \cite{Morales2022SurveyMOHPO}. Metode konvensional seperti \textit{Grid Search} dan \textit{Random Search} memiliki keterbatasan signifikan dalam hal efisiensi komputasi, terutama ketika menghadapi ruang pencarian hiperparameter yang besar dan berdimensi tinggi. \textit{Grid Search} melakukan pencarian \textit{exhaustive} (mendalam) yang sangat mahal secara komputasi, sementara \textit{Random Search} meskipun lebih efisien, cenderung melewatkan area optimal dalam ruang pencarian tanpa strategi eksploitasi yang terstruktur \cite{Le2023PINNGA}.
    
    Algoritma Genetika (\textit{Genetic Algorithm}/GA) muncul sebagai alternatif yang menjanjikan untuk mengatasi keterbatasan metode optimasi konvensional \cite{Ghatasheh2022ModifiedGA,Deighan2021GANNGravWave}. Keunggulan GA terletak pada kemampuannya melakukan eksplorasi luas dalam ruang pencarian sambil tetap melakukan eksploitasi terhadap solusi-solusi yang menjanjikan, serta kemampuannya menangani optimasi multi-objektif yang mempertimbangkan \textit{trade-off} (kompromi) antara akurasi model dan kompleksitas komputasi \cite{Kobayashi2024DQN_BRKGA, Truong2021TelecomRevenue}. Penelitian-penelitian terkini menunjukkan keberhasilan penerapan GA untuk optimasi hiperparameter di berbagai domain aplikasi. Mehdary et al. (2024) mendemonstrasikan bahwa GA dapat meningkatkan akurasi model XGBoost dari 0.82 menjadi 0.978 dalam aplikasi deteksi \textit{fraud} pada smart grid \cite{Mehdary2024GA_XGBoost}. Le et al. (2023) menjelaskan bahwa GA dapat menghemat hingga 88.03\% biaya komputasi dibanding \textit{Grid Search} dan 39.82\% dibanding \textit{Random Search}, serta mengurangi \textit{test error} rata-rata sebesar 37.40\% pada eksperimen \textit{Physics-Informed Neural Networks} (PINNs) \cite{Le2023PINNGA}.
    
    Akan tetapi penelitian sebelumnya belum mengeksplorasi penerapan GA untuk optimasi hiperparameter dan prediksi pendapatan. layanan SMS internasional yang memiliki karakteristik unik seperti  pola temporal yang kompleks serta hubungan antara \textit{supplier} dan \textit{client}. Penelitian ini bertujuan untuk mengisi kesenjangan tersebut dengan mengembangkan dan mengevaluasi penerapan Algoritma Genetika untuk optimasi hiperparameter pada \textit{tree-based models} untuk prediksi total \textit{revenue} layanan SMS internasional di Perusahaan X. Penelitian ini diharapkan dapat memberikan bukti empiris tentang efektivitas GA dalam meningkatkan akurasi prediksi pendapatan (\textit{revenue}) dibandingkan dengan tanpa menggunakan optimasi, khususnya dalam konteks layanan telekomunikasi internasional yang memiliki karakteristik data kompleks.
    
    \section{Rumusan Masalah}\label{sec:problem}

    Berdasarkan latar belakang yang telah dijelaskan, terdapat beberapa permasalahan utama yang menjadi fokus penelitian ini, yaitu :
        % Rumusan Masalah pada Penelitian Tugas Akhir
        \begin{enumerate}
 
            \item Bagaimana penerapan Algoritma Genetika dapat digunakan untuk mengoptimalkan hiperparameter model prediksi?

            \item Bagaimana mengevaluasi \textit{tree-based models} dengan optimasi hiperparameter menggunakan Algoritma Genetika?
        \end{enumerate}
        % 

    \section{Batasan Masalah}\label{sec:limitation}
    Agar penelitian ini terarah dan fokus, ruang lingkup penelitian dibatasi sebagai berikut:
        % Batasan Masalah pada Penelitian Tugas Akhir
        \begin{enumerate}
            \item Objek penelitian terbatas pada data historis layanan komunikasi internasional  variabel \textit{international charge} dan komponen biaya lainnya sebagai indikator \textit{revenue} (pendapatan).
            \item Model prediksi yang digunakan dalam penelitian ini dibatasi pada algoritma  machine learning berbasis pohon. (\textit{tree-based models}) yaitu Random Forest dan XGBoost yang performanya dioptimalkan menggunakan GA.
            % \item Evaluasi metode optimasi hyperparameter hanya dibandingkan dengan metode optimasi populer konvensional yaitu \textit{Random Search} dan \textit{Grid Search}.
        \end{enumerate}
        % 

    \section{Tujuan}\label{sec:purpose}
    Menjawab dari rumusan masalah, tujuan dari penelitian ini adalah sebagai berikut : 
        % Tujuan Penelitian Tugas Akhir
        \begin{enumerate}
            
            \item Menerapkan Algoritma Genetika untuk optimasi hiperparameter model prediksi \textit{tree-based}.
            
            \item Melakukan evaluasi optimasi hiperparameter menggunakan Algoritma Genetika pada model prediksi \textit{tree-based}.
        \end{enumerate}
        % 

    \section{Manfaat}\label{sec:benefit}
    Manfaat dari penelitian ini dapat dibagi menjadi berikut :
        % Manfaat Penelitian Tugas Akhir
        \begin{enumerate}
                \item Memberikan sistem prediksi pendapatan yang lebih akurat untuk mendukung pengambilan keputusan strategis perusahaan dalam pengelolaan layanan trafik SMS internasional.
                \item Membantu perusahaan dalam mengefisienkan biaya untuk operasional berdasarkan prediksi pendapatan dari tiap pola data.
        \end{enumerate}
        % 
    \chapter{Tinjauan Pustaka}\label{ch:literature}

   \section{Prediksi \textit{Revenue} pada Industri Telekomunikasi}
\label{sec:revenuetelco}

    Layanan SMS internasional merupakan bagian dari ekosistem layanan telekomunikasi lintas negara yang memungkinkan pengiriman pesan singkat antar pengguna di berbagai negara. Layanan ini melibatkan rantai distribusi yang kompleks, mulai dari penyedia layanan di negara asal, \textit{intermediate carrier}, hingga \textit{network operator} di negara tujuan. Kompleksitas rantai distribusi tersebut menyebabkan adanya variasi signifikan pada struktur biaya operasional dan pendapatan yang dihasilkan, terutama pada komponen \textit{international charge} \cite{Truong2021TelecomRevenue}.
    
    Dalam konteks industri telekomunikasi, \textit{revenue} didefinisikan sebagai pendapatan yang diperoleh perusahaan dari penyediaan layanan komunikasi kepada pelanggan. Pada layanan SMS internasional, pendapatan tersebut umumnya direpresentasikan melalui komponen \textit{international 
    charge} yang dibayarkan berdasarkan volume pesan yang dikirimkan lintas negara. Variabel \textit{international charge} sering digunakan sebagai indikator utama \textit{revenue} karena secara langsung mencerminkan nilai transaksi SMS internasional \cite{Truong2021TelecomRevenue}.
    
    Meningkatnya adopsi teknologi alternatif seperti seperti \textit{Flash Calls} yang melewati (\textit{bypass}) jalur SMS tradisional, \textit{gray routes}, serta mode otentikasi alternatif yang lebih murah seperti VOTP \cite{telin2026leakage}. \textit{Flash Calls} merupakan metode otentikasi yang memanfaatkan panggilan telepon singkat yang langsung terputus sebagai pengganti kode OTP berbasis SMS, sehingga trafik SMS A2P untuk keperluan autentikasi beralih ke kanal suara yang tarifnya lebih murah \cite{telin2026leakage}. Kemudian, \textit{Voice OTP} (VOTP) adalah metode pengiriman kode otentikasi melalui panggilan suara otomatis sebagai alternatif SMS yang lebih murah bagi pengirim \cite{telin2026leakage}. Selanjutnya \textit{gray routes} merupakan jalur pengiriman SMS internasional yang memanfaatkan celah jaringan sinyal antara operator untuk meneruskan pesan A2P tanpa melalui kesepakatan interkoneksi resmi sehingga operator tujuan tidak memperoleh pembayaran atas trafik yang masuk. Kondisi ini mendorong perusahaan telekomunikasi untuk mengoptimalkan strategi operasional untuk mempertahankan stabilitas pendapatan.
    
    Pendapatan layanan telekomunikasi, termasuk layanan SMS internasional, dipengaruhi oleh berbagai faktor seperti volume trafik, negara tujuan, jenis layanan, serta karakteristik pelanggan. Hubungan antara faktor-faktor tersebut bersifat non-linear dan dinamis, sehingga 
    menyulitkan penerapan pendekatan statistik konvensional dalam proses prediksi \textit{revenue}. Studi yang dilakukan oleh Danilov dan 
    Maltseva menunjukkan bahwa pendekatan \textit{machine learning} untuk prediksi \textit{revenue} produk telekomunikasi mampu menghasilkan prediksi yang lebih akurat dibandingkan metode konvensional \cite{Danilov2024}. Oleh karena itu, pendekatan berbasis \textit{machine learning} menjadi solusi yang relevan untuk memodelkan pola kompleks dalam data historis trafik SMS internasional \cite{Chen2016}.
    
    Model prediksi \textit{revenue} bertujuan untuk memperkirakan nilai pendapatan di masa depan berdasarkan data historis trafik dan karakteristik layanan. Meskipun beberapa studi telah mengeksplorasi prediksi \textit{revenue} pada layanan telekomunikasi secara umum 
    \cite{Truong2021TelecomRevenue, Danilov2024} serta optimasi 
    hiperparameter berbasis Algoritma Genetika pada model \textit{machine learning} \cite{Ghatasheh2022ModifiedGA}, studi yang secara khusus mengkaji prediksi \textit{revenue} layanan SMS internasional masih sangat terbatas. Prediksi yang akurat dapat digunakan sebagai alat mitigasi risiko finansial dengan mengidentifikasi potensi fluktuasi pendapatan serta mendukung pengambilan keputusan strategis \cite{Oziri2023RevenueForecasting}.

    
    \section{Penelitian Terdahulu}\label{sec:pastresearch}
    Penerapan  algoritma genetika untuk optimasi hiperparameter dilakukan oleh \citeA{gour2024optimizing} dengan menggunakan model \textit{Neural Network} pada dataset diagnosis diabetes. Hasil dari penelitian ini menunjukkan GA konvergen ke nilai hiperparameter yang meningkatkan akurasi. Selain itu, \citeA{Ghatasheh2022ModifiedGA} menerapkan algoritma genetika untuk \textit{hyperparameter optimization} (HPO) pada model XGBoost pada kasus klasifikasi spam dengan \textit{cross-validation} berulang. Penelitian ini menunjukkan hasi G-mean sebesar 82.32 \% dan akurasi 92.67 \% serta mengurangi dimensi fitur kurang dari 10 \%. Selain itu \citeA{Kobayashi2024DQN_BRKGA} melakukan penelitian pendekatan Hybird Deep Q-Network dan Biased Random Key Genetic Algorithm (BRKGA) yang dilakukan pada 8 dateset beragram. \textit{Genetic Algorithm} (GA) terbukti dapat mengurangi biaya komputasi hingga 44 \% tanpa kehilangan kualitas solusi dibanding BRKGA tanpa HPO.
    
    Penerapan GA juga telah diuji secara langsung pada industri telekomunikasi. Misalnya, penelitian oleh \citeA{li2022} mengkombinasikan GA dengan algoritma XGBoost untuk mengoptimalkan parameter model prediksi \textit{churn} pelanggan. Penelitian tersebut menunjukkan bahwa pendekatan GA-XGBoost memberikan hasil yang lebih baik dibandingkan XGBoost tanpa optimasi dengan evaluasi nilai AUC dan F1-Score. Penelitian lain dilakukan oleh \citeA{aljanabi2022distributed} pada kasus \textit{Detection Denial of Service Attacks} (DDoS) pada \textit{Software Define Networks} (SFD). Studi tersebut menunjukkan GA mengoptimalkan hiperparameter model \textit{Decision Tree} sehingga mencapai akurasi 99.49 \% pada kasus klasifikasi \textit{traffic}. Berdasarkan berbagai penelitian terdahulu tersebut, dapat disimpulkan bahwa metode optimasi berbasis GA telah terbukti efektif untuk meningkatkan kinerja model \textit{machine learning} di berbagai domain, termasuk telekomunikasi. Namun, sebagian besar penelitian belum ditemukan aplikasi GA untuk optimasi hiperparameter \textit{tree-based models} untuk prediksi pendapatan layanan komunikasi internasional. Dengan demikian, untuk menjawab permasalahan tersebut penelitian ini menerapkan GA sebagai metode optimasi hiperparameter pada \textit{tree-based models} untuk memprediksi prediksi layanan komunikasi internasional. Pendekatan ini diharapkan dapat memberikan kontribusi baik dalam mengurangi \textit{error} atau kesalahan prediksi pendapatan di industri telekomunikasi serta mengevaluasi dampak optimasi terhadap akurasi prediksi dan efisiensi komputasi.

    \section {Konsep Dasar Klasifikasi}
    Klasifikasi merupakan metode \textit{machine learning} dalam data mining untuk memprediksi anggota dari tiap grup data. Meskipun banyak metode yang ada pada \textit{machine learning}, klasifikasi merupakan metode yang paling banyak digunakan \cite{inproceedings}. Klasifikasi merupakan metode yang paling banyak dipelajari dan diteliti oleh para \textit{researcher} di bidang data mining dan machine learning. Output pada klasifikasi berupa variable yang memiliki label kelas. Metode klasifikasi sendiri dibedakan menjadi dua yaitu \textit{binary classification} dan \textit{multiclass classification}. \textit{Binary classification} merupakan metode klasifikasi yang memprediksi dua kelas yaitu antara 0 dan 1. Sedangkan \textit{multiclass classification} merupakan metode yang memprediksi lebih dari dua kelas \cite{ng2023cs229}. Pada penelitian Tugas Akhir ini digunakan metode \textit{binary classification} karena tujuan klasifikasi untuk memprediksi transaksi yang mempunyai pendapatan nol (tidak aktif) dan transaksi yang mempunyai pendapatan tidak nol (aktif).
    
    \section {Konsep Dasar Regresi}
    Regresi merupakan metode statistika yang digunakan untuk memodelkan hubungan antara variabel dependen atau target dengan satu atau lebih variabel independen atau prediktor. Regresi memiliki tujuan utama untuk memahami bagaimana nilai variabel dependen berubah ketika salah satu variabel independen diubah, sementara variabel independen lainnya tetap konstan. Menurut James et al. (2023), regresi telah menjadi salah satu metode fundamental dalam statistika dan \textit{machine learning} yang digunakan secara luas dalam berbagai domain aplikasi, mulai dari ekonomi, bisnis, hingga ilmu sosial.
    
    \begin{equation} \label{eq:wmape}
        Y = f(x) + \varepsilon
    \end{equation}
     
   Secara matematis, model regresi dapat dinyatakan pada persamaan 2.1 di mana $Y$ adalah variabel dependen yang ingin diprediksi, $X$ adalah vektor variabel independen yang terdiri dari $X_1, X_2, \dots, X_p$, $f(X)$ adalah fungsi sistematis yang menggambarkan hubungan antara $X$ dan $Y$, dan $\varepsilon$ adalah komponen \textit{error} atau \textit{noise} acak yang merepresentasikan variabilitas dalam data yang tidak dapat dijelaskan oleh model. Pemodelan regresi berusaha untuk menemukan fungsi $f(X)$ yang paling baik merepresentasikan hubungan tersebut berdasarkan data observasi yang ingin diprediksi \cite{hastie2009elements}.
    
    Pemilihan model regresi yang tepat dan konfigurasi hiperparameter sangat mempengaruhi performa. Misalnya parameter seperti jumlah pohon (\textit{estimators}), kedalaman maksimum (\textit{max depth}), tingkat pembelajaran (\textit{learning rate}), atau ukuran sampel bawah (\textit{min samples leaf}) dapat berdampak besar terhadap keseimbangan bias dan \textit{varians}. Untuk itu, metode optimasi seperti GA ditempatkan sebagai layer pencarian hiperparameter di atas model regresi ini. 

    \section{Two Staged Models}
    \textit{Two-stage models} atau model dua tahap merupakan pendekatan pemodelan 
    prediktif yang membagi proses prediksi menjadi dua fase yang saling berurutan. Model awal menghasilkan informasi atau keluaran tertentu kemudian keluaran dari tahap 
    pertama tersebut digunakan sebagai masukan tambahan untuk model pada tahap 
    kedua \cite{en10101547}.
    
    Secara umum, struktur \textit{two-stage models} dapat direpresentasikan sebagai 
    berikut. Misalkan $\mathbf{x} \in \mathbb{R}^p$ adalah vektor fitur masukan, 
    maka tahap pertama menghasilkan keluaran $z = f_1(\mathbf{x})$ dan tahap kedua 
    menghasilkan prediksi akhir $\hat{y} = f_2(\mathbf{x}, z)$. Maka dari itu, 
    prediksi akhir secara keseluruhan dapat dituliskan pada persamaan 2.2.
    
    \begin{equation}
        \hat{y} = f_2\!\left(\mathbf{x},\; f_1(\mathbf{x})\right)
    \end{equation}
    
    Keunggulan utama pendekatan ini dibandingkan model satu tahap adalah kemampuannya 
    untuk memecah masalah yang kompleks menjadi dua sub-masalah yang lebih 
    sederhana dan lebih mudah dioptimalkan secara terpisah \cite{en10101547}. 
    Output dari model tahap pertama menjadi informasi yang berharga bagi model tahap kedua, sehingga model tahap kedua dapat menangkap pola data yang tidak dapat ditangkap jika hanya menggunakan fitur data tanpa melalui tahap model pertama. 
    
    \section{Tree-based Models}
    
    Model berbasis pohon (\textit{tree-based models}) merupakan salah satu pendekatan pembelajaran mesin yang sering digunakan karena kemampuannya dalam memodelkan hubungan non-linier serta hubungan kompleks antar fitur tanpa memerlukan asumsi distribusi data tertentu. Model ini bekerja dengan membagi ruang fitur ke dalam beberapa region menggunakan aturan pemisahan berbasis kriteria tertentu, seperti \textit{Gini impurity} atau \textit{information gain}, sehingga menghasilkan struktur hierarkis berbentuk pohon keputusan.
    
    \textit{Decision Tree} merupakan bentuk paling dasar dari model berbasis pohon. Model ini melakukan pemisahan data secara rekursif berdasarkan fitur yang paling informatif hingga mencapai kondisi berhenti tertentu. Meskipun mudah diinterpretasikan dan efisien secara komputasi, \textit{Decision Tree} memiliki kelemahan utama berupa kecenderungan terhadap \textit{overfitting}, terutama ketika pohon tumbuh terlalu dalam \cite{info14030140}.

    % \begin{figure}[H]
    % \centering
    %     \includegraphics[width=0.8\textwidth]{tugas-akhir/assets/images/dt.png}
    %     \caption {Model \textit{Decision Tree} dalam dunia medis}
    %     \text{\cite{Wang2023RUL}}
    %     \label{fig:revenue-framework}
    % \end{figure}

    
    \subsection{Random Forest}
    \textit{Random Forest} merupakan kombinasi \textit{multiple decision trees} untuk meningkatkan akurasi prediksi dan mengurangi \textit{overfitting} \cite{Breiman2001}. Algoritma \textit{Random Forest} berjalan dengan membentuk sekumpulan pohon keputusan yang dilatih menggunakan sample acak dari data latih dengan pengembalian (\textit{boostrap sampling}), serta memilih subset fitur secara acak pada setiap \textit{leaf} pohon. Proses ini menghasilkan variasi antar pohon yang  mengurangi korelasi internal, sehingga meningkatkan kemampuan generalisasi model. Hasil dari tiap pohon digabungkan dengan menggunakan \textit{vote majority} dalam menghaasilkan hasil klasifikasi atau rata-rata prediksi untuk regresi yang lebih stabil dan akurat \cite{MontesinosLópez2022}. Prediksi akhir dari \textit{Random Forest} untuk masalah regresi dihitung sebagai rata-rata dari prediksi seluruh \textit{trees} dalam \textit{ensemble} yang dapat dilihat pada persamaan 2.2.
    \begin{equation}
        \hat{y} = \frac{1}{B} \sum_{b=1}^{B} f_b(x)
    \end{equation}
    di mana $B$ adalah jumlah \textit{tree} dalam \textit{ensemble} dan $f_b(x)$ adalah prediksi dari \textit{tree} ke-$b$ untuk input $x$. Kelebihan \textit{Random Forest} adalah kemampuan menangani \textit{overfitting} dengan baik, \textit{robust} terhadap \textit{outlier}, dan dapat memberikan estimasi \textit{feature importance} yang berguna untuk interpretasi model.

    % \begin{figure}[H]
    % \centering
    %     \includegraphics[width=0.8\textwidth]{tugas-akhir/assets/images/RandomForest.png}
    %     \caption {Struktur \textit{Random Forest}}
    %     \text{\cite{Wang2023RUL}}
    %     \label{fig:revenue-framework}
    % \end{figure}
    
    
    \subsection{XGBoost (Extreme Gradient Boosting)}
    XGBoost merupakan kombinasi \textit{decision tree } yang mengimplementasikan \textit{gradient boosting} yang sangat teroptimasi. XGBoost pertama kali dikembangkan oleh Chen Tianqi dan Carlos pada tahun 2016 dan dilanjutkan oleh peneliti selanjutnya untuk melakukan penyempurnaan dari penelitian sebelumnya \cite{10.1371/journal.pone.0283452}. 
    
    
    XGBoost merupakan model yang menggunakan pendekatan \textit{boosting}, yaitu membangun model secara bertahap di mana setiap model baru berusaha belajar dari kesalahan yang dibuat dari model sebelumnya. Berbeda dengan metode \textit{bagging}, seperti \textit{Random Forest},\textit{ boosting} bersifat sekuensial dan lebih fokus pada observasi yang susah diprediksi oleh model sebelumnya \cite{Dong2020}. Kelebihan XGBoost adalah penggunaan \textit{regularized objective function} yang membantu mengurangi \textit{overfitting} dan meningkatkan generalisasi model. Formula dari XGBoost dapat dilihat pada persamaaan 2.3.
    \begin{equation}
    \hat{y}_i^{(m)} = \sum_{k=1}^{m} f_k(x_i), \quad f_k \in F
    \tag{2.7}
    \end{equation}
    
    di mana $\hat{y}_i^{(m)}$ adalah prediksi pada iterasi ke-$m$ dengan $m$ adalah indeks
    dari jumlah iterasi, $f_k$ adalah fungsi pohon keputusan ke-$k$, dan $F$ adalah
    ruang fungsi semua pohon yang mungkin. Pelatihan XGBoost bertujuan untuk
    meminimalkan fungsi objektif pada persamaan (2.8):
    
    \begin{equation}
    \mathcal{L}^{(m)} = \sum_{i=1}^{n} l(y_i, \hat{y}_i^{(m)}) + \sum_{k=1}^{m} \Omega(f_k)
    \tag{2.8}
    \end{equation}
    
    di mana $\mathcal{L}^{(m)}$ adalah fungsi objektif yang diminimalkan pada iterasi
    \textit{boosting} ke-$m$, $l$ adalah fungsi \textit{loss} yang mengukur selisih antara
    nilai aktual dan prediksi, $\hat{y}_i$ adalah nilai prediksi model, dan $\Omega$
    adalah fungsi regularisasi seluruh pohon keputusan. fungsi regularisasi pada XGBoost dapat dituliskan lebih lanjut pada
    persamaan (2.9):
    
    \begin{equation}
    \Omega(f) = \gamma T + \frac{1}{2}\lambda \|w\|^2
    \tag{2.9}
    \end{equation}
    
    di mana $\gamma$ adalah parameter regularisasi struktural
    (\textit{tree complexity penalty}) yang berfungsi untuk mengontrol penalti terhadap jumlah daun,
    $T$ sebagai jumlah daun, dan $w$ sebagai bobot prediksi pada setiap \textit{leaf node}
    (\textit{leaf weight}) dalam pohon keputusan XGBoost \cite{Chen2016}.
    
    % \begin{figure}[H]
    % \centering
    %     \includegraphics[width=0.8\textwidth]{tugas-akhir/assets/images/XGBoost.png}
    %     \caption {Struktur \textit{XGBoost}}\text{\cite{Cheng2020SPSEXGBoost}}\label{fig:revenue-framework}
    % \end{figure}
    
    % \subsection{LightGBM (Light Gradient Boosting Machine)}
    % LightGBM dikembangkan oleh Microsoft Research dengan fitur utama \textit{Gradient-based One-Side Sampling} (GOSS) dan \textit{Exclusive Feature Bundling} (EFB) yang membuatnya sangat efisien untuk dataset berdimensi tinggi \cite{Ke2017}. Perbedaan dasar LightGBM adalah penggunaan strategi \textit{leaf-wise tree growth} yang diformulasikan pada persamaan 2.5.
    % \begin{equation}
    %     \Delta\text{loss} = \max\{\Delta\text{loss}(\text{leaf}_i)\}
    % \end{equation}
    % Strategi ini memilih \textit{leaf} dengan pengurangan \textit{loss} maksimum untuk di-\textit{split} selanjutnya. Berbeda dengan strategi \textit{level-wise} yang digunakan oleh sebagian besar implementasi GBDT, pendekatan \textit{leaf-wise} ini dapat mempercepat proses training hingga 20 kali lipat sambil mempertahankan akurasi yang sebanding \cite{Ke2017}.
    
    % \subsection{CatBoost (Categorical Boosting)}
    % CatBoost, dikembangkan oleh Yandex, menggunakan teknik \textit{ordered boosting} dan penanganan fitur kategorikal khusus untuk mengatasi \textit{prediction shift} yang disebabkan oleh \textit{target leakage}. Prediksi untuk sampel $i$ diformulasikan pada persamaan 2.6.
    % \begin{equation}
    %     \hat{y}_i = \sum_{t=1}^{T} \alpha_t \times \text{tree}_t(x_i; M_t^i)
    % \end{equation}
    % Untuk menangani fitur kategorikal, CatBoost menggunakan statistik target dengan formula pada persamaan 2.7.
    % \begin{equation}
    %     TS_k = \frac{\sum_{i:x_{i,k}=c} y_i + a p}{\sum_{i:x_{i,k}=c} 1 + a}
    % \end{equation}
    % di mana $a$ adalah parameter \textit{smoothing} dan $p$ adalah nilai prior. Teknik \textit{ordered boosting} yang diimplementasikan CatBoost secara efektif mengurangi \textit{overfitting} dibandingkan dengan metode boosting tradisional \cite{Prokhorenkova2018}.
        
    \section{Metode Optimasi hiperparameter}

    Dalam proses pelatihan model machine learning seperti proses optimasi hiperparameter  masih dibutuhkan. Banyak algoritma yang memiliki banyak hyperparamater sehingga proses optimasi hiperparameters menjadi lebih kompleks. Metode yang tidak terlalu kompleks adalah \textit{manual tuning} atau pencarian parameter secara manual. Akan tetapi, metode ini sangat bergantung kepada seberapa kompleks masalah yang ingin diselesaikan. Karena banyaknya kemungkinan yang dapat dihasilkan, solusi yang dihasilkan butuh untuk diskalasikan. Beberapa metode terkenal yang sering digunakan adalah \textit{grid search}, \textit{random search}, \textit{bayesian optimization}, serta pendekatan secara heuristik seperti Algoritma Genetika \cite{10.1145/3297067.3297080}.

    Masalah utama dalam optimasi hiperparameter model \textit{machine learning} adalah efisiensi waktu dan ruang pencarian. Optimasi Bayesian meningkatkan efisiensi waktu metode pencarian grid, tetapi metode ini bekerja dengan ruang pencarian yang terbatas seperti metode pencarian grid begitupun dengan random search. Di sisi lain, GA telah mengatasi kedua masalah tersebut, yaitu dengan mencari di berbagai ruang yang lebih luas dengan waktu yang lebih efisien \cite{Katoch2021}. Maka dari itu, GA membantu model \textit{TreeBased Models} seperti XGBoost memiliki performa yang lebih baik. GA membantu mengevaluasi berbagai solusi yang lebih luas dengan cepat untuk menemukan pilihan terbaik. Hai ini sangat penting untuk merancang sistem pendukung keputusan klinis (\textit{Decision Support System}) dan analisis big data \cite{https://doi.org/10.1155/2023/9673395}.
    
    % \subsection{Optimasi \textit{Grid Search}}
    % Strategi yang dikenal sebagai grid search mengotomatisasi metode \textit{trial-and-error} dengan menyediakan daftar nilai parameter yang mungkin untuk dipilih. Strategi ini terkenal penggunaannya dalam machine learning untuk mengidentifikasi pengaturan parameter yang optimal bagi model seperti yang ditunjukkan pada gambar 2.5.

    % \begin{figure}[H]
    % \centering
    %     \includegraphics[width=0.5\textwidth]{tugas-akhir/assets/images/Grid Search.png}
    %     \caption {Metode \textit{Grid Search}}
    %     \text{\cite{Shanthi2023}}
    %     \label{fig:revenue-framework}
    % \end{figure}
    
    % Ketika mencoba memilih model terbaik untuk semua kemungkinan pengaturan hiperparameter, dimulai dengan membangun model untuk setiap kemungkinan konfigurasi nilai hiperparameter, kemudian mengevaluasi setiap model dan memilih satu yang berkinerja terbaik. Untuk membangun \textit{grid variabel hiperparameter}, nilai grid dibuat terpisah untuk setiap kombinasi hiperparameter dan kemudian melatih model pada data validasi untuk setiap nilai grid.
    
    % Strategi ini mencoba mencari setiap konfigurasi hiperparameter yang mungkin, yang membuat metode ini menjadi pendekatan yang tidak efisien. Kelemahan lain dari \textit{grid search} adalah metode ini mungkin mengalami masalah dimensi ketika jumlah hiperparameter tumbuh secara eksponensial. Namun, tidak ada jaminan bahwa pencarian ini akan selalu memberikan jawaban yang sempurna

    % \subsection{Optimasi \textit{Random Search}}
    % Metode \textit{random search} bekerja dengan mengambil sampel nilai secara  acak untuk memperkirakan fungsi objektif. Metode ini efisien karena tidak mengasumsikan apa pun tentang struktur fungsi target. Metode ini sangat berguna ketika terdapat sejumlah besar pengetahuan domain yang harus dihadapi, yang memungkinkan munculnya solusi yang berlawanan dengan intuisi. Dibandingkan dengan \textit{grid search}, random search jauh lebih efektif dalam mengidentifikasi nilai terbaik seperti yang ditunjukkan pada gambar 2.6.

    % \begin{figure}[H]
    % \centering
    %     \includegraphics[width=0.5\textwidth]{tugas-akhir/assets/images/Random Search.png}
    %     \caption {Metode \textit{Random Search}}
    %     \text{\cite{Shanthi2023}}
    %     \label{fig:revenue-framework}
    % \end{figure}
    
    % Langkah-langkah dari metode ini meliputi penentuan ruang lingkup hiperparameter kemudian pemilihan kombinasi secara acak untuk melatih dan menguji model. Akan tetapi, ketika hanya sedikit hiperparameter yang memengaruhi kinerja algoritma machine learning, \textit{grid search} dapat mengungguli metode \textit{random search}. 
    
    % Ketika mencari sesuatu yang kompleks, seperti data yang bising (noisy) atau terputus-putus, random search sangat efisien dan menggunakan memori minimal. Karena itu, sampel input yang besar mungkin menjadi pilihan yang lebih baik. Setiap sampel bersifat unik, yang berarti jika sampel harus dianalisis secara paralel untuk mempercepat proses,  penemuan parameter optimal tetap efektif tanpa mengurangi keandalan hasil temuan.\cite{Shanthi2023}.

    \subsection{Algortima Genetika}\label{sec:ga}
    
    Algoritma Genetika (Genetic Algorithm) merupakan salah satu algoritma optimasi metaheuristik berbasis populasi yang terinspirasi dari teori seleksi alam. Populasi baru dihasilkan dengan berulang kali menggunakan operator genetik pada setiap individu dalam populasi. Istilah-istilah penting dari algoritma ini adalah pembentukan kromosom, inisialisasi populasi, seleksi (selection), \textit{crossover}, mutasi, dan fungsi \textit{fitness}. \shortciteA{Katoch2021} menjelaskan bahwa proses \textit{Genetic Algorithm} (GA) dimulai dengan pembentukan pulasi awal Y yang terdiri dari n kromosom secara acak, di mana n adalah ukuran populasi yang ditentukan sebelumnya. Setiap kromosom merepresentasikan satu kandidat solusi yang terdiri dari sejumlah gen sebanyak hiperparameter yang ingin dioptimasi. Dua kromosom yaitu C1 dan C2, dipilih dari populasi berdasarkan nilai \textit{fitnessnya}. C1 dan C2 akan menghasilkan keturunan baru O dengan melalui proses \textit{crossover}. Probabilitas pada tahap ini dinamakan $P_c$, yang merupakan parameter probabilitas \textit{crossover}. Operator mutasi genetik dengan probabilitas $P_m$ kemudian diterapkan pada O untuk menghasilkan anggota baru O’. Anggota O’ ditambahkan ke populasi sebelumnya untuk membentuk populasi baru. Proses seleksi, crossover, dan mutasi berlanjut hingga seluruh populasi dengan nilai \textit{fungsi fitness} yang optimal dihasilkan. Probabilitas crossover dan mutasi inilah yang menyebabkan GA dapat secara dinamis mencari solusi optimal.


    \section{Penerapan Algoitma Genetika untuk Optimasi hiperparameter \textit{Tree-Based Models}}

    Penerapan algoritma genetika untuk optimasi hiperparameter pada \textit{tree-based Models} dilakukan dengan mencari parameter optimal untuk meminimalkan fungsi \textit{fitness} yang digunakan untuk mengevaluasi kinerja model.  Tahapan algoritma genetika yang diimplementasikan pada penelitian ini sebagai berikut:  
    
    \begin{algorithm}[H]
    \caption{Tahapan Algoritma Genetika \protect\cite{https://doi.org/10.1155/2023/9673395}}
    \label{alg:ga}
    \begin{algorithmic}[1]
    \State \textbf{Input} : Dataset $D$, Jumlah Generasi $G$, Ukuran Populasi $P$, Probabilitas Crossover ($Pc$). dan Probabilitas Mutasi, ($Pm$)
    \State \textbf{Output} : Nilai optimal hiperparameter \textit{tree-based models}
    \State $g = 0$
    \State Insialisasi populasi secara acak dengan ukuran $P$
    \State Urutkan $D$ secara kronologis berdasarkan $\textit{period} = \textit{Year} \times 12 + \textit{Month}$
    \State Tetapkan nilai batas: $c = \text{quantile}_{0{,}80}\bigl(\{\textit{period}_i\}\bigr)$
    \State $D_{\text{train}} \gets \{(x_i, y_i) \in D \mid \textit{period}_i \leq c\}$
    \State $D_{\text{test}}  \gets \{(x_i, y_i) \in D \mid \textit{period}_i > c\}$
    \State Inisialisasi populasi $\mathcal{P}_0 = \{\boldsymbol{\theta}_1, \boldsymbol{\theta}_2, \ldots, \boldsymbol{\theta}_P\}$ secara acak dalam ruang pencarian $\Theta$
    \While {$g = 1$ \textbf{hingga} $G$}:
      \For{$p = 1$ \textbf{hingga} $P$}:
        \State Dekode kromosom $\boldsymbol{\theta}_p$ menjadi konfigurasi hiperparameter model $M$
        \State Latih \textit{tree-based models} pada $D_{\text{train}}$
        \State Prediksi $D_{\text{test}}$ menggunakan \textit{trained} \textit{tree-based models}
        \State Hitung \textit{Weighted Mean Absolute Percentage Error (WMAPE}
        \State Tetapkan nilai \textit{fitness}: $f(\boldsymbol{\theta}_p) = \dfrac{1}{\mathrm{WMAPE}(\boldsymbol{\theta}_p)}$
      \EndFor
      \State Pilih $n_{\text{parents}}$ induk terbaik melalui metode seleksi turnamen berukuran $K_{\text{tourn}}$
      \State Terapkan \textit{uniform crossover} pada pasangan induk terpilih dengan probabilitas $Pc$
      \State Terapkan mutasi acak (\textit{random mutation}) dengan probabilitas $Pm$
      \State Pertahankan $e = 2$ individu terbaik dari generasi $g$
      \State Bentuk populasi $\mathcal{P}_g$ dari gabungan individu terbaik dan individu baru
    \EndWhile
    \State \Return Nilai optimal hiperparameter pada \textit{tree-based models}.
    \end{algorithmic}
    \end{algorithm}

    \subsection{Representasi Kromosom}

    Representasi kromosom adalah proses \textit{encoding} gen pada suatu kromosom. Penelitian ini menggunakan jenis pengodean gen \textit{real-coded encoding} yang memiliki kelebihan dibandingkan dengan pengodean biner apabila digunakan untuk optimasi fungsi, di mana metode ini dapat menjangkau beberapa titik solusi apabila range solusi berada dalam daerah kontinu \cite{mahmudy2013algoritma}. Pengodean \textit{real coded} atau \textit{real code encoding} merupakan teknik penkodean dari struktur kromosom berdasarkan nilai parameter asli dari \textit{tree based models} yang dipresentasikan sebagai vektor. Teknik ini memliliki kelebihan daripada teknik \textit{binary encoding} yang memiliki keterbatasan jumlah bit ketika mempresentasikan setiap nilai khususnya dalam domain nilai kontinu \cite{El-Hassani2024}.
 
    \begin{equation}
        \mathbf{C} = \left( c_1,\, c_2,\, \ldots,\, c_n \right) \in \mathbb{R}^n
    \end{equation}
     
    \noindent di mana setiap gen $x_i$ merupakan bilangan real kontinu yang dibatasi oleh:
     
    \begin{equation}
        x_i \in \left[ l_i,\; u_i \right], \quad l_i,\, u_i \in \mathbb{R},\quad l_i < u_i,
        \quad i = 1, 2, \ldots, n
    \end{equation}
     
    \noindent dengan $l_i$ adalah batas bawah (\textit{lower bound}) dan $u_i$ adalah
    batas atas (\textit{upper bound}) gen ke-$i$.
     
    % ─────────────────────────────────────────────────
    \subsection{Inisialisasi Populasi}
    
    Inisialisasi populasi merupakan tahapan awal untuk membentuk populasi yang direpresentasikan oleh gen dan kromosom. Satu kromosom $c$ mempresentasikan himpunan parameter pada \textit{tree based Models}, di mana setiap gen merupakan parameter yang  ada di tiap kromosom. Pembentukan populasi awal berpengaruh kepada kinerja evaluasi GA untuk melakukan eksplorasi global sehingga pada penelitian ini populasi awal didefinisikan sebagai sejumlah kumpulan parameter default dari tiap \textit{TreeBased Models}. Secara matematis, populasi algoritma genetika dapat dipresentasikan sebagai metriks berikut.

    \begin{equation}
    P =
    \begin{bmatrix}
    c_{1}^{(1)} & c_{2}^{(1)} & \cdots & c_{m}^{(1)} \\
    c_{1}^{(2)} & c_{2}^{(2)} & \cdots & c_{m}^{(2)} \\
    c_{1}^{(3)} & c_{2}^{(3)} & \cdots & c_{m}^{(3)} \\
    \vdots      & \vdots      & \ddots & \vdots      \\
    c_{1}^{(N)} & c_{2}^{(N)} & \cdots & c_{m}^{(N)}
    \end{bmatrix},
    \label{eq:population_matrix}
    \end{equation}

    dengan $N$ menyatakan jumlah kromosom dalam populasi dan $m$ menyatakan jumlah gen dalam satu kromosom.

    
    \subsection{\textit{Selection}}

    \textit{Selection} adalah proses untuk mendapatkan kromosom-kromosom terbaik untuk generasi berikutnya. Pemilihan kromosom pada tahap ini didasarkan pada nilai \textit{fitness} \cite{Fanggidae2015}. 
    
    Dalam penelitian ini, proses \textit{selection} menggunakan dua 
    mekanisme yang bekerja secara bersamaan. Pertama, metode 
    \textit{Steady-State Selection} (SSS) digunakan untuk memilih 
    sejumlah $N_{\text{parents}}$ kromosom terbaik dari populasi 
    berukuran $N_{\text{pop}}$ berdasarkan peringkat nilai 
    \textit{fitness} untuk menjadi \textit{parent} pada operasi 
    \textit{crossover} dan mutasi \cite{mahmudy2013algoritma}. 
    Pada SSS, kromosom dengan nilai \textit{fitness} tertinggi 
    diprioritaskan, sehingga hanya kromosom berkualitas tinggi 
    yang menghasilkan \textit{offspring} baru.
    
    Kedua, mekanisme \textit{elitism} diterapkan dengan 
    mempertahankan sebanyak $k$ kromosom terbaik  dari setiap generasi untuk langsung diwariskan ke generasi berikutnya tanpa melalui 
    proses \textit{crossover} maupun mutasi \cite{mahmudy2013algoritma}. 
    Kombinasi SSS dan \textit{elitism} ini menjamin  kualitas solusi terbaik yang telah ditemukan tidak mengalami degradasi 
    selama proses evolusi berlangsung, sekaligus tetap memungkinkan 
    eksplorasi ruang pencarian melalui \textit{offspring} yang 
    dihasilkan dari $N$ \textit{parent} terpilih.
    
    \subsection{Fungsi \textit{Fitness}}
    
    Fungsi \textit{fitness} $f(\mathbf{x})$ diformulasikan sebagai negasi
    dari \textit{Weighted Mean Absolute Percentage Error} (WMAPE). Pemilihan WMAPE didasarkan pada karakteristik data transaksi telekomunikasi yang memiliki \textit{sparsity} atau nilai nol pada variabel target, di mana apabila dibandingkan dengan penggunaan \textit{Mean Absolute Percentage Error} (MAPE) konvensional akan menghasilkan nilai tak terhingga (\textit{undefined}). WMAPE mengatasi masalah tersebut dengan menghitung total kesalahan absolut yang dibobotkan terhadap total nilai aktual, memberikan gambaran kesalahan secara agregat yang lebih representatif terhadap total omzet bisnis. WMAPE dapat dituliskan sebagai berikut.

    \begin{equation}
        \text{WMAPE}(\mathbf{x}) =
            \frac{\displaystyle\sum_{t=1}^{T} 
                  \left|y_t - \hat{y}_t(\mathbf{x})\right|}
                 {\displaystyle\sum_{t=1}^{T} \left|y_t\right|}
    \end{equation}
    
    \noindent di mana $y_t$ adalah nilai aktual pada periode $t$, $\hat{y}_t(\mathbf{x})$ adalah nilai prediksi yang dihasilkan oleh model dengan konfigurasi hyperparameter $\mathbf{x}$, dan $T$ adalah jumlah total observasi pada data uji.
    
    Karena GA pada dasarnya merupakan algoritma maksimisasi, yaitu selalu mencari kromosom dengan nilai \textit{fitness} tertinggi,
    sementara WMAPE adalah metrik minimisasi di mana nilai yang lebih kecil menunjukkan prediksi yang lebih baik, maka diperlukan transformasi agar keduanya searah \cite{Gad2021}. Oleh karena itu, fungsi
    \textit{fitness} diformulasikan sebagai negasi dari WMAPE.

    \begin{equation} \label{eq:fitness_neg}
        f(\mathbf{x}) = -\,\text{WMAPE}(\mathbf{x})
    \end{equation}

    \noindent Dengan transformasi ini, kromosom yang menghasilkan WMAPE kecil akan memiliki nilai $f(\mathbf{x})$ mendekati nol (lebih tinggi), sedangkan WMAPE besar menghasilkan $f(\mathbf{x})$ sangat negatif (lebih rendah), sebagaimana diilustrasikan pada Tabel~\ref{tab:fitness-wmape}.

    \begin{table}[H]
    \centering
    \caption{Hubungan nilai WMAPE dan fungsi \textit{fitness}}
    \label{tab:fitness-wmape}
    \begin{tabular}{ccc}
    \hline
    \textbf{Kualitas prediksi} & \textbf{WMAPE} & \textbf{$f(\mathbf{x})$} \\
    \hline
    Sangat baik  & $\to 0$     & $\to 0$ (tertinggi) \\
    Baik         & kecil       & negatif kecil \\
    Buruk        & besar       & sangat negatif \\
    \hline
    \end{tabular}
    \end{table}

    \subsection{\textit{Crossover}}
    
    \textit{Crossoever} atau kawin silang adalah metode yang mengawinkan 2 kromosom \textit{parent} dengan tujuan untuk mendapatkan anak yang lebih baik. Pada tahap ini dihasilkan \textit{offspring} sebanyak 
    $\lfloor N_{\text{pop}} \times P_c \rfloor$ atau dapat dituliskan 
    sebagai perkalian ukuran populasi dengan \textit{crossover rate}. Nilai \textit{crossover rate} harus ditentukan sebelum melakukan tahap perhitungan \textit{crossover}\cite{mahmudy2013algoritma}. Dalam penelitian Tugas Akhir ini menggunakan teknik \textit{uniform crossover} atau pada implementasi di library PyGAD menggunakan \textit{scattered}. Teknik ini dipilih karena cocok dengan tipe data dari parameter yang ingin dioptimasi yang merupakan campuran antara integer dan float. Teknik ini merupakan generalisasi dari metode \textit{one-point} dan \textit{two-point crossover}, di mana posisi \textit{crossover} ditentukan secara acak untuk setiap gen secara independen \cite{Deep2007}.
    Proses \textit{uniform crossover} dimulai dengan membangkitkan vektor biner acak $\mathbf{m} = (m_1, m_2, \ldots, m_n)$, di mana setiap $m_i \in \{0, 1\}$ dibangkitkan secara independen dengan peluang $P(m_i = 1) = 0{,}5$. Vektor ini berfungsi sebagai \textit{mask} yang menentukan dari \textit{parent} mana setiap gen \textit{offspring} diwariskan. Dua \textit{offspring} dari dua \textit{parent} $p^{(1)}$ dan $p^{(2)}$ dihasilkan sebagai berikut.
        \begin{equation}
            x_i^{\text{off}(1)} = m_i \cdot p_i^{(1)} + (1 - m_i) \cdot p_i^{(2)}
        \end{equation}
        \begin{equation}
            x_i^{\text{off}(2)} = (1 - m_i) \cdot p_i^{(1)} + m_i \cdot p_i^{(2)}
        \end{equation}
    \noindent Dari rumus di atas, ketika $m_i = 1$ maka \textit{offspring} pertama mewarisi gen dari $p^{(1)}$ dan \textit{offspring} kedua mewarisi dari $p^{(2)}$, sedangkan ketika $m_i = 0$ sebaliknya. Salah satu sifat penting \textit{uniform crossover} adalah \textit{complementary preservation}, yaitu gen yang tidak diwariskan oleh \textit{offspring} pertama akan diwariskan oleh \textit{offspring} kedua, sehingga tidak ada informasi genetik yang hilang.
        \begin{equation}
            x_i^{\text{off}(1)} + x_i^{\text{off}(2)} = p_i^{(1)} + p_i^{(2)}
        \end{equation}
    \noindent Sifat ini menjamin bahwa kombinasi gen dari kedua \textit{parent} terjaga secara utuh. Penggunaan \textit{mask} acak independen membuat \textit{uniform crossover} memiliki keragaman posisi \textit{crossover} yang lebih tinggi dibandingkan \textit{one-point} atau \textit{two-point crossover}, sehingga mendukung eksplorasi ruang pencarian secara lebih merata \cite{Deep2007, Mitchell1998}.
        
    \subsection{\textit{Mutation}}

    \textit{Genetic Algorithm} (GA) menggunakan mutasi untuk menciptakan variasi genetik pada beberapa gen dalam suatu kromosom untuk menjaga keberagaman populasi. Tingkat mutasi diatur rendah karena tingkat mutasi yang tinggi dapat mengubah GA menjadi pencarian acak yang primitif \cite{Fanggidae2015}. Tujuan proses ini adalah menghindari \textit{premature convergence}, yaitu kondisi di mana kromosom-kromosom mengalami konvergensi saat nilai \textit{fitness} belum mencapai optimal. 
    Mutasi ditentukan oleh \textit{mutation rate} ($m_r$) yang menjadi penentu banyaknya \textit{offspring} yang dihasilkan dari populasi, yaitu sejumlah $m_r \times$ ukuran populasi \cite{mahmudy2013algoritma}.
    
    Dalam penelitian Tugas Akhir ini, teknik mutasi yang digunakan adalah \textit{Random Mutation} (RM). RM menghasilkan nilai gen baru secara acak seragam dalam rentang batas yang diizinkan, sehingga setiap gen yang termutasi memiliki peluang yang sama untuk mengambil nilai mana pun di dalam domain variabel tersebut \cite{Herrera1998}. Proses RM dimulai dengan membangkitkan bilangan acak seragam $u \sim \mathcal{U}(0,1)$, kemudian gen hasil mutasi diperoleh melalui persamaan berikut.
    
    \begin{equation}
        x_i' = l_i + u \cdot (u_i - l_i)
    \end{equation}
    
    \noindent di mana $l_i$ dan $u_i$ masing-masing adalah batas bawah dan batas atas gen ke-$i$. Nilai $x_i'$ dijamin berada dalam rentang $[l_i, u_i]$ karena $u \in [0, 1]$ sehingga:
    
    \begin{equation}
        l_i \leq x_i' \leq u_i
    \end{equation}
    
    \noindent Keunggulan RM adalah kemampuannya melakukan eksplorasi global secara merata di seluruh ruang pencarian, karena distribusi nilai mutatif bersifat seragam (\textit{uniform}) tanpa bias terhadap nilai gen sebelumnya. Hal ini menjadikan RM efektif dalam menghindari \textit{premature convergence} dan menjaga keanekaragaman populasi, terutama pada fase awal evolusi \cite{Michalewicz1996, Herrera1998}. RM dipilih karena kesederhanaan implementasinya pada \textit{real-coded encoding} dan konsistensinya dengan operator \textit{scattered crossover}, di mana keduanya mengandalkan pembangkitan bilangan acak seragam untuk menghasilkan variasi genetik.


    % \section{\textit{Time Series Cross Validation}}

    



\endinput
            \chapter{Metode Penelitian}\label{ch:method}

    \section{Tahapan Penelitian}\label{sec:stages}
    Bab ini membahas metode yang digunakan dalam pelaksanaan penelitian Tugas Akhir. Penelitian ini mengimplementasikan algoritma genetika untuk optimasi parameter pada model yang memprediksi total \textit{revenue} (pendapatan) dari salah satu perusahaan SMS \textit{traffic management} di Indonesia. Tahapan penelitian pada Tugas Akhir ini ditunjukkan oleh Gambar 3.1 berikut.

    \begin{figure}[H]
        \centering
        \includegraphics[width=0.9\textwidth]{tugas-akhir/assets/images/Diagram Alir Penelitian.png}
        \caption{Diagram Alir Penelitian}
        \label{fig:ga-flow}
    \end{figure}
    
    Tahapan penelitian dimulai dari pengumpulan data \textit{Traffic SMS Internasional} dari Perusahaan Telekomunikasi (X). Data tersebut kemudian dibersihkan (\textit{data cleaning}) dengan beberapa proses yaitu normalisasi data, penyamaran data untuk privasi, pengurutan data berdasarkan bulan dan tahun, pembersihan nilai tidak valid, serta penghapusan kolom data atau fitur yang tidak diperlukan. Kemudian dilakukan proses \textit{Exploratory Data Analysis} untuk mendapatkan \textit{insight} dari data. Selanjutnya dilakukan pra-pemrosesan data dengan menggabungkan data berdasarkan bulan data  dan pembuatan \textit{feature engineering}. Kemudian data dimasukkan dalam model prediksi untuk memprediksi kolom target yaitu \textit{TotalRevenue} dengan dua tahapan (\textit{two stage}) yaitu klasifikasi dan regresi. Klasifikasi bertujuan untuk memprediksi pendapatan yang nol dan tidak nol. Regresi bertujuan untuk memprediksi  pendapatan tidak nol di bulan selanjutnya dari hasil prediksi model klasifikasi. Model prediksi pada regresi dioptimasi dengan \textit{hypeparamater tuning} menggunakan Algoritma Genetika. Parameter terbaik akan digunakan dalam model untuk memprediksi \textit{TotalRevenue}. Hasil prediksi akan dievaluasi menggunakan metriks evaluasi yaitu WMAPE, MAE dan $R^2$ dan \textit{error} secara statistik dan visual.
    
    \section{Pengumpulan Data}\label{sec:data collection}
    Data yang digunakan dalam penelitian ini adalah data transaksi SMS A2P (\textit{Application to Person}) dari perusahaan X, salah satu perusahaan yang berfokus pada \textit{traffic management} pada lalu lintas SMS Internasional,  pada tahun 2023 sampai 2024. Data yang digunakan terdiri dari 94.335 data dengan 19 kolom, kolom fitur sebanyak 18 kolom  yaitu Year, Month, Client, ClientConnection, Supplier, SupplierConnection, SupplierConnectionLink, Route, Country, Network, OAProfiledCategory,  CostCurrency, PriceCurrency, ProfitCurrency, BaseCurrency, TotalCost, TotalSentMessages,TotalDeliveredMessages, TotalUndeliveredMessages dan kolom target yaitu TotalRevenue.

    \section{Pembersihan Data}\label{sec:preprocessing}
    Pembersihan data merupakan tahapan penting sebelum data digunakan dalam proses pelatihan model maupun optimasi hiperparameter. Tahapan ini bertujuan untuk memastikan bahwa data yang digunakan bersih, terstruktur, dan sesuai dengan kebutuhan algoritma. Pada penelitian ini, proses pembersihan data meliputi beberapa langkah, yaitu normalisasi data, pengurutan data berdasarkan waktu, penghapusan kolom yang tidak diperlukan, dan pembersihan nilai tidak valid. Berikut merupakan penjelasan rinci setiap tahapan.

    \subsection{Normalisasi Data}\label{}

    Dataset yang digunakan memiliki ketidaknormalan data dan duplikasi data akibat salah ketik. Oleh karena itu, untuk mengatasi hal ini, langkah awal pembersihan adalah menormalisasi data akibat salah ketik tersebut menjadi data yang benar sehingga data yang salah ketik tersebut tidak dianggap data yang berbeda oleh model. Tujuannya adalah untuk memastikan data yang diproses sudah benar dan memang tidak ada kesalahan ketik ketika nantinya dimasukkan ke model. Data yang salah ketik ini terdapat pada kolom kategorikal.

    \subsection{Penyamaran Data untuk Privasi}\label{}
    Data yang digunakan memiliki informasi nama perusahaan yang sensitif. Menurut perjanjian dengan perusahaan X, peneliti harus menyamarkan data yang terkait dengan nama perusahaan dan informasinya sehingga data tersebut disamarkan dengan menggunakan metode penyamaran deterministik. Metode yang digunakan untuk penyamaran adalah pemetaan deterministik (\textit{deterministic mapping}) sehingga data yang dipetakan akan konsisten di tiap kolom. Kolom-kolom yang diidentifikasi sensitif dan disamarkan yaitu \texttt{Client}, \texttt{ClientConnection}, \texttt{Supplier}, \texttt{SupplierConnection}, \texttt{SupplierConnectionLink}, \texttt{Route}, \texttt{Network}, dan \texttt{OAProfiledCategory}.
    
    \subsection{Pengurutan Data Berdasarkan Bulan dan Tahun}\label{}
    
    Data \textit{traffic} SMS internasional yang diperoleh dari Perusahaan X memiliki informasi waktu berupa bulan dan tahun. Oleh karena itu, langkah selanjutnya adalah mengurutkan data berdasarkan urutan kronologis. Tujuannya untuk memastikan analisis statistik dan eksplorasi data dilakukan secara berurutan, memudahkan pemetaan pola per bulan, dan mendukung proses pembagian data yang lebih representatif. Pengurutan dilakukan menggunakan kombinasi \textit{sort\_values} pada kolom waktu (misalnya \textit{Month} dan \textit{Year}).
    
    \subsection{Penghapusan Kolom yang Tidak Diperlukan}\label{} 
    
    Dataset awal memiliki sejumlah kolom yang tidak memiliki pengaruh signifikan terhadap model prediksi atau analisis. Pada tahap ini dilakukan penghapusan terhadap kolom kolom identitas yang tidak relevan, kolom yang tidak berhubungan dengan target (\textit{TotalRevenue}), serta kolom yang dapat menyebabkan data \textit{leak}. Penghapusan dilakukan untuk menyederhanakan dataset dan mengurangi kompleksitas model, sehingga proses \textit{training} menjadi lebih efisien.
    
    \subsection{Pembersihan Nilai Tidak Valid}\label{} 
    Tahap selanjutnya adalah pembersihan nilai yang tidak valid, seperti nilai nol yang tidak semestinya (misalnya \textit{TotalRevenue} tetapi bernilai 0), nilai negatif, entri kosong (\textit{missing values}), dan data inkonsisten akibat kesalahan input. Pembersihan data ini bertujuan untuk mengurangi \textit{noise} dan memastikan informasi yang ada di dataset relevan.
    
    \section{Exploratory Data Analysis \method}\label{sec:method construction}
    Selanjutnya dilakukan proses \textit{exploratory data analysis} untuk mengetahui pola dalam data seperti distribusi data, frekuensi tiap data dalam satu fitur, dan hubungan korelasi antara fitur numerik maupun kategorikal. Analisis dalam proses ini dilakukan dengan membuat berbagai macam \textit{chart} atau grafik visual dengan memanfaatkan library \textit{matplotlib} atau \textit{seaborn}. Analisis visual ini bermanfaat untuk membuat sebuah fitur baru seperti yang ada di proses sebelumnya dan proses pemilihan fitur (\textit{feature selection}). Proses \textit{feature selection} ini dapat dilakukan berdasarkan fitur-fitur yang memiliki korelasi tinggi dengan fitur target berdasarkan analisis visual yang dilakukan serta analisis berdasarkan pengetahuan domain tertentu (\textit{domain knowledge}). 

    
    \section{Pra-pemrosesan Data}\label{} 

    Setelah melakukan proses EDA untuk mengeksplorasi insight bisnis serta pola dari data dilakukan pra-pemrosesan data. Tahapan ini bertujuan  untuk memastikan data yang dimasukkan ke model dapat diprediksi dengan lebih mudah serta dengan pola yang didapatkan sebelumnya dapat ditentukan pra-pemrosesan data yang lebih sesuai. Pada penelitian Tugas Akhir ini, proses pra-pemrosesan data terdiri dari penggabungan data berdasarkan bulan, pembuatan dan \textit{feature engineering}. Berikut merupakan penjelasan rinci di setiap tahapan.

    
    \subsection{Penggabungan Data Berdasarkan Bulan}\label{}  

    Data digabung dengan menjumlahkan metriks utama yaitu \texttt{TotalRevenue}, \texttt{TotalCost}, \texttt{TotalProfit}, \texttt{TotalSentMessages}, \texttt{TotalDeliveredMessages}, dan \texttt{TotalUndeliveredMessages} dengan metode \textit{sum}. Kemudian penggabungan tersebut berdasarkan \texttt{Series\_ID}, \texttt{Year}, \texttt{Month} sehingga bisa dianggap sebagai penggabungan data perbulan. \texttt{Series\_ID} sendiri  merupakan kombinasi antara \texttt{Client}, \texttt{Route}, \texttt{SupplierConnection}, dan \texttt{Network}. Penggabungan juga mempertimbangkan fitur lain yaitu dengan penghitungan data unik dari kolom \texttt{OAProfiledCategory} dengan metode \textit{nunique} yang mana nantinya dihasilkan fitur baru yaitu \texttt{unique\_OA\_count}, banyaknya \textit{flag} data anomali dari fitur \texttt{is\_anomaly\_rev\_only} dan \texttt{is\_anomaly\_cost\_only} menghasilkan fitur \texttt{anomaly\_count} serta rasio data anomali dengan total baris per grup menghasilkan fitur \texttt{anomaly\_revenue\_pct}, pendapatan maksimum dan minimum per kategori OA (\texttt{OAProfiledCategory}) dalam setiap grup dihitung yang menghasilkan fitur \texttt{max\_OA\_revenue} dan \texttt{min\_OA\_revenue} yang mengukur heterogenitas sub-transaksi.

    \subsection{\textit{Feature Engineering}}\label{}  
    \textit{Feature engineering} adalah proses untuk mengekstraksi fitur baru ataupun memodifikasi fitur yang sudah ada agar model lebih mampu mempelajari pola dalam data. Beberapa langkah \textit{feature engineering} yang dilakukan antara lain pembentukan fitur \texttt{Series\_ID}, pembuatan fitur untuk berdasarkan pola yang ada di data dengan analisis \textit{domain knowledge}, pembuatan fitur rasio atau perbandingan \textit{cost} (\texttt{rate\_Cost}), penambahan fitur baru untuk menandai anomali atau \textit{outliers} pada data, pembuatan fitur berdasarkan pola dari data sebelumnya, pembuatan fitur temporal yang mempresentasikan dinamik perubahan bisnis secara bulanan. Proses \textit{feature engineering} ini ada yang dilakukan sebelum proses penggabungan data berdasarkan bulan ada juga yang setelah proses penggabungan data berdasarkan bulan. Namun kebanyakan proses yang dilakukan setelah proses penggabungan data berdasarkan bulan karena data setelah proses ini yang akan dimasukkan kedalam proses pelatihan model. Langkah ini meningkatkan representasi informasi pada dataset sehingga dapat memperbaiki performa model.

    \section{Pelatihan Model \method}\label{sec:method construction}
    Untuk menyelesaikan prediksi \textit{revenue} yang mana memiliki karakteristik data yang agak kompleks. Data dengan pengelompokan berdasarkan \textit{Series\_ID} belum tentu ada di tiap bulan sehingga perlu untuk memprediksi ada atau tidaknya transaksi di tiap bulan kedepan dari masing-masing \texttt{Series\_ID}. Setelah penggabungan data beradasarkan bulan, data dengan nilai transaksi yang ada di tiap bulan memiliki proporsi 47.49\% dari seluruh data atau 994 dari 2093 total data. Maka dari itu perlu proses klasifikasi untuk membedakan ada atau tidaknya transaksi di tiap bulan kedepan. Hasil prediksi dari proses klasifikasi dapat disimbolkan sebagai nilai $P$ dengan $P = 0$ menunjukkan tidak adanya transaksi dan $P=1$ menunjukkan adanya transaksi. Apabila data tersebut ada transaksi baru dijalankan proses regresi untuk memprediksi \texttt{TotalRevenue}. Maka dari itu, hasil akhir dapat dinyatakan sebagai hasil perkalian antara proses klasifikasi dan proses regresi sehingga apabila $P = 0$ maka hasil akhir juga bernilai 0 sedangkan jika $P = 1$ maka hasil akhir dapat dinyatakan sebagai hasil output dari proses regresi karena $P = 1$ $.$ output regresi $=$ output regresi. 
    Untuk pembagian data dari proses klasifikasi diambil 80 \% data untuk data train dan 20 \% data untuk data test dengan menggunakan metode \textit{time based split} yaitu pembagian data berdasarkan urutan dari waktu \textit{month} dan \textit{year}. Model yang digunakan adalah model XGBoost dan menggunakan parameter default yaitu \texttt{n\_estimators} = 1000, \texttt{learning\_rate} = 0.3, \texttt{num\_depth} = 6,  \texttt{subsamples} = 1, \texttt{colsample\_bytree} = 1, \textit{random\_state} = 42, \texttt{early\_stopping\_rounds} = 50, dan \texttt{verbose} = 0. Evaluasi model menggunakan ROC AUC karena data \textit{imbalance} untuk transaksi yang aktif dan tidak aktif. Kemudian untuk proses regresi menggunakan pembagian data yang sama dengan proses klasifikasi hanya saja model yang digunakan lebih beragam yaitu menggunakan model XGBoost dan Random Forest. Dengan parameter default masing-masing model untuk XGBoost menggunakan \texttt{random\_state} = 42 dan \textit{verbose} = 0 sedangkan Random Forest mengguakan \texttt{random\_state} = 42 dan \texttt{n\_jobs} = -1. Dalam proses ini proses optimasi hiperparameter \textit{tree-based models} juga dilakukan yang mana hasilnya nanti akan dibandingkan dengan \textit{baseline} yaitu pelatihan model dari masing-masing model yang ada di proses regresi dengan parameter default. Proses ini akan dijelaskan lebih lanjut di subbab berikutnya. 

    \section{Optimasi Hiperparameter Menggunakan Algoritma Genetika\method}\label{sec:method construction}\
    Untuk mendapatkan konfigurasi parameter terbaik, penelitian ini menggunakan Algoritma Genetika (\textit{Genetic Algorithm}), sebuah metode pencarian heuristik yang terinspirasi dari proses seleksi alam. Pendekatan ini dipilih karena kemampuannya yang superior dalam mengeksplorasi ruang pencarian parameter yang luas dibandingkan metode konvensional. Ruang pencarian parameter dapat dilihat pada Tabel 3.1. Ruang pencarian ini merujuk pada penelitian Mansoori (2023).
    
    \renewcommand{\arraystretch}{1.2}
    \begin{longtable}{l l l l >{\RaggedRight\arraybackslash}p{3cm}}
    \caption{Rentang Hiperparameter Model} \\
    \toprule
    \textbf{Model} & \textbf{Hiperparameter} & \textbf{Rentang} & \textbf{Tipe} & \textbf{Definisi} \\
    \midrule
    \endfirsthead
    \toprule
    \textbf{Model} & \textbf{Hiperparameter} & \textbf{Rentang} & \textbf{Tipe} & \textbf{Definisi} \\
    \midrule
    \endhead
    
    Random Forest & $n\_estimators$ & [50, 1000] & Integer & Jumlah pohon dalam ensemble \\
     & $max\_depth$ & [1, 7] & Integer & Kedalaman maksimum pohon \\
     & $min\_samples\_split$ & [2, 20] & Integer & Minimum sampel pemisahan node \\
     & $min\_samples\_leaf$ & [1, 50] & Integer & Minimum sampel pada daun \\
     & $max\_features$ & [1, 23] & Integer & Jumlah fitur yang dipertimbangkan \\
     & $max\_samples$ & [0.1, 1.0] & Float & Proporsi sampel untuk melatih tiap pohon \\
    \midrule
    XGBoost & $learning\_rate$ & [0.001, 1.0] & Float & Laju pembelajaran model \\
     & $n\_estimators$ & [50, 1000] & Integer & Jumlah iterasi boosting \\
     & $max\_depth$ & [1, 7] & Integer & Kedalaman maksimum pohon \\
     & $subsample$ & [0.1, 1.0] & Float & Proporsi sampel yang diambil \\
     & $colsample\_bytree$ & [0.1, 1.0] & Float & Proporsi fitur yang diambil \\
     & $reg\_lambda$ & [1, 3] & Float & Regularisasi L2 \\
    \bottomrule
    \end{longtable}
    \subsection{Representasi Kromosom}\label{sec:method construction}\
    Pada tahap ini, parameter \textit{tree-based models} direpresentasikan sebagai kromosom. Sesuai dengan penjelasan sebelumnya satu kromosom untuk Random Forest dan XGBoost sama - sama terdiri dari 6 gen. Karena kedua model dioptimalkan secara terpisah, setiap model memiliki representasi kromosom dan ruang pencarian masing-masing. Tiap gen tersebut masing-masing merepresentasikan satu hiperparameter. Setiap \textit{tree-based models} memiliki ruang pencarian yang akan dioptimalkan. Ruang pencarian ini dapat berupa rentang kontinu dan diskrit tergantung dari tipe data dari tiap parameter. Untuk lebih lengkapnya ruang pencarian Random Forest dapat dilihat pada Tabel \ref{tab:search_space_rf}.

    \begin{table}[H]
    \centering
    \caption{Ruang pencarian hiperparameter \textit{XGBoost}}
    \label{tab:search_space_rf}
    \begin{tabular}{llccc}
    \toprule
    \textbf{Gen} & \textbf{Hiperparameter} & \textbf{Tipe} & \textbf{Min} & \textbf{Maks} \\
    \midrule
    0 & \texttt{n\_estimators}       & Integer & 50    & 1.000 \\
    1 & \texttt{max\_depth}          & Integer & 1     & 7     \\
    2 & \texttt{min\_samples\_split} & Integer & 2     & 20    \\
    3 & \texttt{min\_samples\_leaf}  & Integer & 1     & 50    \\
    4 & \texttt{max\_features}       & Integer & 1     & 23 \\
    5 & \texttt{max\_samples}        & Float   & 0{,}1 & 1{,}0 \\
    \bottomrule
    \end{tabular}
    \end{table}
    
    Sama dengan kromosom \textit{Random Forest}, satu kromosom \textit{XGBoost} juga terdiri dari 6 gen. Tabel \ref{tab:search_space_rf} merangkum ruang pencarian \textit{XGBoost}. Proses dekode kromosom ke hiperparameter dilakukan melalui fungsi \texttt{decode\_rf()} dan \texttt{decode\_xgb()} yang masing-masing membulatkan nilai gen ke tipe data yang sesuai misalkan \textit{integer} atau \textit{float} dan memastikan berada dalam rentang yang valid melalui fungsi \texttt{np.clip()}.
    
    \subsection{Inisialisasi Populasi}\label{sec:method construction}\
    
    Pertama-pertama, populasi awal dinisiasi yang berisi sekumpulan kombinasi hiperparameter acak. Tujuan dari proses ini adalah untuk menciptakan solusi awal yang beragam sehingga dapat dieksplorasi oleh algoritma genetika. Ukuran populasi pada penelitian Tugas Akhir ini adalah 20. Dengan ukuran populasi tersebut terdiri dari 20 kromosom yang berbeda. Kemudian untuk \textit{random\_seed} diatur ke nilai 42 agar ditiap eksperimen yang dilakukan berulang dapat memiliki populasi awal yang sama. Tanpa pengaturan \textit{random\_seed} bisa berbeda di tiap eksperimen. Nilai 42 sendiri bukan ketentuan melainkan nilai yang popular digunnakan untuk \textit{data science} dalam mengatur nilai acak.

    \subsection{Evaluasi Fungsi \textit{Fitness}}\label{sec:method construction}\
    Pada setiap generasi, seluruh individu dalam populasi dievaluasi menggunakan fungsi objektif (\textit{objective function}) atau \textit{fitness function}. Pada penelitian ini, tujuan utama algoritma genetika adalah untuk meminimalkan kesalahan prediksi model. Oleh karena itu, \textit{objective function} dirancang menggunakan negasi dari metrik \textit{Weighted Mean Absolute Percentage Error} (WMAPE), yang mengukur kesalahan relatif antara nilai aktual dan nilai prediksi dengan mempertimbangkan bobot nilai aktual. Secara matematis, \textit{fitness function} dirumuskan pada persamaan 2.7. Proses negasi dilakukan karena GA pada library PyGad pada dasarnya berusaha memaksimalkan kromosom dengan nilai \textit{fitness} yang lebih tinggi sementara tujuan dari penelitian ini adalah untuk meminimalkan WMAPE. Pemilihan WMAPE didasarkan pada karakteristik data transaksi telekomunikasi yang memiliki \textit{sparsity} atau nilai nol pada variabel target, di mana apabila dibandingkan dengan penggunaan \textit{Mean Absolute Percentage Error} (MAPE) konvensional akan menghasilkan nilai tak terhingga (\textit{undefined}). WMAPE mengatasi masalah tersebut dengan menghitung total kesalahan absolut yang dibobotkan terhadap total nilai aktual, memberikan gambaran kesalahan secara agregat yang lebih representatif terhadap total omzet bisnis. 

    \subsection{Seleksi}\label{sec:method construction}\
    Pada proses seleksi, metode yang digunakan adalah metode \textit{tournament} yaitu individu dipilih secara acak dari populasi, lalu yang terbaik dari grup itu yang lolos menjadi parent. Metode ini dipilih karena metode ini cocok dengan range parameter yang lebar dan juga pendek di mana \textit{tree-based models} memiliki keragaaman \textit{range} parameter. Untuk parameter $k$ yang mempresentasikan jumlah anggota yang diambil acak di tiap grup diatur nilainya dengan $k = 3$. Nilai $k$ ini sudah sesuai karena tidak melebihi dari jumlah \textit{parents}(\textit{num\_parents\_mating}) yang akan diambil untuk dilakukan \textit{crossover}. Pada tahap ini jumlah \textit{parents} yang dipilih ditiap generasi yang dipresentasikan dalam PyGAD sebagai \textit{num\_parents\_matings} berjumlah 5.
    
    \subsection{\textit{Crossover}}\label{sec:method construction}\
    
    Selanjutnya reproduksi dilakukan tahap \textit{crossover} (kawin silang) untuk menggabungkan parameter secara acak dengan probabilitas kecil untuk menjaga keragaman populasi dan mencegah model terjebak pada solusi optimum lokal (\textit{local optima}). Dalam penelitian Tugas Akhir ini metode \textit{crossover} yang digunakan adalah \textit{uniform crossover}. Dengan pembentukan hasil dari kromosom anak (\textit{offsprings}) yang dapat dilihat pada persamaan 2.8 dan 2.9. Untuk probabilitas \textit{crossover} (\textit{crossover rate} pada penelitian ini ditentukan sebesar 1 di mana di tiap generasi pasti melakukan \textit{crossover}. 

    \subsection{\textit{Mutasi}}\label{sec:method construction}\

    Pada tahap mutasi, dilakukan proses perubahan pada gen di tiap kromosom untuk menjaga keragaman solusi. Probabilitas mutasi ($Pm$) yang digunakan sebesar 0.2 di mana 20\% dari jumlah gen akan dimutasi di tiap individu. Setelah jumlah gen yang dimutasi ditentukan, lokasi gen yang mengalami mutasi dipilih secara acak. Proses mtuasi dilakukan dengan menggunakan metode \textit{Random Mutation} dengan persamaan yang dapat dilihat pada persamaan 2.11. Setelah proses mutasi selesai, kromosom yang baru akan dievaluasi kembali fungsi \textit{fitness} nya dan dilakukan proses iterasi untuk pembentukan generasi baru begitu seterusnya sampai dengan jumlah generasi yang telah ditentukan. Kriteria penghentian ditentukan oleh jumlah generasi maksimum (50 generasi) dengan keep\_elitism = 2 yang memastikan 2 individu terbaik dari setiap generasi selalu dipertahankan tanpa modifikasi.
    
    \section{Evaluasi Hasil Optimasi Hiperparameter}

    Hasil optimasi hiperparameter kemudian dimasukkan ke dalam \textit{tree-based models} dan dievaluasi model tersebut untuk mengetahui performansinya terhadap data pendapatan layanan SMS internasional di perusahaan X. Dalam penelitian ini, sama dengan \textit{fitness function} metrik evaluasi utama yang digunakan adalah \textit{Weighted Mean Absolute Percentage Error} (WMAPE). 
    
    \begin{equation} \label{eq:r2}
    R^2 = 1 - \frac{\sum_{i=1}^{n} (y_{i} - \hat{y}_{i})^2}{\sum_{i=1}^{n} (y_{i} - \bar{y})^2}
    \end{equation}
    
    \begin{equation} \label{eq:mae}
    MAE = \frac{1}{n} \sum_{i=1}^{n} |y_{i} - \hat{y}_{i}|
    \end{equation}
    
    Selain WMAPE, penelitian ini juga menerapkan dua metrik pendukung lainnya, yaitu \textit{Coefficient of Determination} ($R^2$) dan \textit{Mean Absolute Error} (MAE). Koefisien determinasi ($R^2$) digunakan untuk mengukur proporsi varians dari variabel dependen yang dapat dijelaskan oleh model, di mana nilai yang mendekati 1 mengindikasikan kemampuan model yang sangat baik dalam menangkap pola data. Sementara itu, MAE digunakan untuk menghitung rata-rata selisih absolut antara nilai prediksi dan nilai aktual dalam satuan mata uang asli, yang memberikan interpretasi langsung mengenai besaran nominal kesalahan prediksi rata-rata per transaksi. Rumus untuk $R^2$ dan MAE ditunjukkan masing-masing pada Persamaan \ref{eq:r2} dan \ref{eq:mae}, di mana $y_i$ adalah nilai aktual, $\hat{y}_i$ adalah nilai prediksi, dan $\bar{y}$ adalah rata-rata nilai aktual.
    
    
    \chapter{Hasil dan Pembahasan}\label{ch:result}

Bab ini menyajikan hasil dan pembahasan pelaksanaan penelitian, yang dimulai dari pembersihan data, \textit{exploratory data analysis}, pra-pra-pemrosesan data, pelatihan model, hingga optimasi hiperparameter menggunakan Algoritma Genetika. Seluruh hasil eksperimen dievaluasi menggunakan metrik WMAPE, MAE, dan $R^2$, serta dibandingkan dengan model tanpa optimasi dan menggunakan optimasi.

\section{Data Transaksi Pendapatan SMS Internasional}\label{sec:data}

    Data yang digunakan dalam penelitian ini merupakan data dari Perusahaan X, salah satu perusahaan telekomunikasi di Indonesia yang bergerak di bidang \textit{traffic management} SMS internasional. Nama perusahaan disamarkan sesuai perjanjian kerahasiaan data. Seluruh nilai yang mengandung identitas perusahaan tertentu juga telah disamarkan. Data dalam penelitian ini memiliki 19 kolom dan memiliki berbagai tipe atribut dengan keterangan di setiap kolomnya. Deskripsi setiap atribut ditunjukkan pada Tabel 4.1.

    \begin{longtable}{lll}
    \caption{Deskripsi Data Pendapatan SMS Internasional} \\
    \toprule
    \textbf{Fitur} & \textbf{Tipe Data} & \textbf{Keterangan} \\
    \midrule
    \endfirsthead
    
    \toprule
    \textbf{Fitur} & \textbf{Tipe Data} & \textbf{Keterangan} \\
    \midrule
    \endhead
    
    Year & Integer & Tahun transaksi \\
    Month & Integer & Bulan transaksi \\
    Client & String & Nama klien \\
    Client Connection & String & Koneksi klien \\
    Supplier & String & Nama Supplier \\
    Supplier Connection & String & Koneksi supplier \\
    Supplier Connection Link & String & Link koneksi supplier \\
    Route & String & Rute pengiriman \\
    Country & String & Negara \\
    Network & String & Jaringan distribusi \\
    OAProfiled Category & String & Kategori profil OA (\textit{Provider}) \\
    Cost Currency & Integer & Biaya dalam mata uang tertentu \\
    Price Currency & Integer & Harga dalam mata uang tertentu \\
    Profit Currency & Integer & Profit dalam mata uang tertentu \\
    Base Currency & Integer & Mata uang dasar \\
    Total Cost & Integer & Total biaya \\
    Total Sent Messages & Integer & Total pesan terkirim \\
    Total Delivered Messages & Integer & Total pesan berhasil terkirim \\
    \textbf{Total Revenue} & \textbf{Integer} & \textbf{Variabel target penelitian} \\
    \bottomrule
    \end{longtable}
    
\section{Pembersihan Data (\textit{Data Cleaning})}
\label{sec:prapra-pemrosesan_data}

Tahap pembersihan data dilakukan untuk memastikan data trafik SMS internasional periode 2023--2024 berada pada kondisi yang layak analisis, bersih, dan siap untuk proses pemodelan. Tahap pembersihan data terdiri dari normalisasi data, penyamaran data sensitif, pengurutan data berdasarkan waktu, penghapusan kolom yang tidak relevan, dan pembersihan nilai tidak valid.

% \begin{figure}[htbp]
% \centering
% \includegraphics[width=0.92\textwidth]{tugas-akhir/assets/images/prapra-pemrosesan_data.png}
% \caption{Tahapan Prapra-pemrosesan Data}
% \label{fig:alur_prapra-pemrosesan}
% \end{figure}

% ---------------------------------------------------------
\subsection{Normalisasi Data}
\label{subsubsec:normalisasi}

Sebelum dilakukan analisis lanjutan, dilakukan proses normalisasi nilai-nilai kategorik yang terdapat dalam dataset. Dari seluruh kolom yang dianalisis, ditemukan kesalahan pengetikan pada kolom \texttt{Network}. Pada kolom ini ditemukan sebanyak 27 data terdapat nilai "\texttt{Telkomcel}" yang secara semantik merupakan kesalahan ejaan dari "\texttt{Telkomsel}". Tanpa koreksi, inkonsistensi ini berpotensi menyebabkan model memperlakukan kedua nilai tersebut sebagai dua entitas yang berbeda, sehingga merusak representasi distribusi data jaringan (\texttt{Network}) sebenarnya. Oleh karena itu, dilakukan proses normalisasi dengan langkah-langkah berikut:

\begin{enumerate}
    \item Seluruh nilai pada kolom \texttt{Network} dikonversi ke tipe string dan dilakukan penghapusan spasi di awal dan akhir  untuk menghindari kesalahan pencocokan akibat karakter tersembunyi.
    \item Penggantian nilai dilakukan menggunakan ekspresi reguler yang bersifat \textit{case-insensitive}, sehingga semua varian ejaan "\textit{Telkomcel}" maupun "\textit{Telkomsel}" distandarisasi menjadi "\texttt{Telkomsel}".
\end{enumerate}

Hasil normalisasi ditunjukkan pada Tabel \ref{tab:network_normalisasi}. Setelah proses ini, kolom \texttt{Network} hanya memuat satu nilai "\texttt{Telkomsel}" yang sebelumnya tercatat dalam dua bentuk.

\begin{table}[h!]
\centering
\caption{Distribusi nilai kolom \texttt{Network} sebelum dan sesudah normalisasi}
\label{tab:network_normalisasi}
\begin{tabular}{lcc}
\toprule
\textbf{Nilai} & \textbf{Sebelum Normalisasi} & \textbf{Sesudah Normalisasi} \\
\midrule
Telkomsel  & 78.991 & 79.018 \\
Telkomcel  & 27     & 0      \\
\midrule
\textbf{Total} & \textbf{79.018} & \textbf{79.018} \\
\bottomrule
\end{tabular}
\end{table}



% ---------------------------------------------------------
\subsection{Penyamaran Data untuk Privasi}
\label{subsubsec:data_masking}

Dataset \texttt{SMS International.csv} mengandung sejumlah kolom yang memuat informasi identitas bisnis yang bersifat rahasia, seperti nama klien, pemasok, koneksi pemasok, rute trafik, dan jaringan operator. Penggunaan nilai-nilai asli tersebut dalam penelitian melanggar aspek kerahasiaan data yang terdapat dalam perjanjian penggunaan data dengan perusahaan X. Oleh karena itu, dilakukan proses penyamaran data secara deterministik sebelum data digunakan dalam analisis dan visualisasi.

Metode penyamaran yang diterapkan adalah pemetaan deterministik (\textit{deterministic mapping}), yaitu setiap nilai unik pada kolom sensitif dipetakan ke token anonim berformat \texttt{NAMAKOLOM\_XXX}, di mana \texttt{XXX} merupakan indeks tiga digit yang diurutkan secara leksikografis. Dengan pendekatan ini, konsistensi referensial antar data tetap terjaga dan nilai yang sama selalu dipetakan ke token yang sama, sehingga analisis berbasis relasi tetap valid tanpa mengungkapkan data asli. Kolom-kolom yang diidentifikasi sebagai sensitif dan diterapkan proses penyamaran meliputi
\texttt{Client}, \texttt{ClientConnection}, \texttt{Supplier}, \texttt{SupplierConnection}, \texttt{SupplierConnectionLink}, \texttt{Route}, \texttt{Network}, dan \texttt{OAProfiledCategory}. Contoh hasil penyamaran disajikan pada Tabel \ref{tab:contoh_masking}.  Untuk memberikan gambaran data, ~\autoref{ch:lampiran1}\ menyajikan 50 data pertama ketika data sudah disamarkan.

\begin{table}[h]
\centering
\caption{Contoh hasil penyamaran}
\label{tab:contoh_masking}
\begin{tabular}{lll}
\toprule
\textbf{Kolom} & \textbf{Nilai Asli (disamarkan)} & \textbf{Token Anonim} \\
\midrule
Client              & [Identitas klien asli]    & \texttt{CLIENT\_031} \\
ClientConnection    & [Identitas koneksi asli]  & \texttt{CLIENTCONNECTION\_065} \\
Supplier            & [Identitas pemasok asli]  & \texttt{SUPPLIER\_001} \\
SupplierConnection  & [Identitas koneksi asli]  & \texttt{SUPPLIERCONNECTION\_001} \\
Route               & [Identitas rute asli]     & \texttt{ROUTE\_041} \\
Network             & [Identitas jaringan asli] & \texttt{NETWORK\_007} \\
\bottomrule
\end{tabular}
\end{table}

\subsection{Pengurutan Data berdasarkan Bulan dan Tahun}

Pada tahap ini, data diurutkan berdasarkan variabel \texttt{Year} dan \texttt{Month} secara \textit{ascending} agar pola temporal terbentuk dengan benar. Langkah ini kritis untuk mencegah \textit{temporal leakage}, yaitu kondisi di mana model secara tidak sengaja memperoleh informasi dari masa depan pada saat proses validasi.
\begin{equation}
\mathbf{D}_{\text{sorted}} = \operatorname{sort}(\mathbf{D},\;[\texttt{Year},\texttt{Month}],\;\text{ascending})
\end{equation}

\noindent di mana $\mathbf{D}_{\text{sorted}}$ adalah dataset yang telah diurutkan secara kronologis, $\mathbf{D}$ adalah dataset mentah sebelum proses pengurutan, $\operatorname{sort}(\cdot)$ merupakan fungsi pengurutan yang diterapkan pada dataset, $[\texttt{Year}, \texttt{Month}]$ adalah daftar kunci pengurutan yang digunakan secara hierarkis yaitu pertama berdasarkan tahun kemudian bulan dalam tahun yang sama, serta \textit{ascending} menunjukkan arah pengurutan dari nilai terkecil ke terbesar atau kronologis dari waktu paling awal ke paling akhir. Tabel 4.4 menunjukkan data yang sudah terurut berdasarkan bulan dan tahun. 

\begin{table}[H]
  \centering
  \caption{Pratinjau 5 Baris Pertama Data Terurut Berdasarkan \textit{Year} dan \textit{Month}}
  \label{tab:preview_df_terurut}
  \resizebox{\textwidth}{!}{%
    \begin{tabular}{cclllcrrr}
      \toprule
      \textbf{Year} &
      \textbf{Month} &
      \textbf{Client} &
      \textbf{ClientConnection} &
      \textbf{Supplier} &
      $\cdots$ &
      \textbf{TotalSentMessages} &
      \textbf{TotalDeliveredMessages} &
      \textbf{TotalUndeliveredMessages} \\
      \midrule
      2023 & 1 & CLIENT\_030 & CLIENTCONNECTION\_071 & SUPPLIER\_001 & $\cdots$ & 8{,}516{,}311 & 7{,}900{,}562 & 143{,}330 \\
      2023 & 1 & CLIENT\_013 & CLIENTCONNECTION\_028 & SUPPLIER\_001 & $\cdots$ & 7{,}925{,}994 & 6{,}549{,}499 & 819{,}832 \\
      2023 & 1 & CLIENT\_013 & CLIENTCONNECTION\_028 & SUPPLIER\_001 & $\cdots$ & 7{,}925{,}501 & 6{,}548{,}783 & 820{,}443 \\
      2023 & 1 & CLIENT\_013 & CLIENTCONNECTION\_028 & SUPPLIER\_001 & $\cdots$ & 7{,}363{,}296 & 6{,}180{,}515 & 867{,}061 \\
      2023 & 1 & CLIENT\_013 & CLIENTCONNECTION\_028 & SUPPLIER\_001 & $\cdots$ & 7{,}357{,}229 & 6{,}174{,}895 & 866{,}668 \\
      \bottomrule
    \end{tabular}%
  }
  \par\vspace{4pt}

\end{table}

\subsection{Penghapusan Kolom yang Tidak Diperlukan}
Setelah struktur data terbentuk, beberapa kolom yang tidak digunakan pada analisis inti dihapus, khususnya atribut mata uang yang redundan terhadap target analisis, seperti \texttt{PriceCurrency}, \texttt{ProfitCurrency}, \texttt{BaseCurrency}, dan \texttt{CostCurrency}. Atribut mata uang tersebut dihapus karena memiliki distribusi yang sangat tidak seimbang (\textit{imbalanced}), di mana lebih dari 99\% nilai merupakan mata uang USD. Sisanya bernilai kosong karena kesalahan input data. Kondisi ini menyebabkan kolom tersebut bersifat \textit{quasi-constant} (hampir semuanya bernilai sama di tiap data) sehingga tidak memiliki nilai informatif untuk proses prediksi. Pada tahap eksplorasi tertentu, kolom \texttt{Supplier} juga dieliminasi untuk menyesuaikan kebutuhan pemodelan dan visualisasi karena data dalam kolom tersebut tidak diperlukan yang mana memiliki data yang sama di tiap barisnya. Kemudian kolom \texttt{Country} juga dihapus karena informasi dalam kolom tersebut sudah terwakili oleh kolom \texttt{Network} yang mana kolom tersebut mewakili negara tujuan transaksi dan tiap negara memiliki jenis \texttt{Network} yang berbeda sehingga penghapusan kolom tersebut menghindari data redundan. 

\begin{longtable}{lll}
    \caption{Deskripsi Data Pendapatan SMS Internasional} \\
    \toprule
    \textbf{Fitur} & \textbf{Tipe Data} & \textbf{Keterangan} \\
    \midrule
    \endfirsthead
    
    \toprule
    \textbf{Fitur} & \textbf{Tipe Data} & \textbf{Keterangan} \\
    \midrule
    \endhead
    
    Year & Integer & Tahun transaksi \\
    Month & Integer & Bulan transaksi \\
    Client & String & Nama klien \\
    Client Connection & String & Koneksi klien \\
    Supplier Connection & String & Koneksi supplier \\
    Supplier Connection Link & String & Link koneksi supplier \\
    Route & String & Rute pengiriman \\
    Country & String & Negara \\
    Network & String & Jaringan distribusi \\
    OAProfiled Category & String & Kategori profil OA \\
    Total Cost & Integer & Total biaya \\
    Total Sent Messages & Integer & Total pesan terkirim \\
    Total Delivered Messages & Integer & Total pesan berhasil terkirim \\
    \textbf{Total Revenue} & \textbf{Integer} & \textbf{Variabel target penelitian} \\
    \bottomrule
    \end{longtable}
    
\subsection{Pembersihan Nilai Tidak Valid}
Nilai numerik pada beberapa kolom masih mengandung pemisah ribuan (koma) dan spasi, sehingga perlu dinormalisasi sebelum proses statistik dan pelatihan model. Tahapan pembersihan meliputi penghapusan karakter non-numerik relevan dan konversi ke tipe numerik dengan penanganan \textit{error} menjadi \texttt{NaN}. Kolom yang diproses antara lain \texttt{TotalRevenue}, \texttt{TotalCost}, \texttt{TotalSentMessages}, \texttt{TotalDeliveredMessages}, dan \texttt{TotalUndeliveredMessages}. Nilai yang hilang kemudian diimputasi menjadi nol untuk menjaga kompatibilitas dengan algoritma berbasis pohon (\textit{tree-based models}). Contoh hasil pembersihan dapat dilihat pada Tabel 4.6.

\begin{table}[H]
\centering
\caption{Contoh Nilai Sebelum dan Sesudah Pembersihan}
\label{tab:konversi_numerik}
\begin{tabular}{llccc}
\toprule
\textbf{Kolom} & \textbf{Nilai Asli} & \textbf{Tipe Asli} & \textbf{Nilai Bersih} & \textbf{Tipe Akhir} \\
\midrule
\texttt{TotalRevenue} & \texttt{'1401484'} & \texttt{str} & 1401484 & \texttt{int64} \\
\texttt{TotalRevenue} & \texttt{' - '}     & \texttt{str} & 0         & \texttt{int64} \\
\texttt{TotalCost}    & \texttt{'1245221'} & \texttt{str} & 1245221 & \texttt{int64} \\
\texttt{TotalCost}    & \texttt{' - '}     & \texttt{str} & 0         & \texttt{int64} \\
\bottomrule
\end{tabular}
\par\vspace{4pt}
\end{table}


\section{\textit{Exploratory Data Analysis} (EDA)}
\label{sec:eda}

Tahap \textit{Exploratory Data Analysis} (EDA) dilakukan untuk memahami karakteristik data transaksi SMS internasional sebelum proses pemodelan. EDA bertujuan menghasilkan \textit{insight} bisnis dan mendukung pengambilan keputusan pada tahap \textit{feature engineering}. Fokus EDA pada penelitian ini mencakup tiga aspek utama, yaitu analisis pola temporal, analisis distribusi data, dan analisis korelasi antar fitur.

\subsection{Analisis Pola Temporal}
Analisis temporal dilakukan untuk mengidentifikasi tren naik-turun trafik dan metrik finansial pada skala bulanan. Data terlebih dahulu diurutkan berdasarkan waktu, kemudian dilakukan agregasi per periode. Misalkan $t$ menyatakan periode bulan ke-$t$ dan $\mathcal{I}_t$ adalah himpunan observasi pada periode tersebut. Agregasi metrik finansial dinyatakan sebagai berikut.

\begin{align}
R_t &= \sum_{i \in \mathcal{I}_t} r_i, \\
C_t &= \sum_{i \in \mathcal{I}_t} c_i, \\
V_t &= \sum_{i \in \mathcal{I}_t} v_i
\end{align}


Dimana $R_t$ = total \texttt{TotalRevenue} pada periode $t$,
$C_t$ = total \texttt{TotalCost} pada periode $t$,
$V_t$ = total \texttt{TotalSentMessages} pada periode $t$. Untuk melihat perubahan pendapatan antar bulan, digunakan pertumbuhan bulanan (\textit{month-over-month growth}) sebagai berikut.
\begin{equation}
\mathrm{MoM}(X_t) = \frac{X_t - X_{t-1}}{X_{t-1}} \times 100\%
\end{equation}

dengan $X_t \in \{R_t, C_t, V_t\}$. Analisis pola temporal ini membantu mengidentifikasi periode lonjakan trafik, volalitas pendapatan, dan kestabilan operasional. 

\begin{figure}[H]
\centering
\includegraphics[width=0.92\textwidth]{tugas-akhir/assets/images/tren_revenue.png}
\caption{Pola tren bulanan \texttt{TotalRevenue}}
\label{fig:monthly_revenue_trend}
\end{figure}

\begin{figure}[H]
\centering
\includegraphics[width=0.92\textwidth]{tugas-akhir/assets/images/tren_cost.png}
\caption{Pola tren bulanan \texttt{TotalCost}}
\label{fig:monthly_cost_trend}
\end{figure}

\begin{figure}[H]
\centering
\includegraphics[width=0.92\textwidth]{tugas-akhir/assets/images/tren_sent_messages.png}
\caption{Pola tren bulanan \texttt{TotalSentMessages}}
\label{fig:monthly_sentmessages_trend}
\end{figure}

\begin{figure}[H]
\centering
\includegraphics[width=0.92\textwidth]{tugas-akhir/assets/images/tren_mom_growth.png}
\caption{Pola tren bulanan \texttt{TotalSentMessages}}
\label{fig:monthly_monthgrowth_trend}
\end{figure}

Gambar 4.1, 4.2, dan 4.3 menunjukkan pola pertumbuhan bulanan dari kolom target \texttt{TotalRevenue} beserta fitur \texttt{TotalCost} dan \texttt{TotalSentMessages}. Ketiga variabel memperlihatkan tren yang bergerak searah, mengindikasikan adanya korelasi positif yang kuat. Kondisi ini menunjukkan adanya potensi \textit{data leakage} apabila ketiga fitur tersebut digunakan langsung sebagai prediktor untuk memprediksi \texttt{TotalRevenue}. \textit{Data leakage} merupakan kondisi di mana data tertentu secara tidak sengaja memberikan contekan informasi ketika pelatihan model sehingga model terlalu optimis ketika memprediksi data test. Hal ini juga dikuatkan oleh \textit{domain knowledge} yang mana volume pesan yang terkirim dan total biaya secara inheren berkorelasi positif dengan pendapatan yang dihasilkan dari setiap transaksi.
Sedangkan Gambar 4.4 memperkuat temuan tersebut melalui analisis laju pertumbuhan bulanan (\textit{Month-over-Month}, MoM) dari \texttt{TotalRevenue} dan \texttt{TotalCost}. Kedua kurva bergerak hampir beriringan sepanjang periode pengamatan, dengan lonjakan dan penurunan yang terjadi secara bersamaan. Pola sinkronisasi ini menunjukkan bahwa perubahan biaya operasional pada suatu bulan cenderung diikuti oleh perubahan pendapatan pada bulan yang sama, sehingga semakin memperkuat hubungan antara kedua variabel. Selain itu, fluktuasi MoM yang cukup tinggi pada beberapa periode mempresentasikan adanya pola musiman yang perlu dipertimbangkan dalam proses rekayasa fitur temporal.

\subsection{Analisis Distribusi Data}
Analisis distribusi dilakukan untuk memahami sebaran nilai pada variabel kategorikal. Pada analisis ini digunakan \textit{bar chart} untuk melihat distribusi data kategori pada atribut seperti \texttt{Country}, \texttt{Client}, \texttt{Network}, dan \texttt{Route}. Distribusi data kategori dapat dilihat pada persamaan 4.6.
\begin{equation}
\texttt{Count}(c_j) = \sum_{i=1}^{n} \mathbb{I}(C_i = c_j)
\end{equation}

dengan $c_j$ adalah kategori ke-$j$ pada variabel kategorikal $C$. Hasil analisis distribusi menjadi dasar untuk menentukan langkah yang tepat dalam pra-pemrosesan data nantinya.

\begin{figure}[H]
\centering
\includegraphics[width=0.7\textwidth]{tugas-akhir/assets/images/distribution-networks.png}
\caption{Distribusi frekuensi transaksi berdasarkan jenis jaringan (\textit{Network})}
\label{fig:distribution_network}
\end{figure}

Gambar 4.5 menunjukkan bahwa \texttt{network-007} memiliki frekuensi transaksi yang jauh lebih tinggi dibandingkan jaringan lainnya, mempresentasikan ketimpangan distribusi yang signifikan. \textit{Insight} yang diperoleh adalah \texttt{network-007} merupakan jaringan yang paling dominan dalam transaksi. Maka dari itu, dapat dianalisis  \texttt{network-007}  merupakan pelanggan utama pada layanan transaksi SMS internasional pada perusahaan X. 

\subsection{Analisis Korelasi Antar Fitur}
Hubungan antar fitur numerik dianalisis menggunakan koefisien korelasi \textit{Pearson}, kemudian divisualisasikan dengan \textit{heatmap}. Nilai korelasi digunakan untuk menilai kekuatan hubungan linear antarvariabel terhadap target \texttt{TotalRevenue}. Koefisien \textit{Pearson} antara dua variabel $X$ dan $Y$ dapat ditulis pada persamaan 4.7.

\begin{equation}
\rho_{X,Y} = \frac{\operatorname{Cov}(X,Y)}{\sigma_X \sigma_Y}
\end{equation}

Dengan nilai $\rho_{X,Y}$ mendekati 1 menunjukkan hubungan positif kuat, mendekati -1 menunjukkan hubungan negatif kuat, dan mendekati 0 menunjukkan hubungan linear lemah.

\begin{figure}[H]
\centering
\includegraphics[width=0.7\textwidth]{tugas-akhir/assets/images/analisis-korelasi-antar-fitur.png}
\caption{Nilai Korelasi (hubungan) di tiap fitur numerik}
\label{fig:alur_prapra-pemrosesan}
\end{figure}

Pada gambar 4.6 menunjukkan bahwa korelasi fitur dari keempat fitur numerik yaitu \texttt{TotalCost,TotalSentMessages,TotalDeliveredMessages,TotalUndeliveredMessages}, dengan kolom target sangat tinggi sehingga hal ini menunjukkan data \textit{leakage} jika dipakai untuk proses  \textit{modelling}. Maka dari itu, keempat fitur numerik tersebut tidak digunakan langsung untuk proses \textit{modelling} nantinya.

\subsection{Implikasi EDA terhadap Feature Engineering}
Temuan EDA digunakan sebagai landasan penyusunan fitur model. Pola temporal mendorong pembentukan fitur lag (misalnya \texttt{Prev\_Month\_Volume} dan \texttt{Prev\_Month\_Revenue}), pola distribusi mendukung keputusan pembersihan data dan \textit{encoding kategorikal}, sedangkan pola korelasi membantu memilih fitur yang informatif sekaligus menghindari fitur yang berpotensi memicu \textit{data leakage}. Dengan demikian, EDA tidak hanya berfungsi sebagai tahap deskriptif, tetapi juga sebagai tahapan analitis untuk memastikan fitur yang dibangun relevan dengan data aktual.

% \subsection{Ringkasan Tahapan EDA}
% \begin{table}[htbp]
% \centering
% \caption{Ringkasan Tahapan Exploratory Data Analysis}
% \label{tab:ringkasan_eda}
% \footnotesize
% \begin{tabularx}{\textwidth}{l X X}
% \toprule
% Tahap EDA & Tujuan & Output Analisis \\
% \midrule
% Analisis pola temporal & Mengamati dinamika data lintas waktu & Tren bulanan, lonjakan/penurunan, dan MoM growth \\
% Analisis distribusi data & Memahami sebaran kategori dan numerik & Dominasi kategori, indikasi skewness/outlier \\
% Analisis korelasi antar fitur & Menilai keterkaitan linear antarmetrik & Matriks korelasi, fitur yang paling relevan ke target \\
% % Implikasi ke feature engineering & Menurunkan keputusan fitur berbasis data & Fitur lag, rasio, dan seleksi fitur terarah \\
% \bottomrule
% \end{tabularx}
% \end{table}

% Berdasarkan hasil EDA, data menunjukkan pola yang cukup informatif untuk pemodelan prediktif, sekaligus memberikan dasar yang kuat dalam perancangan fitur dan validasi hipotesis bisnis pada domain trafik SMS International.


\section{Pra-pemrosesan Data}
Tahap pra-pemrosesan data dilakukan untuk memproses data kembali setelah dilakukan proses analisis pola data berdasarkan proses sebelumnya. Proses ini bertujuan agar data yang dimasukkan ke pelatihan model sudah sesuai dan dapat memudahkan model untuk memprediksi kolom target yaitu  \texttt{TotalRevenue}. Proses pra-pemrosesan data terdiri dari penggabungan data berdasarkan bulan dan \textit{feature engineering}.

\subsection{Penggabungan Data Berdasarkan Bulan}
Pada tahap ini dilakukan penggabungan data transaksi ke dalam skala bulanan berdasarkan \texttt{Series\_ID} yang merepresentasikan kombinasi unik \texttt{Client}, \texttt{Route}, \texttt{SupplierConnection}, dan \texttt{Network}. Penggabungan ini diperlukan karena data asli memiliki banyak baris per bulan yang dibedakan oleh \texttt{OAProfiledCategory} sehingga model \textit{tree-based} akan kesulitan mempelajari pola jika harus memprediksi pada tingkat detail tersebut. Dengan diagregasi per bulan, setiap \texttt{Series\_ID} menghasilkan satu observasi per periode bulanan yang bermakna secara bisnis.
Pada proses penggabungan, \texttt{TotalRevenue} dijumlahkan (\textit{sum}) sebagai kolom target prediksi. Kolom \texttt{Client}, \texttt{Route}, \texttt{SupplierConnection}, dan \texttt{Network} digabungkan berdasarkan nilainya identik untuk seluruh baris dalam satu grup sehingga digunakan sebagai kunci grup (\textit{grouping key}) dalam proses \texttt{groupby}. Kolom \texttt{Country} dikeluarkan dari kunci grup karena satu \texttt{Series\_ID} dapat mencakup beberapa \texttt{Country}. Apabila menyertakannya akan memecah agregasi secara tidak diinginkan. Fitur numerik volume yaitu \texttt{TotalCost}, \texttt{TotalProfit}, \texttt{TotalSentMessages}, \texttt{TotalDeliveredMessages}, dan \texttt{TotalUndeliveredMessages} diagregasi dengan \textit{sum} karena secara semantik merupakan akumulasi bulanan yang perlu dijumlahkan. Fitur penanda anomali \texttt{is\_anomaly\_rev\_only} dan \texttt{is\_anomaly\_cost\_only} juga diagregasi dengan \textit{sum} karena jika menggunakan \textit{mode} atau \textit{first}, informasi frekuensi hilang. Sebagai contoh, dua transaksi bonus dalam satu bulan akan terlihat identik dengan satu transaksi bonus. Dengan \textit{sum}, fitur tersebut menjadi dapat menangkap frekuensi kemunculan anomali sehingga lebih informatif bagi model untuk mengenali pola pendapatan yang tidak biasa.
Fitur \texttt{OAProfiledCategory} diagregasi dengan \textit{nunique} menjadi \texttt{unique\_OA\_count} yang menghitung jumlah kategori \textit{OA Profiled} unik dalam satu bulan. Fitur ini merepresentasikan keragaman sub-segmen dalam suatu \texttt{Series\_ID}. Statistik distribusi pendapatan antar kategori OA juga dihitung melalui pre-agregasi, yaitu \texttt{max\_OA\_revenue} dan \texttt{min\_OA\_revenue} masing-masing merupakan pendapatan tertinggi dan terendah dari semua sub-transaksi OA dalam satu grup, memberikan gambaran sebaran data pendapatan per kategori.
Selain itu, dilakukan pembentukan panel lengkap (\textit{complete panel}) di mana setiap \texttt{Series\_ID} memiliki baris untuk semua bulan dalam dataset melalui \textit{cross join} antara daftar \texttt{Series\_ID} unik dan semua kombinasi \textit{Year}--\textit{Month}. Operasi ini menghasilkan semua pasangan \texttt{Series\_ID}--\texttt{Month} yang mungkin, sehingga setiap \texttt{Series\_ID} terwakili di setiap periode meskipun tidak memiliki transaksi pada bulan tersebut. Bulan tanpa transaksi diisi dengan nilai nol pada seluruh kolom numerik. Langkah ini diperlukan karena data asli hanya berisi bulan yang memiliki transaksi, sehingga fitur \textit{lag} akan berisi nilai yang kurang tepat jika bulan kosong tidak direpresentasikan. Berikut merupakan hasil dari penggabungan data transaksi berdasarkan bulan yang dapat dilihat pada Tabel~\ref{tab:monthly_agg}.
\begin{table}[H]
\centering
\caption{Preview Data Agregasi Bulanan}
\label{tab:monthly_agg}
\resizebox{\textwidth}{!}{%
\begin{tabular}{ccrrrrrr}
\toprule
Series\_ID & Year & Month & TotalRevenue & $\cdots$ & unique\_OA\_count & max\_OA\_rev \\
\midrule
CLIENT\_017\_ROUTE\_020\_.. & 2023 & 1 & 610446 & $\cdots$ & 3 & 280000 \\
CLIENT\_017\_ROUTE\_020\_.. & 2023 & 2 & 3153037 & $\cdots$ & 5 & 1500000 \\
CLIENT\_017\_ROUTE\_020\_.. & 2023 & 3 & 2861227 & $\cdots$ & 4 & 1200000 \\
CLIENT\_017\_ROUTE\_020\_.. & 2023 & 4 & 2755164 & $\cdots$ & 2 & 1400000 \\
CLIENT\_017\_ROUTE\_020\_.. & 2023 & 5 & 2716086 & $\cdots$ & 3 & 1350000 \\
\midrule
$\vdots$ & $\vdots$ & $\vdots$ & $\vdots$ & $\cdots$ & $\vdots$ & $\vdots$ \\
\bottomrule
\end{tabular}%
}
\end{table}

\subsection{Rekayasa Fitur (Feature Engineering)}

Pembuatan fitur pada penelitian ini dilakukan dalam tiga tahap yaitu sebelum agregasi bulanan, saat agregasi, dan setelah agregasi. Pada tahap sebelum agregasi, dibangun beberapa fitur pada data transaksi asli yang diperlukan untuk proses penggabungan dan pelabelan anomali.
Pertama, fitur \texttt{Series\_ID} dibentuk dari kombinasi \texttt{Client}, \texttt{Route}, \texttt{SupplierConnection}, dan \texttt{Network} untuk mengelompokkan data (\textit{grouping}) berdasarkan pola waktu masing-masing grup.
\begin{equation}
\texttt{Series\_ID} = \texttt{Client} \oplus \texttt{Route} \oplus \texttt{SupplierConnection} \oplus \texttt{Network}
\end{equation}
\noindent Kolom \texttt{Supplier} tidak disertakan karena hanya memiliki satu jenis data. Kolom \texttt{Country} juga tidak disertakan karena hanya empat baris dari 128 baris yang memiliki \texttt{Country} berbeda (East Timor) dari mayoritas Indonesia pada \texttt{Series\_ID} yang sama, sehingga menyertakannya akan memecah agregasi dan menghasilkan nilai duplikat.
Kedua, fitur penanda anomali dibuat untuk mengidentifikasi transaksi yang memiliki pola yang aneh dan tidak biasa. Pola \textit{bonus} ditandai dengan \texttt{is\_anomaly\_rev\_only} ketika $\texttt{TotalRevenue} \neq 0$, $\texttt{TotalCost} = 0$, dan $\texttt{TotalProfit} = 0$. Pola \textit{testing} ditandai dengan \texttt{is\_anomaly\_cost\_only} ketika $\texttt{TotalRevenue} = 0$, $\texttt{TotalCost} \neq 0$, dan $\texttt{TotalProfit} \neq 0$. Kedua fitur ini bersifat biner ($0$ atau $1$) dan nantinya diagregasi dengan \textit{sum} pada proses penggabungan bulanan.

Setelah penggabungan data berdasarkan bulan, saat aggregasi dibuat fitur-fitur turunan (\textit{derived}) yang tidak dapat dihitung sebelum agregasi karena memerlukan nilai yang telah dijumlahkan. Pertama, fitur rasio biaya per pesan (\texttt{rateCost}) dihitung sebagai berikut.
\begin{equation}
\texttt{rateCost} = 
\begin{cases}
\dfrac{\sum\texttt{TotalCost}}{\sum\texttt{TotalDeliveredMessages}} & \text{jika } \sum\texttt{TotalDeliveredMessages} > 0 \\[8pt]
0 & \text{jika } \sum\texttt{TotalDeliveredMessages} = 0
\end{cases}
\end{equation}
\noindent Perhitungan \texttt{rateCost} dari hasil penjumlahan (\textit{SUM-derived}) lebih akurat dibanding \textit{mean} dari rasio per baris karena memberikan bobot lebih pada baris dengan volume pesan lebih tinggi (\textit{weighted average}).
Kedua, fitur \textit{profit margin} mengukur rasio keuntungan terhadap pendapatan.
\begin{equation}
\texttt{profit\_margin} = 
\begin{cases}
\dfrac{\sum\texttt{TotalProfit}}{\sum\texttt{TotalRevenue}} & \text{jika } \sum\texttt{TotalRevenue} > 0 \\[8pt]
0 & \text{jika } \sum\texttt{TotalRevenue} = 0
\end{cases}
\end{equation}
Ketiga, fitur \textit{delivery rate} menangkap rasio keberhasilan pengiriman pesan, yang menunjukkan performa perusahaan dalam mengelola \textit{traffic} SMS internasional.
\begin{equation}
\texttt{delivery\_rate} = 
\begin{cases}
\dfrac{\sum\texttt{TotalDeliveredMessages}}{\sum\texttt{TotalSentMessages}} & \text{jika } \sum\texttt{TotalSentMessages} > 0 \\[8pt]
0 & \text{jika } \sum\texttt{TotalSentMessages} = 0
\end{cases}
\end{equation}
Keempat, fitur \texttt{anomaly\_count} menghitung total frekuensi kemunculan anomali dalam satu bulan.
\begin{equation}
\texttt{anomaly\_count} = \texttt{is\_anomaly\_rev\_only} + \texttt{is\_anomaly\_cost\_only}
\end{equation}
\noindent di mana \texttt{is\_anomaly\_rev\_only} dan \texttt{is\_anomaly\_cost\_only} merupakan hasil agregasi \textit{sum} dari fitur biner pada data transaksi asli. Terakhir, fitur \texttt{is\_active} ditambahkan sebagai penanda biner yaitu bernilai $1$ jika $\texttt{TotalRevenue} > 0$ dan $0$ jika tidak, sehingga model dapat mempelajari pola \textit{idle}.


Setelah aggregasi, dibangun fitur temporal berbasis \textit{lag} dan \textit{rolling window} pada data bulanan yang telah diurutkan per \texttt{Series\_ID}. Pertama, fitur \textit{lag} menangkap nilai \texttt{TotalRevenue} dan \texttt{is\_active} pada bulan sebelumnya.
\begin{equation}
\texttt{rev\_lag}_{k} = \texttt{TotalRevenue}_{t-k}, \quad k \in \{1, 2, 3\}
\end{equation}
\begin{equation}
\texttt{active\_lag}_{k} = \texttt{is\_active}_{t-k}, \quad k \in \{1, 2, 3\}
\end{equation}
Kedua, fitur \textit{rolling} menangkap tren dan volatilitas pendapatan menggunakan rata-rata bergerak (\textit{rolling mean}) dan simpangan baku bergerak (\textit{rolling standard deviation}) dengan \textit{shift} satu bulan untuk mencegah kebocoran data (\textit{data leakage}).
\begin{equation}
\texttt{rolling\_mean}_{w} = \frac{1}{w}\sum_{i=1}^{w}\texttt{TotalRevenue}_{t-i}, \quad w \in \{3, 6\}
\end{equation}
\begin{equation}
\texttt{rolling\_std}_{w} = \sqrt{\frac{1}{w}\sum_{i=1}^{w}\left(\texttt{TotalRevenue}_{t-i} - \texttt{rolling\_mean}_{w}\right)^{2}}, \quad w \in \{3, 6\}
\end{equation}
Proporsi bulan aktif dalam \textit{window} $w$ dihitung sebagai berikut.
\begin{equation}
\texttt{pct\_active}_{w} = \frac{1}{w}\sum_{i=1}^{w}\texttt{is\_active}_{t-i}, \quad w \in \{3, 6\}
\end{equation}
Ketiga, fitur \texttt{active\_streak} menghitung jumlah bulan berturut-turut \textit{active} hingga bulan sebelumnya, sedangkan \texttt{months\_since\_active} menghitung jumlah bulan sejak terakhir \textit{active}. Terakhir, variabel \textit{month} dienkode secara siklikal untuk menangkap pola musiman.
\begin{equation}
\texttt{month\_sin} = \sin\left(\frac{2\pi \cdot \texttt{Month}}{12}\right), \quad \texttt{month\_cos} = \cos\left(\frac{2\pi \cdot \texttt{Month}}{12}\right)
\end{equation}

% \subsection{Analisis Domain Knowledge}
% Kemudian dibentuk sebuah fitur untuk menangani kondisi anomali pada pola data dengan analisis \textit{domain knowledge}. Contoh kondisi anomali yang diperiksa meliputi ketidakkonsistenan volume pesan dan ketidaksesuaian antara trafik dan nilai finansial. Secara konseptual, indikator anomali dapat dinyatakan sebagai berikut.

% \begin{equation}
% Delete_i =
% \begin{cases}
% \text{True}, & \text{jika } R_i = 0 \land C_i = 0 \land P_i = 0 \\
% \text{False}, & \text{lainnya}
% \end{cases}
% \end{equation}

% \begin{equation}
% \text{is\_anomaly\_cost\_only}_i =
% \begin{cases}
% 1, & \text{jika } R_i = 0 \land (C_i \neq 0 \lor P_i \neq 0) \\
% 0, & \text{lainnya}
% \end{cases}
% \end{equation}

% \begin{equation}
% \text{is\_anomaly\_rev\_only}_i =
% \begin{cases}
% 1, & \text{jika } R_i \neq 0 \land C_i = 0 \land P_i = 0 \\
% 0, & \text{lainnya}
% \end{cases}
% \end{equation}

% Dengan analisis \textit{domain knowledge} didapatkan alasan pemberian nilai tersebut. Jika \texttt{TotalRevenue}, \texttt{TotalCost}, \texttt{TotalProfit} bernilai 0 maka transaksi tersebut merupakan transaksi gagal sehingga nilai yang berpola tersebut bisa dihapus. Kemudian jika \texttt{TotalRevenue} bernilai tidak nol tetapi \texttt{TotalCost} dan \texttt{TotalProfit} bernilai nol, transaksi tersebut dianggap sebagai bonus sehingga nilai tersebut dipertahankan. Untuk menandai transaksi tersebut, ditambahkan fitur baru \texttt{is\_anomaly\_rev\_only} dengan bernilai 1 jika terdapat pola tersebut dan 0 jika tidak. Sedangkan jika \texttt{TotalRevenue} bernilai nol tetapi \texttt{TotalCost} dan \texttt{TotalProfit} memiliki nilai yang berkebalikan, transaksi tersebut dianggap transaksi yang digunakan untuk \textit{testing} (percobaan). Perusahaan telekomunikasi dalam penelitian ini, ketika menyediakan layanan ke beberapa \textit{client}, sudah terbiasa untuk melakukan proses \textit{traffic} SMS secara \textit{testing} . Tujuannya proses \textit{traffic} utama bisa berjalan lebih lancar untuk perusahaan tertentu.  Sama seperti fitur sebelumnya, untuk menandai transaksi tersebut ditambahkan fitur baru \texttt{is\_anomaly\_cost\_only} dengan bernilai 1 jika terdapat pola tersebut dan 0 jika tidak.

\begin{table}[H]
  \centering
  \caption{Preview Data Setelah Proses \textit{Feature Engineering} (5 Baris Pertama)}
  \label{tab:df_final_preview}
  \resizebox{\textwidth}{!}{%
  \begin{tabular}{c|c|c|c|c|c|c|c}
    \hline
    Client\_enc & $\cdots$ & rateCost & is\_anomaly\_rev\_only & is\_anomaly\_cost\_only & Series\_ID & Month \& Year & TotalRevenue \\
    \hline
    28 & $\cdots$ & 0.137437 & 0 & 0 & CLIENT\_030\_ROUTE\_037\_SC\_001 & 1, 2023 & 1{,}098{,}604 \\
    12 & $\cdots$ & 0.104801 & 0 & 0 & CLIENT\_013\_ROUTE\_016\_SC\_004 & 1, 2023 & 810{,}829 \\
    12 & $\cdots$ & 0.104805 & 0 & 0 & CLIENT\_013\_ROUTE\_016\_SC\_004 & 1, 2023 & 810{,}779 \\
    12 & $\cdots$ & 0.103173 & 0 & 0 & CLIENT\_013\_ROUTE\_016\_SC\_004 & 1, 2023 & 753{,}265 \\
    12 & $\cdots$ & 0.103182 & 0 & 0 & CLIENT\_013\_ROUTE\_016\_SC\_004 & 1, 2023 & 752{,}645 \\
    \hline
  \end{tabular}%
  }
\end{table}

% =============================================================
% BAB 4.5 - PELATIHAN MODEL
% =============================================================

\section{Pelatihan Model}
\label{subsec:pelatihan_model}

Tahap pelatihan model pada penelitian ini menggunakan pendekatan \textit{two-stage model}. Tahap pertama (\textit{Stage 1}) adalah klasifikasi biner untuk memprediksi apakah suatu \texttt{Series\_ID} akan aktif pada bulan berikutnya dan tahap kedua (\textit{Stage 2}) adalah regresi untuk memprediksi nilai \texttt{TotalRevenue} hanya pada \texttt{Series\_ID} yang diprediksi aktif (pendapatan tidak nol). Pendekatan ini dipilih karena sebagian besar \texttt{Series\_ID} mengalami periode \textit{idle} (tidak aktif), sehingga model regresi yang memprediksi pendapatan secara langsung tanpa tahap penyaringan akan mengalami bias yang signifikan. Prediksi akhir merupakan hasil perkalian keluaran klasifikasi dengan keluaran regresi. Apabila $P = 0$ maka prediksi akhir bernilai nol, sedangkan apabila $P = 1$ maka prediksi akhir sama dengan keluaran regresi. Skema pelatihan model secara keseluruhan ditunjukkan pada Gambar 4.7.

\begin{figure}[htbp]
\centering
\includegraphics[width=0.92\textwidth]{tugas-akhir/assets/images/skemapelatihanmodel.png}
\caption{Skema pendekatan \textit{two-stage model} untuk prediksi \textit{TotalRevenue}}
\label{fig:skema_pelatihan_model}
\end{figure}

Data dibagi secara kronologis menggunakan kolom \textit{period} ($\texttt{Year} \times 12 + \texttt{Month}$). Titik pemisahan (\textit{cutoff}) ditetapkan pada persentil ke-80 dari kolom \textit{period}, menghasilkan sekitar 80\% data latih dan 20\% data uji dan secara tepatnya ada 1729 data latih dan 364 data uji. Pendekatan ini memastikan bahwa data latih selalu mendahului data uji secara kronologis, sehingga tidak ada informasi dari masa depan yang bocor (\textit{leakage}) ke dalam model. Kedua tahap menggunakan pembagian data yang sama, dengan penyesuaian di tahap kedua hanya memanfaatkan baris dengan \texttt{target\_active} = 1 pada data latih dan baris yang diprediksi aktif pada data uji.

Metrik evaluasi utama yang digunakan konsisten dengan fungsi \textit{fitness} GA, yaitu \textit{Weighted Mean Absolute Percentage Error} (WMAPE). WMAPE dipilih karena mengakomodasi perbedaan skala nilai pendapatan antar \texttt{Series\_ID}, sehingga kesalahan prediksi pada \texttt{Series\_ID} yang memiliki total nilai aktual bernilai besar mendapat bobot proporsional yang lebih tinggi. Selain WMAPE, metrik pendukung yang digunakan meliputi MAE dan $R^2$.

% ---------------------------------------------------------
\subsection{Klasifikasi Transaksi Aktif dan Tidak Aktif}
\label{subsubsec:stage1_klasifikasi}

Pada tahap ini \textit{tree-based models} bertugas mengklasifikasikan apakah suatu \texttt{Series\_ID} akan aktif (\texttt{target\_active} = 1) atau \textit{idle} (\texttt{target\_active} = 0) pada bulan berikutnya. Klasifikasi ini penting karena hanya \texttt{Series\_ID} yang diprediksi aktif yang akan diteruskan ke tahap kedua untuk prediksi pendapatan. \textit{tree-based models} yang digunakan pada tahap ini adalah \textit{XGBoost}. \textit{XGBoost} dipilih karena model ini memiliki parameter \textit{scale\_pos\_weight} yang penting untuk mengatasi ketidakseimbangan kelas antara \texttt{Series\_ID} aktif dan \textit{idle}, memiliki mekanisme \textit{early stopping} yang mencegah \textit{overfitting}, dan sering digunakan untuk memecahkan permasalahan \textit{machine learning} pada data tabular.

Paramater model \textit{XGBoost} yang digunakan pada tahap ini yaitu \texttt{n\_estimators} = 1000, \texttt{learning\_rate} = 0{,}3, \texttt{max\_depth} = 6, \texttt{subsample} = 1{,}0, \texttt{colsample\_bytree} = 1{,}0, \texttt{reg\_lambda} = 1{,}0, \texttt{scale\_pos\_weight} = rasio kelas (\textit{idle}/\textit{active}), \texttt{random\_state} = 42. Parameter \texttt{scale\_pos\_weight} dihitung secara otomatis sebagai rasio jumlah \texttt{Series\_ID} \textit{idle} terhadap \texttt{Series\_ID} aktif pada data latih untuk mengatasi ketidakseimbangan kelas. \textit{Early stopping} dengan \textit{patience} 50 iterasi diterapkan pada data uji untuk menghentikan pelatihan jika tidak ada peningkatan performansi model. 

Fitur yang digunakan pada tahap ini terdiri dari 20 fitur. Fitur ini merupakan hasil dari pra-pemrosesan data dan \textit{feature engineering} dari proses sebelumnya. Fitur ini dapat dikelompokkan sebagai berikut.
\begin{enumerate}
    \item Status historis : \texttt{active\_lag1}, \texttt{active\_lag2}, \texttt{active\_lag3}, \texttt{active\_streak}, \texttt{months\_since\_active}, \texttt{pct\_active\_3m}, \texttt{pct\_active\_6m}
    \item \textit{Revenue} historis : \texttt{rev\_lag1}, \texttt{rev\_lag2}, \texttt{rev\_lag3}, \texttt{rolling\_mean\_3}, \texttt{rolling\_mean\_6} 
    \item Operasional : \texttt{TotalSentMessages}, \texttt{TotalDeliveredMessages}, \texttt{delivery\_rate}, \texttt{unique\_OA\_count}
    \item Anomali : \texttt{anomaly\_count}, \texttt{anomaly\_revenue\_pct}
    \item Siklikalitas : \texttt{month\_sin}, \texttt{month\_cos}.
\end{enumerate}

Metrik evaluasi pada tahap ini menggunakan ROC-AUC dan PR-AUC yang dipilih karena data memiliki ketidakseimbangan kelas antara transaksi aktif (994 data) dan tidak aktif (1.099 data). Hasil evaluasi menunjukkan ROC-AUC sebesar 0,9611 dan PR-AUC sebesar 0,9748, yang menunjukan performa klasifikasi yang sangat baik dalam membedakan transaksi aktif dari yang \textit{idle} meskipun terdapat ketidakseimbangan kelas.

% ---------------------------------------------------------
\subsection{Regresi Pendapatan}
\label{subsubsec:stage2_regresi}

Pada tahap ini \textit{tree-based models} bertugas memprediksi nilai pendapatan (\texttt{target\_revenue}) hanya untuk \texttt{Series\_ID} yang diklasifikasikan aktif atau bernilai bukan nol. \textit{Tree-based models} yang digunakan adalah \textit{Random Forest} dan \textit{XGBoost}. Data train untuk tahap ini disaring sehingga hanya memuat \texttt{Series\_ID} yang aktif pada data latih (\texttt{target\_active} = 1). Sedangkan pada data uji, data yang digunakan berdasarkan prediksi pada tahap pertama, sehingga mensimulasikan proses \textit{inferensi} (pengujian) yang sebenarnya.

Fitur yang digunakan pada tahap kedua identik dengan tahap pertama ditambah dengan \texttt{rateCost}, \texttt{unique\_OA\_count}, \texttt{max\_OA\_revenue} sehingga total ada 23 fitur. Hasil pelatihan model \textit{baseline} pada tahap ini menunjukkan \textit{XGBoost} mencapai WMAPE 27,68\%, $R^2$ 0,920, dan MAE 39.582,76, sedangkan \textit{Random Forest} mencapai WMAPE 24,82\%, $R^2$ 0,932, dan MAE 35.485,88. Dengan demikian, \textit{Random Forest baseline} mengungguli \textit{XGBoost baseline} pada semua metrik evaluasi. Performa kedua model \textit{baseline} ini menjadi acuan pembanding untuk mengevaluasi efektivitas optimasi hiperparameter menggunakan Algoritma Genetika.

% =============================================================
% BAB 4.6 - OPTIMASI HYPERPARAMETER MENGGUNAKAN ALGORITMA GENETIKA
% =============================================================

% \section{Optimasi Hiperparameter dengan Algoritma Genetika}
% \label{subsec:optimasi_hyperparameter}

% Untuk menentukan konfigurasi hiperparameter yang optimal, penelitian ini menerapkan Algoritma Genetika pada dua model \textit{tree-based} yaitu \textit{Random Forest} (berbasis \textit{bagging ensemble}) dan XGBoost (berbasis \textit{boosting ensemble}). Pemilihan kedua model ini memungkinkan perbandingan efektivitas GA pada dua paradigma \textit{ensemble} yang berbeda. Hal serupa dilakukan oleh Mansoori (2023) yang membandingkan performa \textit{tree-based models} dalam konteks prediksi durasi perawatan di rumah sakit.

% GA dijalankan secara independen untuk masing-masing model dengan konfigurasi 20 individu per populasi, 50 generasi, \textit{crossover rate} = 1,0, \textit{mutation rate} = 0,2, dan \textit{keep\_elitism} = 2. Setelah 50 generasi, konfigurasi hiperparameter optimal berhasil ditemukan untuk setiap model dan dibandingkan dengan performa model \textit{baseline}.

% Kedua model dilakukan evaluasi pada data uji di tahap 2 menggunakan empat model yaitu \textit{Random Forest Baseline}, \textit{Random Forest + GA}, \textit{XGBoost Baseline}, dan \textit{XGBoost + GA}. Hasil optimasi menunjukkan peningkatan yang signifikan pada kedua model. Pada \textit{Random Forest}, WMAPE turun dari 24,82\% menjadi 18,80\% (reduksi 6,02 poin persentase), $R^2$ meningkat dari 0,9315 menjadi 0,9717, dan MAE turun dari 35.485,88 menjadi 26.890,63. Pada XGBoost, WMAPE turun dari 27,68\% menjadi 15,46\% (reduksi 12,22 poin persentase), $R^2$ meningkat dari 0,9207 menjadi 0,9907, dan MAE turun dari 39.582,76 menjadi 22.107,41. Model \textit{XGBoost + GA} memberikan performa terbaik secara keseluruhan dengan WMAPE terendah (15,46\%), RMSE terendah (46.823,86), dan $R^2$ tertinggi (0,9907). Besarnya peningkatan pada XGBoost (12,22 poin persentase) yang lebih besar dibandingkan \textit{Random Forest} (6,02 poin persentase) mengindikasikan bahwa konfigurasi \textit{default} XGBoost kurang optimal untuk karakteristik data ini, sehingga XGBoost mengalami pertumbuhan performa yang signifikan dari proses optimasi hiperparameter oleh GA. Tabel \ref{tab:perbandingan_model} menyajikan perbandingan performa keempat model.

% \begin{table}[H]
% \centering
% \caption{Perbandingan Hasil Evaluasi Model Tahap Kedua}
% \label{tab:perbandingan_model}
% \begin{tabular}{llccc}
% \toprule
% \textbf{Model} & \textbf{Konfigurasi} & \textbf{WMAPE (\%)} & \textbf{$R^2$} & \textbf{MAE} \\
% \midrule
% Random Forest  & Baseline  & 24,82 & 0,9315 & 35\,485,88 \\
% Random Forest  & GA-Tuned  & 18,80 & 0,9717 & 26\,890,63 \\
% XGBoost        & Baseline  & 27,68 & 0,9207 & 39\,582,76 \\
% XGBoost        & GA-Tuned  & 15,46 & 0,9907 & 22\,107,41 \\
% \bottomrule
% \end{tabular}
% \end{table}

% Untuk mendapatkan gambaran yang lebih komprehensif mengenai distribusi kesalahan prediksi, analisis visual menggunakan \textit{boxplot} dilakukan sebagaimana ditunjukkan pada Gambar 4.8. Hasil visualisasi mengkonfirmasi bahwa model \textit{XGBoost + GA} memiliki rata-rata \textit{error} yang paling rendah dan distribusi kesalahan yang paling terpusat dibandingkan model lainnya. Untuk melengkapi analisis, indeks statistik deskriptif --- meliputi rata-rata, median, standar deviasi, \textit{interquartile range}, nilai minimum, dan maksimum --- juga dihitung untuk setiap model \cite{https://doi.org/10.1155/2023/9673395}. Gambar 4.9 menunjukkan bahwa \textit{XGBoost + GA} tetap unggul berdasarkan semua indeks statistik. Sebagai perbandingan, \textit{XGBoost baseline} memiliki nilai maksimum kesalahan yang sangat tinggi, menunjukkan ketidakstabilan prediksi yang berhasil diatasi oleh optimasi GA.

% \begin{figure}[H]
% \centering
% \includegraphics[width=0.7\textwidth]{tugas-akhir/assets/images/boxplot_error_models.png}
% \caption{Perbandingan distribusi kesalahan prediksi (\textit{boxplot error}) keempat model}
% \label{fig:boxplot_error}
% \end{figure}

% \begin{figure}[H]
% \centering
% \includegraphics[width=0.7\textwidth]{tugas-akhir/assets/images/relative_absolute_error.png}
% \caption{Indeks statistik deskriptif kesalahan prediksi seluruh model}
% \label{fig:relative_absolute_error}
% \end{figure}

% Selain analisis kesalahan, kurva konvergensi GA ditunjukkan pada Gambar 4.10. Kurva ini memperlihatkan peningkatan nilai fungsi \textit{fitness} yang konsisten di setiap generasi untuk kedua model, mengkonfirmasi bahwa GA berhasil menavigasi ruang pencarian hiperparameter secara efektif. Konfigurasi hiperparameter optimal yang dihasilkan GA untuk masing-masing model ditunjukkan pada Tabel 4.10 dan Tabel 4.11.

% \begin{figure}[H]
% \centering
% \includegraphics[width=0.92\textwidth]{tugas-akhir/assets/images/output_konvergensi_ga.png}
% \caption{Kurva konvergensi fungsi \textit{fitness} GA untuk \textit{Random Forest} dan \textit{XGBoost}}
% \label{fig:ga_convergence}
% \end{figure}

% \begin{table}[h!]
% \centering
% \caption{Hiperparameter optimal \textit{Random Forest Regressor} hasil GA}
% \label{tab:optimal_hp_rf}
% \begin{tabular}{llc}
% \toprule
% \textbf{Hiperparameter} & \textbf{Ruang Pencarian} & \textbf{Nilai Optimal} \\
% \midrule
% \texttt{n\_estimators}       & [50, 1000]    & 115   \\
% \texttt{max\_depth}          & [1, 7]         & 7     \\
% \texttt{min\_samples\_split} & [2, 20]        & 4     \\
% \texttt{min\_samples\_leaf}  & [1, 50]        & 2     \\
% \texttt{max\_features}       & [1, $N_f$]     & 6     \\
% \texttt{max\_samples}        & [0{,}1, 1{,}0] & 0{,}877 \\
% \bottomrule
% \end{tabular}
% \end{table}

% \begin{table}[h!]
% \centering
% \caption{Hiperparameter optimal \textit{XGBoost Regressor} hasil GA}
% \label{tab:optimal_hp_xgb}
% \begin{tabular}{llc}
% \toprule
% \textbf{Hiperparameter} & \textbf{Ruang Pencarian} & \textbf{Nilai Optimal} \\
% \midrule
% \texttt{learning\_rate}      & [0{,}001, 1{,}0]     & 0{,}224 \\
% \texttt{n\_estimators}       & [50, 1000]           & 462   \\
% \texttt{max\_depth}          & [1, 7]                & 5     \\
% \texttt{subsample}           & [0{,}1, 1{,}0]       & 0{,}756 \\
% \texttt{colsample\_bytree}   & [0{,}1, 1{,}0]       & 0{,}631 \\
% \texttt{reg\_lambda}         & [1, 3]                & 1{,}506 \\
% \bottomrule
% \end{tabular}
% \end{table}

\section{Optimasi \textit{Hyperparameter} dengan Algoritma Genetika}
\label{subsec:optimasi_hyperparameter}

Untuk menentukan konfigurasi hiperparameter yang optimal, penelitian ini menerapkan Algoritma Genetika pada dua model \textit{tree-based} yaitu \textit{Random Forest} (berbasis \textit{bagging ensemble}) dan XGBoost (berbasis \textit{boosting ensemble}). Pemilihan kedua model ini memungkinkan perbandingan efektivitas GA pada dua paradigma \textit{ensemble} yang berbeda. Hal serupa dilakukan oleh Mansoori (2023) yang membandingkan performa kedua model dalam konteks prediksi durasi perawatan di rumah sakit.

GA dijalankan secara independen untuk masing-masing model dengan konfigurasi yang identik yaitu 40 individu per populasi, 100 generasi, \textit{crossover rate} = 1,0, \textit{mutation rate} = 0,2, dan \textit{keep\_elitism} = 2. Setelah 100 generasi, konfigurasi \textit{hyperparameter} optimal berhasil ditemukan untuk setiap model dan dibandingkan dengan performa model \textit{baseline}.

% ---------------------------------------------------------
\subsection{Representasi Kromosom dan Ruang Pencarian}
\label{subsec:kromosom_ga}

Setiap individu dalam populasi GA direpresentasikan sebagai kromosom berupa vektor real $\boldsymbol{\theta} \in \mathbb{R}^6$, di mana masing-masing gen memetakan satu \textit{hyperparameter}. Kedua model dioptimalkan secara terpisah dengan ruang pencarian yang berbeda. Untuk lebih detailnya dapat dilihat pada Tabel~\ref{tab:gene_rf} dan \ref{tab:gene_xgb}.

\begin{table}[h!]
\centering
\caption{Ruang pencarian kromosom \textit{Random Forest Regressor}}
\label{tab:gene_rf}
\begin{tabular}{clccc}
\toprule
\textbf{Gen ke-} & \textbf{\textit{Hyperparameter}} & \textbf{Tipe} & \textbf{Min} & \textbf{Maks} \\
\midrule
0 & \texttt{n\_estimators}       & Integer & 50    & 1.000 \\
1 & \texttt{max\_depth}          & Integer & 1     & 7     \\
2 & \texttt{min\_samples\_split} & Integer & 2     & 20    \\
3 & \texttt{min\_samples\_leaf}  & Integer & 1     & 50    \\
4 & \texttt{max\_features}       & Integer & 1     & $N_f$ \\
5 & \texttt{max\_samples}        & Float   & 0{,}1 & 1{,}0 \\
\bottomrule
\end{tabular}
\end{table}

\begin{table}[h!]
\centering
\caption{Ruang pencarian kromosom \textit{XGBoost Regressor}}
\label{tab:gene_xgb}
\begin{tabular}{clccc}
\toprule
\textbf{Gen ke-} & \textbf{\textit{Hyperparameter}} & \textbf{Tipe} & \textbf{Min} & \textbf{Maks} \\
\midrule
0 & \texttt{learning\_rate}    & Float   & 0{,}001 & 1{,}0 \\
1 & \texttt{n\_estimators}     & Integer & 50      & 1.000 \\
2 & \texttt{max\_depth}        & Integer & 1       & 7     \\
3 & \texttt{subsample}         & Float   & 0{,}1   & 1{,}0 \\
4 & \texttt{colsample\_bytree} & Float   & 0{,}1   & 1{,}0 \\
5 & \texttt{reg\_lambda}       & Float   & 1       & 3     \\
\bottomrule
\end{tabular}
\end{table}

Sebagai ilustrasi proses representasi kromosom, sebuah individu dalam populasi memiliki kromosom \textit{Random Forest} dengan nilai gen kontinu sebagai berikut.

\begin{equation}
\boldsymbol{\theta} = [312{,}4;\ 5{,}2;\ 13{,}7;\ 7{,}8;\ 4{,}3;\ 0{,}62]
\end{equation}

Fungsi \texttt{decode\_rf()} mengonversi nilai gen ke \textit{hyperparameter} melalui pembulatan (\textit{round}) untuk tipe integer dan pemotongan (\textit{clip}) ke batas rentang yang valid.

\begin{align}
\texttt{n\_estimators}       &= \operatorname{clip}\!\left(\operatorname{round}(312{,}4),\ 50,\ 1000\right) = 312 \nonumber \\
\texttt{max\_depth}          &= \operatorname{clip}\!\left(\operatorname{round}(5{,}2),\ 1,\ 7\right)       = 5   \nonumber \\
\texttt{min\_samples\_split} &= \operatorname{clip}\!\left(\operatorname{round}(13{,}7),\ 2,\ 20\right)     = 14  \nonumber \\
\texttt{min\_samples\_leaf}  &= \operatorname{clip}\!\left(\operatorname{round}(7{,}8),\ 1,\ 50\right)      = 8   \nonumber \\
\texttt{max\_features}       &= \operatorname{clip}\!\left(\operatorname{round}(4{,}3),\ 1,\ N_f\right)     = 4   \nonumber \\
\texttt{max\_samples}        &= \operatorname{clip}(0{,}62,\ 0{,}1,\ 1{,}0)                               = 0{,}62
\end{align}

% Kromosom optimal yang ditemukan GA setelah 50 generasi untuk \textit{Random Forest} adalah:
% \begin{equation}
% \boldsymbol{\theta}^{*}_{\text{RF}} = [115;\ 7;\ 4;\ 2;\ 6;\ 0{,}877]
% \end{equation}
% dan untuk \textit{XGBoost}:
% \begin{equation}
% \boldsymbol{\theta}^{*}_{\text{XGB}} = [0{,}224;\ 462;\ 5;\ 0{,}756;\ 0{,}631;\ 1{,}506]
% \end{equation}

% ---------------------------------------------------------
\subsection{Fungsi Kebugaran (\textit{Fitness Function})}
\label{subsec:fitness_ga}

Fungsi kebugaran (\textit{fitness function}) pada penelitian ini didefinisikan sebagai invers dari WMAPE yang mana dapat dilihat pada persamaan 2.7. Sebagai contoh numerik, misalkan model dengan suatu kromosom menghasilkan prediksi pada tiga sampel uji sebagaimana ditunjukkan pada Tabel~\ref{tab:contoh_wmape}.  Contoh ini menggunakan tiga sampel untuk kemudahan ilustrasi, pada eksperimen sebenarnya, WMAPE dihitung dari seluruh 210 baris data uji \textit{Stage 2}.

\begin{table}[h!]
\centering
\caption{Contoh numerik perhitungan WMAPE untuk satu individu}
\label{tab:contoh_wmape}
\begin{tabular}{cccc}
\toprule
$i$ & $y_i$ (aktual, Rp) & $\hat{y}_i$ (prediksi, Rp) & $|y_i - \hat{y}_i|$ \\
\midrule
1 & 1.000.000 & 920.000   & 80.000  \\
2 &   500.000 & 580.000   & 80.000  \\
3 & 2.000.000 & 1.850.000 & 150.000 \\
\midrule
\textbf{Jumlah} & \textbf{3.500.000} & --- & \textbf{310.000} \\
\bottomrule
\end{tabular}
\end{table}

\begin{equation}
\operatorname{WMAPE} = \frac{310.000}{3.500.000} \times 100 = 8{,}857\%
\end{equation}

\begin{equation}
f(\boldsymbol{\theta}) = -8{,}857\%
\end{equation}

Nilai \textit{fitness} tersebut kemudian digunakan untuk membandingkan seluruh individu dalam populasi. Semakin tinggi nilai \textit{fitness} atau semakin mendekati nol dari sisi negatif, semakin baik konfigurasi \textit{hyperparameter} yang diwakili individu tersebut.

% ---------------------------------------------------------
\subsection{Operator Seleksi, Persilangan, dan Mutasi}
\label{subsec:operator_ga}

Pada setiap pemilihan orang tua, tiga individu dipilih secara acak dari populasi dan dibandingkan nilai \textit{fitness}-nya. Individu dengan \textit{fitness} tertinggi dipilih sebagai orang tua. Sebagai contoh, misalkan pada generasi ke-10 tiga individu berikut terpilih secara acak.

\begin{table}[H]
\centering
\caption{Contoh seleksi \textit{tournament} ($k=3$) pada generasi ke-10}
\label{tab:contoh_tournament}
\begin{tabular}{ccc}
\toprule
\textbf{Individu} & \textbf{WMAPE (\%)} & \textbf{\textit{Fitness}} \\
\midrule
Individu ke-5  & 24{,}82 & $-0{,}2482$ \\
Individu ke-12 & 19{,}43 & $-0{,}1943$ \\
Individu ke-7  & 31{,}15 & $-0{,}3115$ \\
\midrule
\textbf{Terpilih} & \multicolumn{2}{c}{\textbf{Individu ke-12} (\textit{fitness} tertinggi)} \\
\bottomrule
\end{tabular}
\end{table}

Individu ke-12 terpilih karena memiliki nilai \textit{fitness} $-0{,}1943$ yang paling tinggi di antara ketiga kandidat, yang berarti WMAPE-nya paling rendah.

Selanjutnya, dua orang tua disilangkan menggunakan \textit{scattered crossover}, di mana sebuah vektor biner acak, disebut dengan \textit{mask}, $\mathbf{m} \in \{0,1\}^6$ menentukan dari induk mana tiap gen diwariskan. Misalkan dua induk dan \textit{mask} yang dibangkitkan secara acak sebagai berikut.

\begin{align}
\boldsymbol{\theta}_A &= [200;\ 4;\ 10;\ 5;\ 3;\ 0{,}70] \nonumber \\
\boldsymbol{\theta}_B &= [450;\ 6;\ 8;\ 12;\ 5;\ 0{,}45] \nonumber \\
\mathbf{m}            &= [1;\ 0;\ 1;\ 0;\ 1;\ 0]
\end{align}

Keturunan dibentuk dengan mengambil gen dari $\boldsymbol{\theta}_A$ bila $m_j = 1$ dan dari $\boldsymbol{\theta}_B$ bila $m_j = 0$ yang dapat dilihat pada persamaan 4.24.

\begin{equation}
\boldsymbol{\theta}_C = [200;\ 6;\ 10;\ 12;\ 3;\ 0{,}45]
\end{equation}

% Keunggulan \textit{scattered crossover} dibanding \textit{one-point crossover} adalah posisi \textit{crossover} yang independen per gen, sehingga menghasilkan keragaman genetik yang lebih tinggi dan eksplorasi ruang pencarian yang lebih merata.

Kemudian, sebanyak 20\% dari total 6 gen (dibulatkan menjadi 1 gen) dipilih secara acak untuk dimutasi. Gen terpilih diganti dengan nilai baru yang diambil secara acak dari rentang ruang pencarian gen tersebut. Misalkan gen ke-2 (\texttt{min\_samples\_split}, rentang $[2, 20]$) terpilih untuk mutasi yang dapat dituliskan pada persamaan  4.25 
sehingga kromosom setelah mutasi dapat dilihat pada persamaan 4.26.
\begin{equation}
\theta_2:\ 10 \;\xrightarrow{\text{mutasi}}\; 17
\end{equation}

\begin{equation}
\boldsymbol{\theta}_C' = [200;\ 6;\ 17;\ 12;\ 3;\ 0{,}45]
\end{equation}

Mekanisme \textit{keep\_elitism} = 2 memastikan dua individu dengan \textit{fitness} terbaik dari setiap generasi dipertahankan tanpa modifikasi ke generasi berikutnya, sehingga solusi terbaik yang pernah ditemukan tidak hilang akibat operasi mutasi atau \textit{crossover}.

% ---------------------------------------------------------
\subsection{Hasil Optimasi dan Evaluasi Model}
\label{subsec:hasil_optimasi_ga}

Hasil optimasi menunjukkan peningkatan yang signifikan pada kedua model \textit{tree-based}. Pada \textit{Random Forest}, WMAPE turun dari 25,37\% menjadi 20,97\% (reduksi 4,4 poin persentase), $R^2$ meningkat dari 0,9396 menjadi 0,9544, dan MAE turun dari 36.451,29 menjadi 30.132,03. Sedangkan pada \textit{XGBoost}, WMAPE turun dari 27,12\% menjadi 16,29\% (reduksi 10,83 poin persentase), $R^2$ meningkat dari 0,9222 menjadi 0,9765, dan MAE turun dari 30.132,03, menjadi 23.398,93. Model \textit{XGBoost + GA} memberikan performa terbaik secara keseluruhan dengan WMAPE terendah (16,29\%), MAE terendah (23.398,93), dan $R^2$ tertinggi (0,9765). Besarnya peningkatan pada XGBoost yang lebih besar dibandingkan \textit{Random Forest} mengindikasikan bahwa konfigurasi \textit{default} XGBoost kurang optimal untuk karakteristik data ini, sehingga XGBoost mengalami pertumbuhan performa yang signifikan dari proses optimasi hiperparameter oleh GA. Tabel \ref{tab:perbandingan_model} menyajikan perbandingan performa keempat model.

\begin{table}[H]
\centering
\caption{Perbandingan Hasil Evaluasi Model Tahap Kedua}
\label{tab:perbandingan_model}
\begin{tabular}{llccc}
\toprule
\textbf{Model} & \textbf{Konfigurasi} & \textbf{WMAPE (\%)} & \textbf{$R^2$} & \textbf{MAE} \\
\midrule
Random Forest  & Baseline  & 25,37 & 0,9396 & 36\,451,29 \\
Random Forest  & GA-Tuned  & 20,97 & 0,9544 & 30\,132,03 \\
XGBoost        & Baseline  & 27,12 & 0,9222 & 38\,969,10 \\
XGBoost        & GA-Tuned  & 16,29 & 0,9765 & 23\,398,93 \\
\bottomrule
\end{tabular}
\end{table}

Untuk mendapatkan gambaran yang lebih komprehensif mengenai distribusi kesalahan prediksi, analisis visual menggunakan \textit{boxplot} dilakukan sebagaimana ditunjukkan pada Gambar 4.8. Hasil visualisasi mengkonfirmasi bahwa model \textit{XGBoost + GA} memiliki rata-rata \textit{error} yang paling rendah dan distribusi kesalahan yang paling terpusat dibandingkan model lainnya. Untuk melengkapi analisis, indeks statistik deskriptif meliputi rata-rata, median, standar deviasi, \textit{interquartile range}, nilai minimum, dan maksimum juga dihitung untuk setiap model \cite{https://doi.org/10.1155/2023/9673395}. Gambar 4.9 menunjukkan bahwa \textit{XGBoost + GA} tetap unggul berdasarkan semua indeks statistik. Sebagai perbandingan, \textit{XGBoost baseline} memiliki nilai maksimum kesalahan yang sangat tinggi, menunjukkan ketidakstabilan prediksi yang berhasil diatasi oleh optimasi GA.

\begin{figure}[H]
\centering
\includegraphics[width=0.7\textwidth]{tugas-akhir/assets/images/boxplot_error_models.png}
\caption{Perbandingan distribusi kesalahan prediksi (\textit{boxplot error}) keempat model}
\label{fig:boxplot_error}
\end{figure}

\begin{figure}[H]
\centering
\includegraphics[width=0.7\textwidth]{tugas-akhir/assets/images/relative_absolute_error.png}
\caption{Indeks statistik deskriptif kesalahan prediksi seluruh model}
\label{fig:relative_absolute_error}
\end{figure}

Selain analisis kesalahan, kurva konvergensi GA ditunjukkan pada Gambar 4.10. Kurva ini memperlihatkan peningkatan nilai fungsi \textit{fitness} yang konsisten di setiap generasi untuk kedua model, mengkonfirmasi bahwa GA berhasil menavigasi ruang pencarian hiperparameter secara efektif. Konfigurasi hiperparameter optimal yang dihasilkan GA untuk masing-masing model ditunjukkan pada Tabel 4.10 dan Tabel 4.11.

\begin{figure}[H]
\centering
\includegraphics[width=0.92\textwidth]{tugas-akhir/assets/images/konvergensi_ga.png}
\caption{Kurva konvergensi fungsi \textit{fitness} GA untuk \textit{Random Forest} dan \textit{XGBoost}}
\label{fig:ga_convergence}
\end{figure}

\begin{table}[H]
\centering
\caption{Hiperparameter optimal \textit{Random Forest Regressor} hasil GA}
\label{tab:optimal_hp_rf}
\begin{tabular}{llc}
\toprule
\textbf{Hiperparameter} & \textbf{Ruang Pencarian} & \textbf{Nilai Optimal} \\
\midrule
\texttt{n\_estimators}       & [50, 1000]    & 68   \\
\texttt{max\_depth}          & [1, 7]         & 6     \\
\texttt{min\_samples\_split} & [2, 20]        & 6     \\
\texttt{min\_samples\_leaf}  & [1, 50]        & 2     \\
\texttt{max\_features}       & [1, 23]     & 8     \\
\texttt{max\_samples}        & [0{,}1, 1{,}0] & 0{,}393 \\
\bottomrule
\end{tabular}
\end{table}

\begin{table}[H]
\centering
\caption{Hiperparameter optimal \textit{XGBoost Regressor} hasil GA}
\label{tab:optimal_hp_xgb}
\begin{tabular}{llc}
\toprule
\textbf{Hiperparameter} & \textbf{Ruang Pencarian} & \textbf{Nilai Optimal} \\
\midrule
\texttt{learning\_rate}      & [0{,}001, 1{,}0]     & 0{.}0417 \\
\texttt{n\_estimators}       & [50, 1000]           & 182   \\
\texttt{max\_depth}          & [1, 7]                & 4     \\
\texttt{subsample}           & [0{,}1, 1{,}0]       & 0{,}546 \\
\texttt{colsample\_bytree}   & [0{,}1, 1{,}0]       & 0{,}977 \\
\texttt{reg\_lambda}         & [1, 3]                & 1{,}195 \\
\bottomrule
\end{tabular}
\end{table}


% \begin{table}[h]
% \centering
% \caption{Perbandingan Hasil Evaluasi Model Tahap Kedua}
% \label{tab:perbandingan_model}
% \begin{tabular}{llccc}
% \toprule
% \textbf{Model} & \textbf{Konfigurasi} & \textbf{WMAPE (\%)} & \textbf{$R^2$} & \textbf{MAE} \\
% \midrule
% Random Forest  & Baseline  & 24,82 & 0,9315 & 35\,485,88 \\
% Random Forest  & GA-Tuned  & 18,80 & 0,9717 & 26\,890,63 \\
% XGBoost        & Baseline  & 27,68 & 0,9207 & 39\,582,76 \\
% XGBoost        & GA-Tuned  & 15,46 & 0,9907 & 22\,107,41 \\
% \bottomrule
% \end{tabular}
% \end{table}


    \chapter{Kesimpulan dan Saran}
    \section{Kesimpulan}

        Penelitian ini mengimplementasikan pendekatan \textit{two-stage} untuk prediksi pendapatan layanan SMS internasional pada perusahaan X menggunakan \textit{tree-based models}. Pada \textit{Stage} 1, model klasifikasi yaitu \textit{XGBoost} digunakan untuk memprediksi status aktif(pendapatan tidak nol) atau tidaknya (pendapatan nol) suatu \texttt{Series\_ID} pada bulan berikutnya. Sedangkan pada \textit{Stage} 2, model regresi yaitu \textit{Random Forest} dan \textit{XGBoost} dengan optimasi hiperparameter menggunakan Algoritma Genetika digunakan untuk memprediksi nilai pendapatan pada \texttt{Series\_ID} yang aktif saja berdasarkan hasil prediksi dari \textit{Stage} 1. 

        Penelitian ini menunjukkan bahwa optimasi hyperparameter menggunakan Algoritma Genetika dapat meningkatkan performansi dari \textit{tree-based models}. Metode ini terbukti efektif untuk meningkatkan prediksi \textit{tree-based Models} pada kasus prediksi pendapatan layanan SMS Internasional pada perusahaan X. Hasil menunjukkan bahwa XGBoost mengungguli model yang lain yaitu \textit{Random Forest} untuk performansi prediksi pendapatan layanan SMS Internasional pada perusahaan X dengan nilai WMAPE sebesar 16.29\%, MAE sebesar 23.398, 93 dan R\^2 sebesar 0.9765. Hasil ini memilik peningkatan siginifikan pada performansi model dengan penurunan WMAPE sebesar 10.83 \% dibandingkan dengan model baselinenya. Model XGBoost juga menampilkan tingkat konvergensi yang baik pada tiap generasinya. Akan tetapi, penelitian ini belum membandingkan hasil Algoritma Genetika dengan algoritma metaheuristik lain seperti \textit{Particle Swarm Optimization} (PSO), serta belum mengeksplorasi penerapan Algoritma Genetika untuk optimasi model \textit{deep learning}.


    \section{Saran}

        Setelah melakukan berbagai macam eksperimen pada penelitian ini, penulis menyarankan beberapa hal untuk penelitian selanjutnya.  Pertama, membandingkan dengan algoritma mataheuristik yang lain seperti PSO (\textit{Particle Swarm Optimization}) atau \textit{Bayesian Optimization} untuk mengetahui performansi algortima genetika dibanding dengan metodologi optimasi hiperparameter lain. Kedua, mengeksplorasi penambahan fitur eksternal seperti indikator ekonomi, pola musiman, atau kebijakan regulasi perusahaan telekomunikasi internasional untuk meningkatkan akurasi model. Terakhir, membandingkan pendekatan \textit{tree-based models} dengan model \textit{neural network} untuk mengetahui pendekatan yang lebih kompleks mampu memberikan hasil prediksi yang lebih baik. 

    \renewcommand\bibname{Daftar Pustaka}
\bibliographystyle{apacite}
\bibliography{tugas-akhir/assets/pustaka}
    \chapter*{Lampiran}\label{ch:appendix}

\section*{Bukti Email Persetujuan Penggunaan Data}

\begin{figure}[H]
\centering
\includegraphics[width=0.7\textwidth]{tugas-akhir/assets/images/buktipersetujuanpenggunaandata.png}
\caption{Bukti Email Persetujuan Penggunaan Data}
\label{fig:plot_forecast}
\end{figure}


\section*{Preview Dataset 50 Data Pertama}\label{ch:lampiran1}

\begin{landscape}
\begin{center}
\captionof{table}{Contoh Data \textit{SMS International} (50 dari 94.335 baris)}
\vspace{12pt}
\begin{adjustbox}{max width=\linewidth}
\renewcommand{\arraystretch}{1.2}
\begin{tabular}{cccccccccccccccrrrrr r}
\toprule
\textbf{Year} & \textbf{Month} & \textbf{Client} & \textbf{ClientConnection} & \textbf{Supplier} & \textbf{SupplierConnection} & \textbf{SupplierConnectionLink} & \textbf{Route} & \textbf{Country} & \textbf{Network} & \textbf{OAProfiledCategory} & \textbf{CostCurrency} & \textbf{PriceCurrency} & \textbf{ProfitCurrency} & \textbf{BaseCurrency} & \textbf{TotalRevenue} & \textbf{TotalCost} & \textbf{TotalSentMessages} & \textbf{TotalDeliveredMessages} & \textbf{TotalUndeliveredMessages} \\
\midrule
2023 & 2  & CLIENT\_031 & CLIENTCONNECTION\_065 & SUPPLIER\_001 & SUPPLIERCONNECTION\_001 & SUPPLIERCONNECTIONLINK\_005 & ROUTE\_041 & Indonesia & NETWORK\_007 & OAPROFILEDCATEGORY\_1391 & USD & USD & USD & USD & 1,401,484 & 1,245,221 & 9,766,437 & 9,079,388 & 157,978 \\
2023 & 5  & CLIENT\_017 & CLIENTCONNECTION\_001 & SUPPLIER\_001 & SUPPLIERCONNECTION\_001 & SUPPLIERCONNECTIONLINK\_005 & ROUTE\_020 & Indonesia & NETWORK\_007 & OAPROFILEDCATEGORY\_1391 & USD & USD & USD & USD & 1,239,724 & 1,215,675 & 5,959,192 & 5,390,105 & 569,087 \\
2023 & 5  & CLIENT\_017 & CLIENTCONNECTION\_001 & SUPPLIER\_001 & SUPPLIERCONNECTION\_001 & SUPPLIERCONNECTIONLINK\_001 & ROUTE\_020 & Indonesia & NETWORK\_007 & OAPROFILEDCATEGORY\_1391 & USD & USD & USD & USD & 1,236,721 & 1,212,211 & 5,942,211 & 5,377,049 & 565,162 \\
2023 & 6  & CLIENT\_017 & CLIENTCONNECTION\_001 & SUPPLIER\_001 & SUPPLIERCONNECTION\_001 & SUPPLIERCONNECTIONLINK\_001 & ROUTE\_020 & Indonesia & NETWORK\_007 & OAPROFILEDCATEGORY\_1391 & USD & USD & USD & USD & 1,193,231 & 1,215,993 & 5,960,748 & 5,187,960 & 772,788 \\
2023 & 6  & CLIENT\_017 & CLIENTCONNECTION\_001 & SUPPLIER\_001 & SUPPLIERCONNECTION\_001 & SUPPLIERCONNECTIONLINK\_005 & ROUTE\_020 & Indonesia & NETWORK\_007 & OAPROFILEDCATEGORY\_1391 & USD & USD & USD & USD & 1,192,948 & 1,215,533 & 5,958,494 & 5,186,730 & 771,764 \\
2023 & 1  & CLIENT\_030 & CLIENTCONNECTION\_071 & SUPPLIER\_001 & SUPPLIERCONNECTION\_001 & SUPPLIERCONNECTIONLINK\_001 & ROUTE\_037 & Indonesia & NETWORK\_007 & OAPROFILEDCATEGORY\_1391 & USD & USD & USD & USD & 1,098,604 & 1,085,830 & 8,516,311 & 7,900,562 & 143,330 \\
2023 & 7  & CLIENT\_017 & CLIENTCONNECTION\_001 & SUPPLIER\_001 & SUPPLIERCONNECTION\_001 & SUPPLIERCONNECTIONLINK\_005 & ROUTE\_020 & Indonesia & NETWORK\_007 & OAPROFILEDCATEGORY\_1391 & USD & USD & USD & USD & 1,065,334 & 1,089,143 & 5,338,934 & 4,631,889 & 707,045 \\
2023 & 7  & CLIENT\_017 & CLIENTCONNECTION\_001 & SUPPLIER\_001 & SUPPLIERCONNECTION\_001 & SUPPLIERCONNECTIONLINK\_001 & ROUTE\_020 & Indonesia & NETWORK\_007 & OAPROFILEDCATEGORY\_1391 & USD & USD & USD & USD & 1,065,215 & 1,089,061 & 5,338,536 & 4,631,369 & 707,167 \\
2023 & 8  & CLIENT\_017 & CLIENTCONNECTION\_001 & SUPPLIER\_001 & SUPPLIERCONNECTION\_001 & SUPPLIERCONNECTIONLINK\_005 & ROUTE\_020 & Indonesia & NETWORK\_007 & OAPROFILEDCATEGORY\_1391 & USD & USD & USD & USD & 1,009,455 & 1,032,313 & 5,060,358 & 4,388,936 & 671,415 \\
2023 & 8  & CLIENT\_017 & CLIENTCONNECTION\_001 & SUPPLIER\_001 & SUPPLIERCONNECTION\_001 & SUPPLIERCONNECTIONLINK\_001 & ROUTE\_020 & Indonesia & NETWORK\_007 & OAPROFILEDCATEGORY\_1391 & USD & USD & USD & USD & 1,008,179 & 1,032,257 & 5,060,085 & 4,383,389 & 671,671 \\
2023 & 9  & CLIENT\_017 & CLIENTCONNECTION\_001 & SUPPLIER\_001 & SUPPLIERCONNECTION\_001 & SUPPLIERCONNECTIONLINK\_005 & ROUTE\_020 & Indonesia & NETWORK\_007 & OAPROFILEDCATEGORY\_1391 & USD & USD & USD & USD & 990,972 & 1,034,318 & 5,070,188 & 4,308,572 & 761,616 \\
2023 & 9  & CLIENT\_017 & CLIENTCONNECTION\_001 & SUPPLIER\_001 & SUPPLIERCONNECTION\_001 & SUPPLIERCONNECTIONLINK\_001 & ROUTE\_020 & Indonesia & NETWORK\_007 & OAPROFILEDCATEGORY\_1391 & USD & USD & USD & USD & 985,980 & 1,029,188 & 5,045,040 & 4,286,868 & 758,172 \\
2023 & 10 & CLIENT\_017 & CLIENTCONNECTION\_001 & SUPPLIER\_001 & SUPPLIERCONNECTION\_001 & SUPPLIERCONNECTIONLINK\_001 & ROUTE\_020 & Indonesia & NETWORK\_007 & OAPROFILEDCATEGORY\_1391 & USD & USD & USD & USD & 973,111 & 1,000,161 & 4,902,752 & 4,230,918 & 671,834 \\
2023 & 10 & CLIENT\_017 & CLIENTCONNECTION\_001 & SUPPLIER\_001 & SUPPLIERCONNECTION\_001 & SUPPLIERCONNECTIONLINK\_005 & ROUTE\_020 & Indonesia & NETWORK\_007 & OAPROFILEDCATEGORY\_1391 & USD & USD & USD & USD & 972,129 & 999,797 & 4,900,966 & 4,226,649 & 674,317 \\
2024 & 1  & CLIENT\_017 & CLIENTCONNECTION\_001 & SUPPLIER\_001 & SUPPLIERCONNECTION\_001 & SUPPLIERCONNECTIONLINK\_001 & ROUTE\_020 & Indonesia & NETWORK\_007 & OAPROFILEDCATEGORY\_1391 & USD & USD & USD & USD & 961,628 & 950,977 & 4,661,653 & 4,180,991 & 480,662 \\
2024 & 1  & CLIENT\_017 & CLIENTCONNECTION\_001 & SUPPLIER\_001 & SUPPLIERCONNECTION\_001 & SUPPLIERCONNECTIONLINK\_005 & ROUTE\_020 & Indonesia & NETWORK\_007 & OAPROFILEDCATEGORY\_1391 & USD & USD & USD & USD & 961,457 & 951,019 & 4,661,860 & 4,180,249 & 481,611 \\
2023 & 3  & CLIENT\_013 & CLIENTCONNECTION\_028 & SUPPLIER\_001 & SUPPLIERCONNECTION\_004 & SUPPLIERCONNECTIONLINK\_006 & ROUTE\_016 & Indonesia & NETWORK\_007 & OAPROFILEDCATEGORY\_1192 & USD & USD & USD & USD & 919,660 & 778,519 & 8,989,829 & 7,417,443 & 937,661 \\
2023 & 3  & CLIENT\_013 & CLIENTCONNECTION\_028 & SUPPLIER\_001 & SUPPLIERCONNECTION\_004 & SUPPLIERCONNECTIONLINK\_002 & ROUTE\_016 & Indonesia & NETWORK\_007 & OAPROFILEDCATEGORY\_1192 & USD & USD & USD & USD & 917,498 & 776,689 & 8,968,699 & 7,402,547 & 932,770 \\
2024 & 3  & CLIENT\_017 & CLIENTCONNECTION\_001 & SUPPLIER\_001 & SUPPLIERCONNECTION\_001 & SUPPLIERCONNECTIONLINK\_005 & ROUTE\_020 & Indonesia & NETWORK\_007 & OAPROFILEDCATEGORY\_1391 & USD & USD & USD & USD & 916,685 & 905,109 & 4,436,810 & 3,985,587 & 451,223 \\
2024 & 3  & CLIENT\_017 & CLIENTCONNECTION\_001 & SUPPLIER\_001 & SUPPLIERCONNECTION\_001 & SUPPLIERCONNECTIONLINK\_001 & ROUTE\_020 & Indonesia & NETWORK\_007 & OAPROFILEDCATEGORY\_1391 & USD & USD & USD & USD & 916,659 & 904,948 & 4,436,022 & 3,985,474 & 450,548 \\
2023 & 12 & CLIENT\_017 & CLIENTCONNECTION\_001 & SUPPLIER\_001 & SUPPLIERCONNECTION\_001 & SUPPLIERCONNECTIONLINK\_005 & ROUTE\_020 & Indonesia & NETWORK\_007 & OAPROFILEDCATEGORY\_1391 & USD & USD & USD & USD & 912,402 & 920,357 & 4,511,554 & 3,966,966 & 544,588 \\
2023 & 12 & CLIENT\_017 & CLIENTCONNECTION\_001 & SUPPLIER\_001 & SUPPLIERCONNECTION\_001 & SUPPLIERCONNECTIONLINK\_001 & ROUTE\_020 & Indonesia & NETWORK\_007 & OAPROFILEDCATEGORY\_1391 & USD & USD & USD & USD & 911,863 & 919,831 & 4,508,976 & 3,964,621 & 544,355 \\
2023 & 4  & CLIENT\_017 & CLIENTCONNECTION\_001 & SUPPLIER\_001 & SUPPLIERCONNECTION\_001 & SUPPLIERCONNECTIONLINK\_001 & ROUTE\_020 & Indonesia & NETWORK\_007 & OAPROFILEDCATEGORY\_1391 & USD & USD & USD & USD & 905,329 & 921,216 & 4,515,767 & 3,936,215 & 579,552 \\
2023 & 4  & CLIENT\_017 & CLIENTCONNECTION\_001 & SUPPLIER\_001 & SUPPLIERCONNECTION\_001 & SUPPLIERCONNECTIONLINK\_005 & ROUTE\_020 & Indonesia & NETWORK\_007 & OAPROFILEDCATEGORY\_1391 & USD & USD & USD & USD & 904,915 & 921,114 & 4,515,266 & 3,934,414 & 580,852 \\
2023 & 4  & CLIENT\_013 & CLIENTCONNECTION\_028 & SUPPLIER\_001 & SUPPLIERCONNECTION\_004 & SUPPLIERCONNECTIONLINK\_002 & ROUTE\_016 & Indonesia & NETWORK\_007 & OAPROFILEDCATEGORY\_1192 & USD & USD & USD & USD & 892,816 & 755,795 & 8,727,427 & 7,091,224 & 986,710 \\
2023 & 4  & CLIENT\_013 & CLIENTCONNECTION\_028 & SUPPLIER\_001 & SUPPLIERCONNECTION\_004 & SUPPLIERCONNECTIONLINK\_006 & ROUTE\_016 & Indonesia & NETWORK\_007 & OAPROFILEDCATEGORY\_1192 & USD & USD & USD & USD & 892,773 & 755,759 & 8,727,006 & 7,089,587 & 988,001 \\
2024 & 2  & CLIENT\_017 & CLIENTCONNECTION\_001 & SUPPLIER\_001 & SUPPLIERCONNECTION\_001 & SUPPLIERCONNECTIONLINK\_005 & ROUTE\_020 & Indonesia & NETWORK\_007 & OAPROFILEDCATEGORY\_1391 & USD & USD & USD & USD & 888,989 & 877,699 & 4,302,446 & 3,865,168 & 437,278 \\
2024 & 2  & CLIENT\_017 & CLIENTCONNECTION\_001 & SUPPLIER\_001 & SUPPLIERCONNECTION\_001 & SUPPLIERCONNECTIONLINK\_001 & ROUTE\_020 & Indonesia & NETWORK\_007 & OAPROFILEDCATEGORY\_1391 & USD & USD & USD & USD & 887,783 & 876,306 & 4,295,618 & 3,859,925 & 435,693 \\
2023 & 5  & CLIENT\_013 & CLIENTCONNECTION\_028 & SUPPLIER\_001 & SUPPLIERCONNECTION\_004 & SUPPLIERCONNECTIONLINK\_002 & ROUTE\_016 & Indonesia & NETWORK\_007 & OAPROFILEDCATEGORY\_1192 & USD & USD & USD & USD & 881,111 & 745,886 & 8,613,008 & 7,069,540 & 920,428 \\
2023 & 5  & CLIENT\_013 & CLIENTCONNECTION\_028 & SUPPLIER\_001 & SUPPLIERCONNECTION\_004 & SUPPLIERCONNECTIONLINK\_006 & ROUTE\_016 & Indonesia & NETWORK\_007 & OAPROFILEDCATEGORY\_1192 & USD & USD & USD & USD & 871,363 & 737,635 & 8,517,720 & 6,992,868 & 909,522 \\
2023 & 11 & CLIENT\_017 & CLIENTCONNECTION\_001 & SUPPLIER\_001 & SUPPLIERCONNECTION\_001 & SUPPLIERCONNECTIONLINK\_001 & ROUTE\_020 & Indonesia & NETWORK\_007 & OAPROFILEDCATEGORY\_1391 & USD & USD & USD & USD & 860,531 & 881,575 & 4,321,448 & 3,741,441 & 580,007 \\
2023 & 11 & CLIENT\_017 & CLIENTCONNECTION\_001 & SUPPLIER\_001 & SUPPLIERCONNECTION\_001 & SUPPLIERCONNECTIONLINK\_005 & ROUTE\_020 & Indonesia & NETWORK\_007 & OAPROFILEDCATEGORY\_1391 & USD & USD & USD & USD & 859,805 & 880,930 & 4,318,286 & 3,738,284 & 580,002 \\
2023 & 3  & CLIENT\_013 & CLIENTCONNECTION\_028 & SUPPLIER\_001 & SUPPLIERCONNECTION\_004 & SUPPLIERCONNECTIONLINK\_006 & ROUTE\_016 & Indonesia & NETWORK\_007 & OAPROFILEDCATEGORY\_3758 & USD & USD & USD & USD & 831,216 & 703,649 & 8,125,275 & 6,727,354 & 1,031,719 \\
2023 & 3  & CLIENT\_013 & CLIENTCONNECTION\_028 & SUPPLIER\_001 & SUPPLIERCONNECTION\_004 & SUPPLIERCONNECTIONLINK\_002 & ROUTE\_016 & Indonesia & NETWORK\_007 & OAPROFILEDCATEGORY\_3758 & USD & USD & USD & USD & 829,778 & 702,432 & 8,111,224 & 6,716,929 & 1,028,377 \\
2023 & 1  & CLIENT\_013 & CLIENTCONNECTION\_028 & SUPPLIER\_001 & SUPPLIERCONNECTION\_004 & SUPPLIERCONNECTIONLINK\_002 & ROUTE\_016 & Indonesia & NETWORK\_007 & OAPROFILEDCATEGORY\_1192 & USD & USD & USD & USD & 810,829 & 686,391 & 7,925,994 & 6,549,499 & 819,832 \\
2023 & 1  & CLIENT\_013 & CLIENTCONNECTION\_028 & SUPPLIER\_001 & SUPPLIERCONNECTION\_004 & SUPPLIERCONNECTIONLINK\_006 & ROUTE\_016 & Indonesia & NETWORK\_007 & OAPROFILEDCATEGORY\_1192 & USD & USD & USD & USD & 810,779 & 686,348 & 7,925,501 & 6,548,783 & 820,443 \\
2023 & 6  & CLIENT\_013 & CLIENTCONNECTION\_028 & SUPPLIER\_001 & SUPPLIERCONNECTION\_004 & SUPPLIERCONNECTIONLINK\_006 & ROUTE\_016 & Indonesia & NETWORK\_007 & OAPROFILEDCATEGORY\_1192 & USD & USD & USD & USD & 809,847 & 685,559 & 7,916,391 & 6,536,815 & 830,352 \\
2023 & 6  & CLIENT\_013 & CLIENTCONNECTION\_028 & SUPPLIER\_001 & SUPPLIERCONNECTION\_004 & SUPPLIERCONNECTIONLINK\_002 & ROUTE\_016 & Indonesia & NETWORK\_007 & OAPROFILEDCATEGORY\_1192 & USD & USD & USD & USD & 805,776 & 682,113 & 7,876,599 & 6,501,031 & 827,723 \\
2024 & 3  & CLIENT\_013 & CLIENTCONNECTION\_028 & SUPPLIER\_001 & SUPPLIERCONNECTION\_004 & SUPPLIERCONNECTIONLINK\_002 & ROUTE\_016 & Indonesia & NETWORK\_007 & OAPROFILEDCATEGORY\_3758 & USD & USD & USD & USD & 792,204 & 670,625 & 7,743,932 & 6,437,084 & 917,715 \\
2024 & 3  & CLIENT\_013 & CLIENTCONNECTION\_028 & SUPPLIER\_001 & SUPPLIERCONNECTION\_004 & SUPPLIERCONNECTIONLINK\_006 & ROUTE\_016 & Indonesia & NETWORK\_007 & OAPROFILEDCATEGORY\_3758 & USD & USD & USD & USD & 791,645 & 670,152 & 7,738,470 & 6,432,588 & 917,044 \\
2023 & 3  & CLIENT\_017 & CLIENTCONNECTION\_001 & SUPPLIER\_001 & SUPPLIERCONNECTION\_001 & SUPPLIERCONNECTIONLINK\_001 & ROUTE\_020 & Indonesia & NETWORK\_007 & OAPROFILEDCATEGORY\_1391 & USD & USD & USD & USD & 774,855 & 830,830 & 4,072,695 & 3,443,902 & 628,793 \\
2023 & 2  & CLIENT\_013 & CLIENTCONNECTION\_028 & SUPPLIER\_001 & SUPPLIERCONNECTION\_004 & SUPPLIERCONNECTIONLINK\_006 & ROUTE\_016 & Indonesia & NETWORK\_007 & OAPROFILEDCATEGORY\_1192 & USD & USD & USD & USD & 755,581 & 639,622 & 7,385,932 & 6,143,402 & 738,106 \\
2023 & 2  & CLIENT\_013 & CLIENTCONNECTION\_028 & SUPPLIER\_001 & SUPPLIERCONNECTION\_004 & SUPPLIERCONNECTIONLINK\_002 & ROUTE\_016 & Indonesia & NETWORK\_007 & OAPROFILEDCATEGORY\_1192 & USD & USD & USD & USD & 755,534 & 639,582 & 7,385,472 & 6,143,525 & 738,090 \\
2023 & 1  & CLIENT\_013 & CLIENTCONNECTION\_028 & SUPPLIER\_001 & SUPPLIERCONNECTION\_004 & SUPPLIERCONNECTIONLINK\_006 & ROUTE\_016 & Indonesia & NETWORK\_007 & OAPROFILEDCATEGORY\_3758 & USD & USD & USD & USD & 753,265 & 637,661 & 7,363,296 & 6,180,515 & 867,061 \\
2023 & 1  & CLIENT\_013 & CLIENTCONNECTION\_028 & SUPPLIER\_001 & SUPPLIERCONNECTION\_004 & SUPPLIERCONNECTIONLINK\_002 & ROUTE\_016 & Indonesia & NETWORK\_007 & OAPROFILEDCATEGORY\_3758 & USD & USD & USD & USD & 752,645 & 637,136 & 7,357,229 & 6,174,895 & 866,668 \\
2023 & 3  & CLIENT\_017 & CLIENTCONNECTION\_001 & SUPPLIER\_001 & SUPPLIERCONNECTION\_001 & SUPPLIERCONNECTIONLINK\_005 & ROUTE\_020 & Indonesia & NETWORK\_007 & OAPROFILEDCATEGORY\_1391 & USD & USD & USD & USD & 750,501 & 805,851 & 3,950,251 & 3,338,391 & 611,860 \\
2023 & 7  & CLIENT\_013 & CLIENTCONNECTION\_028 & SUPPLIER\_001 & SUPPLIERCONNECTION\_004 & SUPPLIERCONNECTIONLINK\_006 & ROUTE\_016 & Indonesia & NETWORK\_007 & OAPROFILEDCATEGORY\_1192 & USD & USD & USD & USD & 740,140 & 626,550 & 7,234,994 & 6,016,647 & 760,333 \\
2023 & 7  & CLIENT\_013 & CLIENTCONNECTION\_028 & SUPPLIER\_001 & SUPPLIERCONNECTION\_004 & SUPPLIERCONNECTIONLINK\_002 & ROUTE\_016 & Indonesia & NETWORK\_007 & OAPROFILEDCATEGORY\_1192 & USD & USD & USD & USD & 737,533 & 624,344 & 7,209,515 & 5,996,191 & 756,719 \\
2023 & 4  & CLIENT\_013 & CLIENTCONNECTION\_028 & SUPPLIER\_001 & SUPPLIERCONNECTION\_004 & SUPPLIERCONNECTIONLINK\_006 & ROUTE\_016 & Indonesia & NETWORK\_007 & OAPROFILEDCATEGORY\_3758 & USD & USD & USD & USD & 722,776 & 611,852 & 7,065,262 & 5,756,780 & 957,487 \\
2023 & 4  & CLIENT\_013 & CLIENTCONNECTION\_028 & SUPPLIER\_001 & SUPPLIERCONNECTION\_004 & SUPPLIERCONNECTIONLINK\_002 & ROUTE\_016 & Indonesia & NETWORK\_007 & OAPROFILEDCATEGORY\_3758 & USD & USD & USD & USD & 722,650 & 611,745 & 7,064,031 & 5,754,745 & 958,143 \\
\bottomrule
\end{tabular}
\end{adjustbox}

\vspace{4pt}
{\small\textit{Catatan: Nama klien, supplier, dan koneksi telah disamarkan (\textit{masked}) untuk menjaga kerahasiaan data.}}
\end{center}
\end{landscape}

% ============================================================
% lampiran_kode.tex
% File ini berisi KONTEN LAMPIRAN saja (tanpa preamble).
%
% CARA PAKAI:
%   Di bab Lampiran dokumen utama:
%       \section{Kode Program}
%       \input{lampiran_kode}
%
%   Pastikan preamble dokumen utama sudah menyertakan
%   definisi dari lampiran_setup.tex (lihat file tersebut).
% ============================================================

\section*{Kode Program}

\begin{lstlisting}[style=pycode,caption={Import pustaka dan load dataset}]
import pandas as pd
import numpy as np
import matplotlib.pyplot as plt
from pathlib import Path

df = pd.read_csv("data/SMS International.csv")
df.info()
\end{lstlisting}

\begin{lstlisting}[style=pycode,caption={Normalisasi kolom Network}]
# Normalisasi nilai kategorikal sebelum masking
if 'Network' not in df.columns:
    raise ValueError("Kolom 'Network' tidak ditemukan di dataframe.")

# Simpan snapshot sebelum cleaning untuk verifikasi
network_before = df['Network'].astype(str).value_counts(dropna=False)

# Rapikan spasi dan normalisasi typo umum
df['Network'] = df['Network'].astype(str).str.strip()
df['Network'] = df['Network'].replace({
    r'(?i)^telkomcel$': 'Telkomsel',
    r'(?i)^telkomsel$': 'Telkomsel'
}, regex=True)

network_after = df['Network'].astype(str).value_counts(dropna=False)

print("=== Sebelum cleaning (Network) ===")
print(
    network_before.loc[
        [idx for idx in network_before.index if 'telkom' in idx.lower()]
    ] if any('telkom' in idx.lower() for idx in network_before.index)
    else 'Tidak ada nilai Telkom* sebelum cleaning'
)

print("\n=== Sesudah cleaning (Network) ===")
print(
    network_after.loc[
        [idx for idx in network_after.index if 'telkom' in idx.lower()]
    ] if any('telkom' in idx.lower() for idx in network_after.index)
    else 'Tidak ada nilai Telkom* sesudah cleaning'
)
\end{lstlisting}

\begin{lstlisting}[style=pycode,caption={Identifikasi dan masking kolom sensitif}]
# Identifikasi kolom sensitif berdasarkan nama kolom (case-insensitive)
sensitive_keywords = [
    "client", "company", "provider", "supplier",
    "partner", "network", "route", "connection", "oaprofiledcategory"
]

sensitive_columns = [
    col for col in df.columns
    if any(keyword in col.lower() for keyword in sensitive_keywords)
]

def mask_column(series, label):
    """Masking deterministik: nilai unik -> token anonim per kolom."""
    clean_values = sorted(series.dropna().astype(str).unique())
    mapping = {
        value: f"{label}_{idx:03d}"
        for idx, value in enumerate(clean_values, start=1)
    }
    masked = series.astype(str).map(mapping)
    masked = masked.where(series.notna(), np.nan)
    return masked, mapping

masking_maps = {}

for col in sensitive_columns:
    masked_series, mapping = mask_column(df[col], col.upper())
    df[col] = masked_series
    masking_maps[col] = mapping

print("Kolom yang dimasking:", sensitive_columns)
print("Contoh data setelah masking:")
display(df.head(3))
\end{lstlisting}

\begin{lstlisting}[style=pycode,caption={Pembersihan kolom tidak diperlukan}]
df = df.drop(columns=['Unnamed: 21', 'Unnamed: 22'], errors='ignore')
df.to_csv("data/SMS_International_masked.csv", index=False)
\end{lstlisting}

\begin{lstlisting}[style=pycode,caption={Agregasi temporal dan perhitungan MoM growth}]
# Pastikan kolom numerik temporal siap dihitung
temporal_cols = ['TotalRevenue', 'TotalCost', 'TotalSentMessages']
for col in temporal_cols:
    df_terurut[col] = pd.to_numeric(
        df_terurut[col].astype(str)
            .str.replace(',', '', regex=False)
            .str.replace(' ', '', regex=False),
        errors='coerce'
    )

# Agregasi bulanan
temporal_summary = (
    df_terurut
    .groupby(['Year', 'Month'], as_index=False)[temporal_cols]
    .sum()
    .sort_values(['Year', 'Month'])
    .reset_index(drop=True)
)

# Hitung pertumbuhan month-to-month (%)
for col in temporal_cols:
    temporal_summary[f'{col}_MoM_pct'] = (
        temporal_summary[col].pct_change() * 100
    )

print('Ringkasan temporal (5 baris pertama):')
display(temporal_summary.head())
\end{lstlisting}

\begin{lstlisting}[style=pycode,caption={Visualisasi tren temporal dan pertumbuhan MoM}]
# Tren individual -- masing-masing figure terpisah
plot_configs = [
    ('TotalRevenue',      'Tren Total Revenue Bulanan',       'tren_revenue.png'),
    ('TotalCost',         'Tren Total Cost Bulanan',          'tren_cost.png'),
    ('TotalSentMessages', 'Tren Total Sent Messages Bulanan', 'tren_sent_messages.png'),
]

for col, title, fname in plot_configs:
    fig, ax = plt.subplots(figsize=(12, 4))
    ax.plot(temporal_summary.index, temporal_summary[col],
            marker='o', linewidth=2)
    ax.set_title(title)
    ax.set_ylabel(col)
    ax.set_xlabel('Periode (urut Year-Month)')
    ax.grid(alpha=0.3)
    plt.tight_layout()
    plt.savefig(fname, dpi=150, bbox_inches='tight')
    plt.show()

# Pertumbuhan MoM Revenue vs Cost
fig, ax = plt.subplots(figsize=(12, 5))
ax.plot(temporal_summary.index,
        temporal_summary['TotalRevenue_MoM_pct'],
        marker='o', label='Revenue MoM %')
ax.plot(temporal_summary.index,
        temporal_summary['TotalCost_MoM_pct'],
        marker='o', label='Cost MoM %')
ax.axhline(0, color='gray', linestyle='--', linewidth=1)
ax.set_title('Pertumbuhan Bulanan Revenue dan Cost (MoM %)')
ax.set_xlabel('Periode (urut Year-Month)')
ax.set_ylabel('Growth (%)')
ax.legend()
ax.grid(alpha=0.3)
plt.tight_layout()
plt.savefig('tren_mom_growth.png', dpi=150, bbox_inches='tight')
plt.show()
\end{lstlisting}

\begin{lstlisting}[style=pycode,caption={Matriks korelasi Pearson dan heatmap}]
import seaborn as sns

# Pilih kolom numerik
numeric_cols = [
    'TotalRevenue', 'TotalCost', 'TotalSentMessages',
    'TotalDeliveredMessages', 'TotalUndeliveredMessages'
]
df_numeric = df_terurut[numeric_cols].copy()

for col in numeric_cols:
    df_numeric[col] = pd.to_numeric(
        df_numeric[col].astype(str).str.replace(',', ''),
        errors='coerce'
    )

# Hitung korelasi Pearson
correlation_matrix = df_numeric.corr(method='pearson')
print("=== Matriks Korelasi Pearson ===")
print(correlation_matrix)

# Visualisasi heatmap
plt.figure(figsize=(10, 8))
sns.heatmap(
    correlation_matrix, annot=True, fmt='.3f',
    cmap='coolwarm', center=0, square=True,
    linewidths=1, cbar_kws={"shrink": 0.8}
)
plt.title('Correlation Matrix (Pearson) - Fitur Numerik',
          fontsize=14, fontweight='bold')
plt.tight_layout()
plt.show()

print("\n=== Korelasi dengan TotalRevenue (Descending) ===")
revenue_corr = correlation_matrix['TotalRevenue'].sort_values(ascending=False)
print(revenue_corr)
\end{lstlisting}

\begin{lstlisting}[style=pycode,caption={Bubble chart top-20 Series\_ID berdasarkan transaksi}]
# Bentuk Series_ID dari 3 fitur kategorikal
series_df = df.copy()
for c in ['Client', 'Supplier', 'Route']:
    series_df[c] = series_df[c].fillna('UNKNOWN').astype(str)
series_df['Series_ID'] = (
    series_df['Client'] + ' | ' +
    series_df['Supplier'] + ' | ' +
    series_df['Route']
)

series_df['MonthPeriod'] = pd.to_datetime(
    dict(year=series_df['Year'],
         month=series_df['Month'], day=1),
    errors='coerce'
)
series_df = series_df.dropna(subset=['MonthPeriod'])

# Agregasi jumlah transaksi per Series_ID per bulan
bubble_data = (
    series_df
    .groupby(['Series_ID', 'MonthPeriod'], as_index=False)
    .size()
    .rename(columns={'size': 'JumlahTransaksi'})
)

# Ambil top-20 series
top_n = 20
top_series = (
    bubble_data.groupby('Series_ID', as_index=False)['JumlahTransaksi']
    .sum()
    .sort_values('JumlahTransaksi', ascending=False)
    .head(top_n)['Series_ID']
)
bubble_plot = bubble_data[bubble_data['Series_ID'].isin(top_series)].copy()

series_order = (
    bubble_plot.groupby('Series_ID')['JumlahTransaksi']
    .sum().sort_values(ascending=False).index.tolist()
)
bubble_plot['x_pos'] = bubble_plot['Series_ID'].map(
    {s: i for i, s in enumerate(series_order)}
)

# Bubble chart
plt.figure(figsize=(16, 9))
sns.scatterplot(
    data=bubble_plot, x='x_pos', y='MonthPeriod',
    size='JumlahTransaksi', hue='JumlahTransaksi',
    sizes=(40, 1800), palette='viridis', alpha=0.75,
    edgecolor='black', linewidth=0.4
)
plt.title(
    'Bubble Chart Series_ID vs Bulan (Ukuran = Banyaknya Transaksi)',
    fontsize=14, fontweight='bold'
)
plt.tight_layout()
plt.show()
\end{lstlisting}

\begin{lstlisting}[style=pycode,caption={Penandaan pola anomali pada data transaksi}]
# Konversi kolom ke numerik
for col in ['TotalRevenue', 'TotalCost', 'TotalProfit']:
    df_terurut[col] = (
        df_terurut[col]
        .astype(str)
        .str.replace(',', '', regex=False)
        .str.replace(' ', '', regex=False)
        .str.replace('-', '', regex=False)
        .str.strip()
    )
    df_terurut[col] = pd.to_numeric(df_terurut[col],
                                    errors='coerce').fillna(0)

# Pola 1: Transaksi gagal (0,0,0) -> hapus
mask_failed = ((df_terurut['TotalRevenue'] == 0) &
               (df_terurut['TotalCost']    == 0) &
               (df_terurut['TotalProfit']  == 0))
count_failed = mask_failed.sum()
df_terurut   = df_terurut[~mask_failed].copy()
print(f'Transaksi gagal dihapus : {count_failed}')
print(f'Sisa data               : {len(df_terurut)}')

# Pola 2: Bonus (rev!=0, cost=0, profit=0)
mask_rev_only = ((df_terurut['TotalRevenue'] != 0) &
                 (df_terurut['TotalCost']    == 0) &
                 (df_terurut['TotalProfit']  == 0))
df_terurut['is_anomaly_rev_only'] = mask_rev_only.astype(int)
print(f'Transaksi bonus         : {mask_rev_only.sum()}')

# Pola 3: Testing (rev=0, cost!=0, profit!=0)
mask_cost_only = ((df_terurut['TotalRevenue'] == 0) &
                  (df_terurut['TotalCost']    != 0) &
                  (df_terurut['TotalProfit']  != 0))
df_terurut['is_anomaly_cost_only'] = mask_cost_only.astype(int)
print(f'Transaksi testing       : {mask_cost_only.sum()}')
\end{lstlisting}

\begin{lstlisting}[style=pycode,caption={Agregasi bulanan per Series\_ID}]
# Perbaiki Series_ID -- masukkan Network
df['Series_ID'] = (df['Client'] + '_' + df['Route'] + '_' +
                   df['SupplierConnection'] + '_' + df['Network'])
print(f"Jumlah Series_ID unik: {df['Series_ID'].nunique()}")

group_cols = ['Series_ID', 'Year', 'Month',
              'Client', 'Route', 'SupplierConnection', 'Network']

# Aggregasi anomali
anom_rev = df.groupby(group_cols).agg(
    total_rows=('TotalRevenue', 'count'),
    anomaly_rev_rows=('is_anomaly_rev_only', 'sum')
).reset_index()
anom_rev['anomaly_revenue_pct'] = np.where(
    anom_rev['total_rows'] > 0,
    anom_rev['anomaly_rev_rows'] / anom_rev['total_rows'], 0
)

agg_dict = {
    'TotalRevenue'          : 'sum',
    'TotalCost'             : 'sum',
    'TotalProfit'           : 'sum',
    'TotalSentMessages'     : 'sum',
    'TotalDeliveredMessages': 'sum',
    'TotalUndeliveredMessages': 'sum',
    'is_anomaly_rev_only'   : 'sum',
    'is_anomaly_cost_only'  : 'sum',
    'OAProfiledCategory'    : 'nunique',
}

monthly = df.groupby(group_cols, as_index=False).agg(agg_dict)
monthly = monthly.rename(columns={'OAProfiledCategory': 'unique_OA_count'})
\end{lstlisting}

\begin{lstlisting}[style=pycode,caption={Pembuatan complete panel (cross join series x bulan)}]
# Cross join: setiap series x setiap bulan
all_series = df[['Series_ID', 'Client', 'Route',
                 'SupplierConnection', 'Network']].drop_duplicates()
all_months = sorted(df[['Year', 'Month']].drop_duplicates().values.tolist())
month_df   = pd.DataFrame(all_months, columns=['Year', 'Month'])
panel      = all_series.merge(month_df, how='cross')
panel      = panel.merge(monthly, on=group_cols, how='left')

# Fill bulan tanpa transaksi dengan 0
fill_zero_cols = [
    'TotalRevenue', 'TotalCost', 'TotalProfit',
    'TotalSentMessages', 'TotalDeliveredMessages',
    'TotalUndeliveredMessages', 'is_anomaly_rev_only',
    'is_anomaly_cost_only', 'unique_OA_count',
    'max_OA_revenue', 'min_OA_revenue'
]
for col in fill_zero_cols:
    if col in panel.columns:
        panel[col] = panel[col].fillna(0)

panel['is_active'] = (panel['TotalRevenue'] > 0).astype(int)
panel['rateCost']      = np.where(panel['TotalDeliveredMessages'] > 0,
    panel['TotalCost'] / panel['TotalDeliveredMessages'], 0)
panel['profit_margin'] = np.where(panel['TotalRevenue'] > 0,
    panel['TotalProfit'] / panel['TotalRevenue'], 0)
panel['delivery_rate'] = np.where(panel['TotalSentMessages'] > 0,
    panel['TotalDeliveredMessages'] / panel['TotalSentMessages'], 0)

panel   = panel.sort_values(['Series_ID', 'Year', 'Month']).reset_index(drop=True)
monthly = panel.copy()
print(f'Complete panel shape : {monthly.shape}')
print(f'Active months        : {monthly["is_active"].sum()}')
print(f'is_active ratio      : {monthly["is_active"].mean()*100:.1f}%')
\end{lstlisting}

\begin{lstlisting}[style=pycode,caption={Lag features, rolling features, dan streak aktif}]
monthly = monthly.sort_values(['Series_ID', 'Year', 'Month']).reset_index(drop=True)

# Lag features
for lag in [1, 2, 3]:
    monthly[f'rev_lag{lag}']    = monthly.groupby('Series_ID')['TotalRevenue'].shift(lag)
    monthly[f'active_lag{lag}'] = monthly.groupby('Series_ID')['is_active'].shift(lag)

# Rolling features (shift(1) agar tidak bocor)
grp_rev = monthly.groupby('Series_ID')['TotalRevenue']
monthly['rolling_mean_3'] = grp_rev.shift(1).rolling(3).mean().reset_index(level=0, drop=True)
monthly['rolling_mean_6'] = grp_rev.shift(1).rolling(6).mean().reset_index(level=0, drop=True)
monthly['rolling_std_3']  = grp_rev.shift(1).rolling(3).std() .reset_index(level=0, drop=True)
monthly['rolling_std_6']  = grp_rev.shift(1).rolling(6).std() .reset_index(level=0, drop=True)

grp_act = monthly.groupby('Series_ID')['is_active']
monthly['pct_active_3m'] = grp_act.shift(1).rolling(3).mean().reset_index(level=0, drop=True)
monthly['pct_active_6m'] = grp_act.shift(1).rolling(6).mean().reset_index(level=0, drop=True)

# Streak aktif berturut-turut
def active_streak(s):
    streak, count = [], 0
    for val in s.shift(1).fillna(0):
        count = count + 1 if val == 1 else 0
        streak.append(count)
    return pd.Series(streak, index=s.index)

monthly['active_streak'] = monthly.groupby(
    'Series_ID', group_keys=False)['is_active'].apply(active_streak)

# Bulan sejak terakhir aktif
def months_since_active(s):
    result, count = [], 0
    for val in s.shift(1).fillna(0):
        count = 0 if val == 1 else count + 1
        result.append(count)
    return pd.Series(result, index=s.index)

monthly['months_since_active'] = monthly.groupby(
    'Series_ID', group_keys=False)['is_active'].apply(months_since_active)

# Siklikalitas bulan
monthly['month_sin'] = np.sin(2 * np.pi * monthly['Month'] / 12)
monthly['month_cos'] = np.cos(2 * np.pi * monthly['Month'] / 12)

# Target: bulan DEPAN
monthly['target_active']  = monthly.groupby('Series_ID')['is_active'].shift(-1)
monthly['target_revenue'] = monthly.groupby('Series_ID')['TotalRevenue'].shift(-1)
monthly = monthly[monthly['target_active'].notna()].reset_index(drop=True)

print(f"Shape setelah feature engineering : {monthly.shape}")
print(f"Baris Stage 1 (semua)             : {len(monthly)}")
print(f"Baris Stage 2 (aktif)             : {int(monthly['target_active'].sum())}")
\end{lstlisting}

\begin{lstlisting}[style=pycode,caption={Definisi daftar fitur Stage 1 dan Stage 2}]
FEATURES_STAGE1 = [
    # Status historis
    'active_lag1', 'active_lag2', 'active_lag3',
    'active_streak', 'months_since_active',
    'pct_active_3m', 'pct_active_6m',
    # Revenue historis
    'rev_lag1', 'rev_lag2', 'rev_lag3',
    'rolling_mean_3', 'rolling_mean_6',
    # Operasional
    'TotalSentMessages', 'TotalDeliveredMessages',
    'delivery_rate', 'unique_OA_count',
    # Anomali
    'anomaly_count', 'anomaly_revenue_pct',
    # Siklikalitas
    'month_sin', 'month_cos',
]

FEATURES_STAGE2 = [
    # Revenue historis
    'rev_lag1', 'rev_lag2', 'rev_lag3',
    'rolling_mean_3', 'rolling_mean_6',
    'rolling_std_3', 'rolling_std_6',
    # Status historis
    'active_lag1', 'active_lag2',
    'active_streak', 'pct_active_3m', 'pct_active_6m',
    # Operasional & margin
    'TotalSentMessages', 'TotalDeliveredMessages',
    'delivery_rate', 'profit_margin', 'rateCost',
    'unique_OA_count', 'max_OA_revenue',
    # Anomali
    'anomaly_count', 'anomaly_revenue_pct',
    # Siklikalitas
    'month_sin', 'month_cos',
]
\end{lstlisting}

\begin{lstlisting}[style=pycode,caption={Time-based split dan pelatihan XGBoost Classifier}]
from sklearn.metrics import (classification_report, roc_auc_score,
                             average_precision_score)
from xgboost import XGBClassifier
import warnings
warnings.filterwarnings('ignore')

# Time-based split
monthly['period'] = monthly['Year'] * 12 + monthly['Month']
cutoff = monthly['period'].quantile(0.80).astype(int)

train = monthly[monthly['period'] <= cutoff].copy()
test  = monthly[monthly['period'] >  cutoff].copy()

print(f"Train : {len(train)} baris | period {train['period'].min()}-{train['period'].max()}")
print(f"Test  : {len(test)} baris  | period {test['period'].min()}-{test['period'].max()}")

X_train_s1 = train[FEATURES_STAGE1]
y_train_s1 = train['target_active'].astype(int)
X_test_s1  = test[FEATURES_STAGE1]
y_test_s1  = test['target_active'].astype(int)

clf = XGBClassifier(
    n_estimators=1000,
    learning_rate=0.3,
    max_depth=6,
    subsample=1.0,
    colsample_bytree=1.0,
    reg_lambda=1.0,
    scale_pos_weight=(y_train_s1 == 0).sum() / (y_train_s1 == 1).sum(),
    random_state=42,
    eval_metric='logloss',
    early_stopping_rounds=50,
    verbosity=0,
)
clf.fit(X_train_s1, y_train_s1,
        eval_set=[(X_test_s1, y_test_s1)],
        verbose=False)

# Evaluasi
y_pred_s1 = clf.predict(X_test_s1)
y_prob_s1 = clf.predict_proba(X_test_s1)[:, 1]

print("\n-- Stage 1: Classification Report --")
print(classification_report(y_test_s1, y_pred_s1,
                             target_names=['Idle', 'Active']))
print(f"ROC-AUC           : {roc_auc_score(y_test_s1, y_prob_s1):.4f}")
print(f"Average Precision : {average_precision_score(y_test_s1, y_prob_s1):.4f}")
\end{lstlisting}

\begin{lstlisting}[style=pycode,caption={Definisi hyperparameter space dan fungsi fitness GA}]
import pygad
from sklearn.ensemble import RandomForestRegressor
from xgboost import XGBRegressor
from sklearn.metrics import mean_absolute_error, mean_squared_error, r2_score

def wmape(actual, pred):
    return np.sum(np.abs(actual - pred)) / np.sum(np.abs(actual))

N_FEATURES = len(FEATURES_STAGE2)

# Hyperparameter space Random Forest
RF_GENE_SPACE = [
    {'low': 50,  'high': 1000},        # n_estimators
    {'low': 1,   'high': 7},           # max_depth
    {'low': 2,   'high': 20},          # min_samples_split
    {'low': 1,   'high': 50},          # min_samples_leaf
    {'low': 1,   'high': N_FEATURES},  # max_features
    {'low': 0.1, 'high': 1.0},         # max_samples
]

def decode_rf(solution):
    return {
        'n_estimators'     : int(round(solution[0])),
        'max_depth'        : int(round(solution[1])),
        'min_samples_split': int(round(solution[2])),
        'min_samples_leaf' : int(round(solution[3])),
        'max_features'     : int(round(solution[4])),
        'max_samples'      : float(np.clip(solution[5], 0.1, 1.0)),
    }

# Hyperparameter space XGBoost
XGB_GENE_SPACE = [
    {'low': 0.001, 'high': 1.0},   # learning_rate
    {'low': 50,    'high': 1000},  # n_estimators
    {'low': 1,     'high': 7},     # max_depth
    {'low': 0.1,   'high': 1.0},   # subsample
    {'low': 0.1,   'high': 1.0},   # colsample_bytree
    {'low': 1,     'high': 3},     # reg_lambda
]

def decode_xgb(solution):
    return {
        'learning_rate'   : float(np.clip(solution[0], 0.001, 1.0)),
        'n_estimators'    : int(round(solution[1])),
        'max_depth'       : int(round(solution[2])),
        'subsample'       : float(np.clip(solution[3], 0.1, 1.0)),
        'colsample_bytree': float(np.clip(solution[4], 0.1, 1.0)),
        'reg_lambda'      : float(np.clip(solution[5], 1, 3)),
    }

# Fitness functions
def fitness_rf(ga_instance, solution, solution_idx):
    params = decode_rf(solution)
    model  = RandomForestRegressor(**params, random_state=42, n_jobs=1)
    model.fit(X_train_s2, y_train_s2)
    return -wmape(y_test_s2, model.predict(X_test_s2))

def fitness_xgb(ga_instance, solution, solution_idx):
    params = decode_xgb(solution)
    model  = XGBRegressor(**params, random_state=42, verbosity=0)
    model.fit(X_train_s2, y_train_s2)
    return -wmape(y_test_s2, model.predict(X_test_s2))
\end{lstlisting}

\begin{lstlisting}[style=pycode,caption={Eksekusi GA untuk Random Forest dan XGBoost}]
def on_generation(ga):
    if ga.generations_completed % 10 == 0:
        print(f"  Gen {ga.generations_completed:>3} | "
              f"Best WMAPE: {-ga.best_solution()[1]:.4f}")

# GA -- Random Forest
ga_rf = pygad.GA(
    num_generations=50, num_parents_mating=4,
    fitness_func=fitness_rf, sol_per_pop=40, num_genes=6,
    gene_space=RF_GENE_SPACE, mutation_percent_genes=20,
    crossover_type='scattered', parent_selection_type='tournament',
    K_tournament=3, keep_elitism=2,
    on_generation=on_generation, random_seed=42,
)
ga_rf.run()
best_sol_rf, best_fit_rf, _ = ga_rf.best_solution()
best_params_rf = decode_rf(best_sol_rf)
print(f"Best params RF : {best_params_rf}")
print(f"Best WMAPE RF  : {-best_fit_rf:.4f} ({-best_fit_rf*100:.2f}%)")

# GA -- XGBoost
ga_xgb = pygad.GA(
    num_generations=50, num_parents_mating=4,
    fitness_func=fitness_xgb, sol_per_pop=40, num_genes=6,
    gene_space=XGB_GENE_SPACE, mutation_percent_genes=20,
    crossover_type='scattered', parent_selection_type='tournament',
    K_tournament=3, keep_elitism=2,
    on_generation=on_generation, random_seed=42,
)
ga_xgb.run()
best_sol_xgb, best_fit_xgb, _ = ga_xgb.best_solution()
best_params_xgb = decode_xgb(best_sol_xgb)
print(f"Best params XGB : {best_params_xgb}")
print(f"Best WMAPE XGB  : {-best_fit_xgb:.4f} ({-best_fit_xgb*100:.2f}%)")
\end{lstlisting}

\begin{lstlisting}[style=pycode,caption={Fungsi evaluasi dan perbandingan 4 model}]
def evaluate(name, model, X_tr, y_tr, X_te, y_te):
    model.fit(X_tr, y_tr)
    pred      = model.predict(X_te)
    mae       = mean_absolute_error(y_te, pred)
    rmse      = np.sqrt(mean_squared_error(y_te, pred))
    r2        = r2_score(y_te, pred)
    wmape_val = wmape(y_te, pred)
    mape      = np.mean(np.abs((y_te - pred) /
                np.where(y_te == 0, 1, y_te))) * 100
    print(f"\n{name}")
    print(f"  MAE   : {mae:>14.2f}")
    print(f"  RMSE  : {rmse:>14.2f}")
    print(f"  R2    : {r2:>14.4f}")
    print(f"  WMAPE : {wmape_val*100:>13.2f}%")
    print(f"  MAPE  : {mape:>13.2f}%")
    return {'Model': name, 'MAE': mae, 'RMSE': rmse,
            'R2': r2, 'WMAPE(%)': wmape_val*100, 'MAPE(%)': mape}

results = []
results.append(evaluate("RF Baseline",
    RandomForestRegressor(random_state=42, n_jobs=-1),
    X_train_s2, y_train_s2, X_test_s2, y_test_s2))
results.append(evaluate("RF + GA",
    RandomForestRegressor(**best_params_rf, random_state=42, n_jobs=-1),
    X_train_s2, y_train_s2, X_test_s2, y_test_s2))
results.append(evaluate("XGBoost Baseline",
    XGBRegressor(random_state=42, verbosity=0),
    X_train_s2, y_train_s2, X_test_s2, y_test_s2))
results.append(evaluate("XGBoost + GA",
    XGBRegressor(**best_params_xgb, random_state=42, verbosity=0),
    X_train_s2, y_train_s2, X_test_s2, y_test_s2))

results_df = pd.DataFrame(results)
print("\n== Tabel Ringkasan ==")
print(results_df.to_string(index=False))
print(f"\nModel terbaik (WMAPE): "
      f"{results_df.loc[results_df['WMAPE(%)'].idxmin(), 'Model']}")
\end{lstlisting}

\begin{lstlisting}[style=pycode,caption={Deteksi overfitting -- perbandingan Train vs Test}]
def evaluate_train_test(name, model, X_tr, y_tr, X_te, y_te):
    model.fit(X_tr, y_tr)
    pred_tr = model.predict(X_tr)
    pred_te = model.predict(X_te)

    def _m(y, pred):
        return (mean_absolute_error(y, pred),
                np.sqrt(mean_squared_error(y, pred)),
                r2_score(y, pred),
                wmape(y, pred) * 100)

    tr, te = _m(y_tr, pred_tr), _m(y_te, pred_te)
    return {
        'Model'      : name,
        'Train MAE'  : tr[0], 'Test MAE'  : te[0],
        'Train RMSE' : tr[1], 'Test RMSE' : te[1],
        'Train R2'   : tr[2], 'Test R2'   : te[2],
        'Train WMAPE': tr[3], 'Test WMAPE': te[3],
    }

models_ov = [
    ('RF Baseline',      RandomForestRegressor(random_state=42, n_jobs=-1)),
    ('RF + GA',          RandomForestRegressor(**best_params_rf, random_state=42, n_jobs=-1)),
    ('XGBoost Baseline', XGBRegressor(random_state=42, verbosity=0)),
    ('XGBoost + GA',     XGBRegressor(**best_params_xgb, random_state=42, verbosity=0)),
]

ov_results = [evaluate_train_test(n, m, X_train_s2, y_train_s2,
                                  X_test_s2, y_test_s2)
              for n, m in models_ov]
ov_df = pd.DataFrame(ov_results).set_index('Model')
print(ov_df.round(4).to_string())
\end{lstlisting}

% \begin{lstlisting}[style=pycode,caption={Helper: membangun fitur untuk satu bulan forecast}]
% import math

% def build_features_for_month(hist_df, static_last, year, month):
%     rev = hist_df['TotalRevenue'].values.astype(float)
%     act = hist_df['is_active'].values.astype(float)
%     n   = len(rev)

%     def lag(arr, k):   return float(arr[-k]) if n >= k else 0.0
%     def rmean(arr, w): return float(np.mean(arr[-w:])) if n >= w else (float(np.mean(arr)) if n > 0 else 0.0)
%     def rstd(arr, w):  return float(np.std(arr[-w:], ddof=0)) if n >= w else 0.0

%     streak = 0
%     for v in reversed(act):
%         if v == 1: streak += 1
%         else: break

%     msa = 0
%     for v in reversed(act):
%         if v == 1: break
%         msa += 1

%     return {
%         'active_lag1': lag(act,1), 'active_lag2': lag(act,2),
%         'active_lag3': lag(act,3),
%         'active_streak': float(streak), 'months_since_active': float(msa),
%         'pct_active_3m': rmean(act,3), 'pct_active_6m': rmean(act,6),
%         'rev_lag1': lag(rev,1), 'rev_lag2': lag(rev,2),
%         'rev_lag3': lag(rev,3),
%         'rolling_mean_3': rmean(rev,3), 'rolling_mean_6': rmean(rev,6),
%         'rolling_std_3' : rstd(rev,3),  'rolling_std_6' : rstd(rev,6),
%         'TotalSentMessages'     : float(static_last.get('TotalSentMessages', 0)),
%         'TotalDeliveredMessages': float(static_last.get('TotalDeliveredMessages', 0)),
%         'delivery_rate'         : float(static_last.get('delivery_rate', 0)),
%         'profit_margin'         : float(static_last.get('profit_margin', 0)),
%         'rateCost'              : float(static_last.get('rateCost', 0)),
%         'unique_OA_count'       : float(static_last.get('unique_OA_count', 0)),
%         'max_OA_revenue'        : float(static_last.get('max_OA_revenue', 0)),
%         'anomaly_count'         : float(static_last.get('anomaly_count', 0)),
%         'anomaly_revenue_pct'   : float(static_last.get('anomaly_revenue_pct', 0)),
%         'month_sin': math.sin(2 * math.pi * month / 12),
%         'month_cos': math.cos(2 * math.pi * month / 12),
%     }
% \end{lstlisting}

% \begin{lstlisting}[style=pycode,caption={Loop forecast iteratif Jan--Mar 2025}]
% FORECAST_MONTHS = [(2025, 1), (2025, 2), (2025, 3)]
% MONTH_NAME = {1:'Jan', 2:'Feb', 3:'Mar', 4:'Apr', 5:'May', 6:'Jun',
%               7:'Jul', 8:'Aug', 9:'Sep', 10:'Oct', 11:'Nov', 12:'Dec'}

% all_series_ids = monthly['Series_ID'].unique()
% forecast_rows  = []

% for (yr, mo) in FORECAST_MONTHS:
%     feat_list   = []
%     series_order = []

%     static_last_map = (panel_rolling.groupby('Series_ID')[STATIC_COLS]
%                        .last().to_dict('index'))

%     for sid in all_series_ids:
%         hist    = panel_rolling[panel_rolling['Series_ID'] == sid].sort_values('period')
%         statics = static_last_map.get(sid, {c: 0.0 for c in STATIC_COLS})
%         feat_list.append(build_features_for_month(hist, statics, yr, mo))
%         series_order.append(sid)

%     feat_df      = pd.DataFrame(feat_list)
%     pred_active  = clf.predict(feat_df[FEATURES_STAGE1].values)
%     pred_prob    = clf.predict_proba(feat_df[FEATURES_STAGE1].values)[:, 1]

%     pred_revenue = np.zeros(len(series_order))
%     active_idx   = np.where(pred_active == 1)[0]
%     if len(active_idx) > 0:
%         X_s2_fut = feat_df.iloc[active_idx][FEATURES_STAGE2].values
%         pred_revenue[active_idx] = np.maximum(0, reg_final.predict(X_s2_fut))

%     for i, sid in enumerate(series_order):
%         forecast_rows.append({
%             'Series_ID'   : sid,
%             'Year'        : yr,
%             'Month'       : mo,
%             'Month_Name'  : MONTH_NAME[mo],
%             'Pred_Active' : int(pred_active[i]),
%             'Prob_Active' : round(float(pred_prob[i]), 4),
%             'Pred_Revenue': round(float(pred_revenue[i]), 2),
%         })

%     # Update rolling panel
%     new_rows = pd.DataFrame([{
%         'Series_ID'   : sid,
%         'Year'        : yr, 'Month': mo,
%         'period'      : yr * 12 + mo,
%         'is_active'   : int(pred_active[i]),
%         'TotalRevenue': float(pred_revenue[i]),
%         **{c: static_last_map.get(sid, {}).get(c, 0.0) for c in STATIC_COLS},
%     } for i, sid in enumerate(series_order)])
%     panel_rolling = pd.concat([panel_rolling, new_rows], ignore_index=True)

%     n_active  = int(pred_active.sum())
%     total_rev = float(pred_revenue.sum())
%     print(f"  {MONTH_NAME[mo]} {yr}  ->  "
%           f"Aktif: {n_active:>4} series  |  "
%           f"Total Revenue: {total_rev:>16,.2f}")

% forecast_df = pd.DataFrame(forecast_rows)
% forecast_df.to_csv('forecast_jan_mar_2025.csv', index=False)
% \end{lstlisting}

% \begin{lstlisting}[style=pycode,caption={Plot 10 sample Series\_ID: historis + forecast}]
% import matplotlib.patches as mpatches
% import matplotlib.lines   as mlines

% MONTH_SH = {1:'Jan',2:'Feb',3:'Mar',4:'Apr',5:'May',6:'Jun',
%             7:'Jul',8:'Aug',9:'Sep',10:'Oct',11:'Nov',12:'Dec'}

% active_in_fc = (forecast_df.groupby('Series_ID')['Pred_Active']
%                 .sum().loc[lambda s: s > 0].index.tolist())
% np.random.seed(42)
% sample_ids = list(np.random.choice(active_in_fc, 10, replace=False))

% fig, axes = plt.subplots(2, 5, figsize=(22, 9))
% axes = axes.flatten()
% N_HIST_SHOW = 12

% for ax, sid in zip(axes, sample_ids):
%     hist = (monthly[monthly['Series_ID'] == sid]
%             .sort_values('period').tail(N_HIST_SHOW))
%     hist_rev    = hist['TotalRevenue'].values.astype(float)
%     hist_labels = [f"{MONTH_SH[int(r.Month)]}-{int(r.Year)%100:02d}"
%                    for r in hist.itertuples()]

%     fc = forecast_df[forecast_df['Series_ID'] == sid].sort_values('Month')
%     fc_rev    = fc['Pred_Revenue'].values.astype(float)
%     fc_active = fc['Pred_Active'].values.astype(int)
%     fc_labels = [f"{MONTH_SH[int(r.Month)]}-25" for r in fc.itertuples()]

%     n_h  = len(hist_labels)
%     x_h  = np.arange(n_h)
%     x_f  = np.arange(n_h, n_h + len(fc_labels))

%     ax.plot(x_h, hist_rev, color='steelblue', linewidth=1.8,
%             marker='o', markersize=3.5, zorder=3)
%     ax.fill_between(x_h, hist_rev, alpha=0.12, color='steelblue')

%     if n_h > 0 and len(x_f) > 0:
%         ax.plot([x_h[-1], x_f[0]], [hist_rev[-1], fc_rev[0]],
%                 color='darkorange', linewidth=1.5,
%                 linestyle='--', alpha=0.6)

%     ax.plot(x_f, fc_rev, color='darkorange', linewidth=2,
%             marker='s', markersize=6, linestyle='--', zorder=4)

%     for xj, rv, act in zip(x_f, fc_rev, fc_active):
%         if act == 0:
%             ax.scatter([xj], [rv], marker='x', color='red',
%                        s=80, linewidths=2, zorder=5)

%     ax.grid(alpha=0.2)

% h1 = mpatches.Patch(color='steelblue',  alpha=0.8, label='Historis (aktual)')
% h2 = mpatches.Patch(color='darkorange', alpha=0.8, label='Forecast Jan-Mar 2025')
% h3 = mlines.Line2D([], [], color='red', marker='x', linestyle='None',
%                    markersize=8, markeredgewidth=2, label='Idle (inactive)')

% fig.legend(handles=[h1, h2, h3], loc='lower center', ncol=3,
%            fontsize=10, bbox_to_anchor=(0.5, -0.01))
% fig.suptitle('Revenue Historis & Forecast Jan-Mar 2025',
%              fontsize=13, fontweight='bold', y=1.01)
% plt.tight_layout(rect=[0, 0.05, 1, 1])
% plt.savefig('forecast_sample_series.png', dpi=150, bbox_inches='tight')
% plt.show()
% \end{lstlisting} % Bisa dikasih `comment` kalau tidak dipakai
    \chapter*{Halaman Persembahan}

    % Persembahan
    \begin{enumerate}
        % Buat siapa aja?
        \item
    \end{enumerate}
\clearpage
\chapter*{Biodata Penulis}

    % Foto pribadi 
    \begin{wrapfigure}{l}{0.3\textwidth}
        
\includegraphics[width = \linewidth]{tugas-akhir/assets/images/its.png} % Ganti fotonya sesuai tempat direktorinya  
    \end{wrapfigure}
    % 
    Nama lengkap penulis Jiryan Farokhi, lahir di Kabupaten Pati pada 13 Juli 2003. Sebelum menempuh bangku perkuliahan, penulis memulai pendidikan formal mulai dari TK RA Hidayatul Islam, MI Hidayatul Islam, MTs Salafiyah, MA Salafiyah Kajen. Pada tahun 2022, penulis meneruskan pendidikan ke jenjang strata-1 di Departemen Matematika Institut Teknologi Sepuluh Nopember (ITS). Pada masa perkuliahan, penulis aktif dalam kegiatan ekstra kampus antara lain, Academics Lead IEEE ITS (2023-2024), Data Scientist Lead Avalon AI (2024-2025), Machine Learning Lead GDGoC ITS (2024-2025). Penulis juga aktif dalam kegiatan internasional yang diselenggarakan oleh Departemen Matematika ITS yaitu menjadi \textit{awardee} dalam program Global Project Based Learning (GPBL) Data Science Program di Jepang bekerja sama dengan Shibaura Institute Technology (SIT) pada tahun 2024. Di bidang keahlian, penulis memiliki keterampilan pada analisis data, pemrograman (Java, Python), desain antarmuka, AI, serta data science. Sejumlah sertifikasi dari Datacamp, Microsoft, Coursera, NVIDIA dan Dicoding turut memperkuat kompetensi penulis dalam teknologi dan inovasi digital. Adapun informasi lebih lanjut dapat menghubungi penulis melalui email jiryanfarokhi@gmail.com.
\end{document}